\documentclass[11pt]{article}

\usepackage[utf8]{inputenc}
\usepackage[T1]{fontenc}
\usepackage[margin=1in]{geometry}
\usepackage{amsmath,amssymb}
\usepackage{graphicx}
\usepackage{longtable,booktabs}
\usepackage{microtype}
% hyperref goes last; it makes \cite clickable through to the bibliography.
\usepackage[colorlinks=true,linkcolor=blue,citecolor=teal,urlcolor=magenta]{hyperref}

% pandoc emits \tightlist but expects the preamble to define it.
\providecommand{\tightlist}{%
  \setlength{\itemsep}{0pt}\setlength{\parskip}{0pt}}

% If the compile times out on Overleaf, comment out \tableofcontents below first.
\title{Retrieval-Augmented Generation for Large Language Models: A Comprehensive Survey of Paradigms, Techniques, Evaluation, and Future Directions}
\author{Generated by AutoSurvey}
\date{}

\begin{document}
\maketitle
\tableofcontents
\newpage

\section{Foundations, Evolution, and Core Methodologies of Retrieval-Augmented Generation}

\subsection{The Evolution of RAG Paradigms: From Naive to Advanced and Modular Architectures}

{\def\LTcaptype{none} % do not increment counter
\begin{longtable}[]{@{}
  >{\raggedright\arraybackslash}p{(\linewidth - 0\tabcolsep) * \real{0.0556}}@{}}
\toprule\noalign{}
\endhead
\bottomrule\noalign{}
\endlastfoot
The evolution of Retrieval-Augmented Generation (RAG) represents a fundamental shift in how large language models (LLMs) interact with external knowledge, moving from static, parametric-only systems to dynamic architectures capable of integrating real-time, non-parametric information. This progression is not a single leap but rather a continuous refinement that can be broadly categorized into three sequential paradigms: Naive RAG, Advanced RAG, and Modular RAG. Each paradigm marks a distinct advancement in addressing the core limitations of LLMs, namely hallucination, outdated knowledge, and the lack of traceable reasoning processes \cite{ref1}. Understanding this historical trajectory is essential for contextualizing the tripartite architecture detailed above; the modules of retrieval, generation, and augmentation are not static components but have co-evolved through these paradigms, with each stage refining how these components interact and where the primary design effort is concentrated. \\
The foundational stage of this evolution is known as Naive RAG. This paradigm represents the earliest attempt to combine retrieval with generation, operating through a straightforward ``retrieve-then-generate'' pipeline. In this model, a user query is first processed by a retrieval module---typically a sparse or dense retriever---that searches for the most relevant document chunks from an external knowledge base. These retrieved chunks are then concatenated and prepended to the original query, forming an enriched prompt that is fed to the LLM for final answer generation. The core strength of Naive RAG lies in its conceptual simplicity: it provides the model with the most relevant context, thereby grounding generation in specific, verifiable sources. However, this paradigm suffers from several pronounced limitations that restrict its effectiveness. The \cite{ref2} outlines that Naive RAG faces significant challenges in retrieval precision, often retrieving documents that are partially or entirely irrelevant, and in recall, frequently missing crucial passages needed for the answer, leading to hallucinated or incomplete responses. Additionally, the pipeline is vulnerable to retrieval noise, where redundant or conflicting passages can be passed to the LLM without the system being able to select or rank them, ultimately degrading output quality. Furthermore, Naive RAG lacks adaptive capabilities; it always retrieves a fixed number of documents regardless of whether the query actually requires external augmentation, leading to computational waste or insufficient context in complex scenarios. The strictly phased approach also offers no opportunity for iterative refinement, making it a rigid architecture that struggles to handle multi-hop reasoning or query complexities. \\
Recognizing these flaws, the research community developed a more sophisticated framework: \textbf{Advanced RAG}. This paradigm is characterized by a set of optimizations applied before and after the retrieval step---often labeled pre-retrieval and post-retrieval processing---to correct the shortcomings of the Naive approach \cite{ref1} and to improve the quality of the context used for generation. In the pre-retrieval phase, the primary goal is to improve the quality of the query and the indexed data, directly refining the augmentation component described in the previous section. Query-side enhancements include techniques like query rewriting, where the user's original question is rephrased to improve clarity; query expansion, which adds related terms to the original query to bridge the lexical gap; and query decomposition, which breaks down complex multi-hop questions into smaller, simpler sub-queries. On the data side, indexing optimizations involve sophisticated chunking strategies that overcome fixed-size splitting, an especially critical challenge in domains like financial documents where information is densely packed within long, continuous passages \cite{ref3}. These advanced chunking algorithms break documents into sections based on semantic coherence, structural boundaries, and metadata annotations, ensuring that each chunk is self-contained and contextually rich. The post-retrieval phase, conversely, focuses on refining the retrieved documents \emph{after} they have been fetched. This includes a crucial step of reranking. Instead of presenting the retriever's top-k results directly to the LLM, Advanced RAG uses a cross-encoder or LLM-based reranker to score and reorder the documents, filtering out noise and eliminating irrelevant passages before they reach the generation module. Advanced RAG can also employ context compression techniques, which distill the retrieved paragraphs into QA-specific contexts that retain only the most critical sentences or phrases. These optimizations mitigate the noise sensitivity noted in the generation module's capacity limitations, making the context more precise and information-dense, leading to improved accuracy and reduced hallucination rates. For instance, experimental evaluations comparing Naive vs Advanced RAG architectures show that the use of Hypothetical Document Embedding (HyDE) and LLM reranking significantly improves retrieval precision \cite{ref4}. While the same report notes that Maximal Marginal Relevance (MMR) and Cohere rerankers showed no significant advantage over the baseline Naive RAG, certain context methods such as Sentence Window Retrieval and Document Summary Index demonstrated variable results, highlighting the variability within Advanced RAG \cite{ref4}. \\
The latest and most adaptive stage of this evolution is the \textbf{Modular RAG} paradigm, which moves beyond the linear sequence of Advanced RAG towards a flexible and composable architecture. Modular RAG treats the RAG system as a collection of interchangeable and reconfigurable modules, allowing it to be adapted to the demands of the task \cite{ref1}. It also incorporates specialized modules such as memory, routing, and fusion, providing a higher degree of control and enabling the integration of adaptive strategies. For example, in a modular RAG system, the retrieval module can support iterative searching---where the LLM may run multiple retrieval rounds based on initial results---and an active retrieval mechanism, where the model self-determines whether to query a knowledge base or rely on its parametric knowledge \cite{ref1}. This is a direct response to issues like hallucination and knowledge staleness. Moreover, Modular RAG supports a routing module to send queries to specialized retrievers and to fuse results from multiple sources, enabling a broad range of domain-specific adaptations. It seamlessly adds ``post-refinement'' modules, such as re-rankers or task-specific aggregators, which can be tuned to adapt to the domain. This approach is exemplified in the Retrieve-Augmented Iterative Self-Feedback (RA-ISF) framework, which iteratively decomposes tasks and refines outputs via self-critique and selective retrieval, showing enhanced performance over existing baselines \cite{ref5}. \\
In essence, the trajectory from Naive \(\rightarrow\) Advanced \(\rightarrow\) Modular marks the progression from a simplistic, static pipeline to a dynamic, interactive thought process. While Advanced RAG focuses on optimizing the individual components---particularly the augmentation and retrieval modules---to improve the quality of retrieved documents, Modular RAG focuses on creating a more flexible and adaptable architecture, where the system becomes integrated into the LLM's reasoning loop. This evolution represents the explored direction of hybrid augmentation, where the LM augments its output by navigating, refining, and synthesizing the retrieved knowledge in a structured yet flexible flow. The latest stride, often revolving around the concept of fine-tuning entire RAG architectures, including the retriever, shows that not only is the modular system capable of planning and retrieval, but the very internal components can also be jointly trained for better performance \cite{ref6}. This end-to-end fine-tuning approach leverages the accuracy of supervised adaptation while exploiting the overarching flexibility of modular orchestrations. The outcome, as will be examined in the following subsection, is a solid progression towards modern RAG systems that serve as a robust foundation for future advances across multimodal, knowledge-intensive, and dynamic environments, built upon the finely tuned interaction of retriever, generator, and augmentation strategies detailed above. \\
\#\# Core Architectural Components: Retrieval, Generation, and Augmentation \\
The tripartite architecture of Retrieval-Augmented Generation (RAG) represents a fundamental departure from purely parametric language modeling, establishing a framework where the strengths of large language models (LLMs) are synergistically combined with dynamic, non-parametric knowledge sources. This architectural foundation, which has evolved significantly since its inception through the paradigms of Naive, Advanced, and Modular RAG outlined previously, comprises three critical, interlocking modules: the retrieval module, the generation module, and the augmentation strategies that bridge them. Understanding the distinct roles and complex interactions of these components is essential for grasping the capabilities and limitations of modern RAG systems \cite{ref1,ref2}. Furthermore, the design choices made within each module---particularly the selection of retrieval strategies---directly set the stage for the detailed taxonomy of techniques examined in the subsequent subsection. \\
The retrieval module serves as the gateway to external knowledge, responsible for sourcing relevant information from a corpus, knowledge base, or datastore. Its primary objective is to transform a user query and a large, often unstructured, collection of documents into a concise set of snippets that are most likely to contain the information required to pursue the query. The effectiveness of this module is paramount, as it establishes an upper bound on the quality of the information that the generation module can utilize. The design of the retrieval module encompasses a diverse array of techniques, ranging from sparse methods like BM25, which rely on exact lexical matching, to dense retrieval methods that encode both queries and documents into a shared semantic vector space---a spectrum that will be exhaustively categorized in the following subsection. The choice of retrieval strategy is highly consequential, and it is the subject of continual innovation. A critical insight from empirical research is that retrievers are not monolithic and their performance is deeply interwoven with the entire RAG architecture. For example, the ``RAGGED'' framework demonstrated that different large language models (LMs) have vastly different sensitivity to the retrieval component; encoder-decoder models, for instance, are shown to rely heavily on retrieved context, making them more sensitive to retrieval quality, while decoder-only models tend to rely on knowledge already memorized during pre-training, making them less responsive to the exact content of the retrieved set \cite{ref7}. The capabilities of the retriever are also heavily influenced by the nature of the queries it faces. Benchmarks like MultiHop-RAG have highlighted the standard retrievers' inadequacy in handling complex queries that require reasoning over multiple pieces of supporting evidence, motivating the iterative, multi-hop, and graph-based retrieval methods discussed later \cite{ref8}. \\
The generation module represents the second pillar of the RAG architecture. It consists of a large language model that integrates the user's query and the context provided by the retrieval module to generate a natural language response. In the foundational understanding of RAG, its role was considered to be a passive ``knowledge receptor'' and answer generator. However, more recent thinking postulates that this is a passive role; LLMs have the capability to actively participate in the process by querying or reasoning about external knowledge \cite{ref1}. In this context, the generation module not only synthesizes a fluent and coherent answer but performs resource radical interaction, seamlessly fusing the parametric knowledge acquired during training with the non-parametric knowledge derived from the retrieved context. This process of knowledge fusion is rarely straightforward and often involves navigating potential conflicts between internal parametric knowledge and conflicting external evidence, as well as deciding which information source is reliable \cite{ref9}. The generation module's capacity, especially in terms of its context length, also defines the maximum amount of retrieved information that can be considered for a response \cite{ref10}. For instance, certain decoder-only models can only effectively utilize less than five documents before performance degrades, despite having larger context windows, showing that some models are less adept at managing irrelevant information \cite{ref7}. This capacity constraint directly influences which retrieval strategies---such as the compression and reranking methods described in the next section---are needed to ensure the generated output remains both accurate and grounded. \\
The ``augmentation'' strategies constitute the most critical and integrative component of the RAG architecture, serving as a spectrum of techniques that fuse the parametric and non-parametric knowledge. These strategies define the exact pipeline for how the retrieved information is processed and integrated with the user's query to produce a compelling prompt for the generation model. This process is far from being a static task of simple concatenation; it is a domain with a large range of design decisions that can heavily influence the final output \cite{ref1}. The augmentation stage can be broadly divided into pre-retrieval and post-retrieval processes. Pre-retrieval processes involve optimizing the index and the original user prompt---techniques such as query expansion, rewriting, or decomposition---while post-retrieval focus on reranking and compressing the retrieved content to reduce noise and select the most relevant information \cite{ref1}. The augmentation component is responsible for addressing the rich complexity of the context and can be dynamically modified to address different tasks. For example, post-retrieval augmentation can address the challenge of noisy, irrelevant documents, and studies have shown that, surprisingly, adding slightly relevant or irrelevant documents can occasionally provide performance boosts, underscoring the need for sophisticated integration methods that guide the model on how to use the retrieved context \cite{ref10}. The augmentation stage is also the primary integration point for query processing techniques, such as query expansion, decomposition, or multi-query generation, that bridge the gap between user intent and retrieval inspect \cite{ref1}. \\
In summary, the tripartite architecture of RAG is a dynamic and heterogeneous system in which each component is not only important on its own but also part of a larger, finely orchestrated system that requires careful co-design. The architectural integrity depends on continuous interaction and feedback among the modules \cite{ref1}. Studies such as ``RAGGED'' and ``The Power of Noise'' demonstrate that the retrieval module cannot be considered in a vacuum; its output quality is only as useful as the capacity of the generation module to ingest and utilize it \cite{ref7}. The augmentation component orchestrates the entire pipeline, from optimizing the query and context to teaching the model how to best use the retrieved knowledge, ensuring both the system's efficiency and robustness \cite{ref1}. Furthermore, the iterative refinement seen in Modular RAG---where retrieval can be repeated and adapted based on generation feedback---underscores how the modules' interaction is bidirectional and task-dependent. The next subsection builds directly on this architectural foundation by systematically categorizing retrieval techniques---sparse, dense, multi-vector, iterative, graph-based, and hybrid---and demonstrating how these choices interact with the generation and augmentation components described above to balance accuracy, latency, and index size in practice. \\
\#\# Retrieval Strategies and Techniques \\
Retrieval strategies form the backbone of Retrieval-Augmented Generation (RAG) systems, determining the quality, relevance, and efficiency with which external knowledge is sourced to augment large language models. Over the years, a rich taxonomy of retrieval techniques has emerged, each with distinct strengths, trade-offs, and applicability across various task settings. This subsection systematically organizes the retrieval landscape into six major families: sparse retrieval, single-representation dense retrieval, multi-representation (multi-vector) retrieval, iterative and active retrieval, multi-hop and graph-based retrieval, and hybrid retrieval. \\
\textbf{Sparse Retrieval: The Traditional Workhorse.} Sparse lexical methods, epitomized by BM25, remain a foundational baseline in RAG pipelines due to their interpretability, computational efficiency, and strong performance on exact-match queries. BM25 operates on bag-of-words representations, making it highly effective when user queries share significant lexical overlap with relevant passages. Its rigidity, however, is its limitation: lexical matching systems frequently miss semantically relevant passages that lack exact term overlap, particularly for abstract or definitional queries \cite{ref11,ref12}. Despite being long-standing, BM25 continues to anchor many modern systems, often used as a first-stage ranker or as a check on dense retrieval outputs. It is especially useful in hybrid pipelines, where its diverse candidate sourcing compensates for dense models' blind spots. \\
\textbf{Single-Representation Dense Retrieval (DR).} Dense retrieval overcomes lexical sparsity by mapping both queries and passages into a shared continuous embedding space, enabling semantic rather than surface-form matching. In the single-representation family, each passage is collapsed into one vector---often the CLS token from a BERT encoder---and queried using nearest-neighbor search. Semantic search beyond lexical overlap allows dense methods to address the vocabulary mismatch problem that query processing techniques aim to bridge. The seminal Dense Passage Retriever (DPR) \cite{ref12} demonstrated that this approach could outperform BM25 in top-20 retrieval accuracy when trained end-to-end on a small number of question-passage pairs, leading to state-of-the-art results in open-domain question answering \cite{ref12}. Follow-up work on ANCE and other models further improved dense training stability and efficiency through hard-negative mining and asynchronous embedding updates. \\
Nevertheless, single-vector encoding involves a loss of local detail. To mitigate this, a dual-encoder architecture that allows independent encoding is often adopted to maintain throughput, though this structure forfeits query-document interaction during encoding. To regain some interaction without sacrificing efficiency, GNN-encoder introduces a graph neural network that fuses query information into passage representations and vice versa at the representation level, creating a denser, interaction-enriched embedding at low computational cost \cite{ref13}. Similarly, Topic-DPR tackles the problem of collapsed semantic space: rather than optimizing a single continuous prompt per query, it learns multiple topic-based prompts over a probabilistic simplex, using contrastive learning to spread the representation space, and finds this yields more uniform, and semantically separated, embeddings \cite{ref14}. \\
Recent mechanistic studies have further illuminated dense retriever internal behavior. Probing and model editing experiments on DPR-trained models show that training decentralizes knowledge storage across the network, creating multiple access pathways---but that the pre-trained model's internal knowledge places an upper bound on what the retrieval model can retrieve. This suggests directions for injecting facts as decentralized representations, or building connections to a knowledge base \cite{ref15}. Moreover, a replication study of DPR confirmed the effectiveness of dense passage retrieval against BM25; however, highlighted the original paper under-reports the strength of BM25, and that combining dense and sparse signals gives substantial hybrid boosts \cite{ref16}. \\
\textbf{Multi-Vector (Multi-Representation) Retrieval: Late Interaction and Beyond.} Multi-vector methods combine the semantic richness of dense retrieval with the precision of lexical matching. Unlike single-vector representations, these methods encode each token of a passage and query into separate vectors. ColBERT leads this family within the RAG literature, using a late interaction mechanism that only matches query tokens against passage tokens at the end, preserving selectivity while outperforming single-vector methods on MAP/MRR and showing greater robustness in queries that are harder for BM25 \cite{ref11}. CITADEL moves beyond with conditional token interaction via dynamic lexical routing: token embeddings are routed into predicted lexical keys, such that a query's token interacts only with document tokens routed to the same key. This design achieves accuracy comparable to ColBERT-v2, but is about 40x faster on both in-domain and out-of-domain evaluations, making it a practical multi-vector for industrial-scale settings \cite{ref17}. \\
\textbf{Iterative, Active, and Recursive Retrieval.} Static retrieval is facing limitations in RAG: for complex questions a single retrieval step may be insufficient. The emergence of complex, multi-hop queries---which multi-hop benchmarks have shown to defeat standard retrievers---motivates a move toward iterative strategies. A promising direction is active retrieval---where the model decides whether to fetch additional evidence---and recursive retrieval, where each cycle continues from the previous outcome to form a chain of evidence. For instance, M3 demonstrates the feasibility of advanced recursive multi-hop dense sentence retrieval built on a multi-task mixed-objective learning framework, where contrastive learning and other objectives press the retriever to find evidence for the next step \cite{ref18}. This shows the retrieval across dependent evidence can outperform single-shot retrieval in bridging the gap between simple and multi-step reasoning. \\
\textbf{Multi-Hop and Graph-based Retrieval.} Multi-hop reasoning requires retrieving several documents whose facts compose together to produce an answer. The M3 system above is built for multi-hop, using multi-hop refinement to achieve state-of-the-art performance on the FEVER benchmark \cite{ref18}. Graph-based retrieval extends this notion coherently: edges between closely related documents are pre-computed, and retrieval walks the graph to gather connected evidence. This is appropriate for questions that necessitate compositional answers, or when the evidential network is non-latent. However, these graphs require frequent updates to remain useful in dynamically changing corpora. \\
\textbf{Hybrid Retrieval: Combining Signals for Recall and Precision.} Another popular direction leverages the combination of sparse lexical matching and dense semantic matching (a natural partnership inside RAG pipelines). Hybrid setups reduce false negatives in both directions---with a term-level retriever also offering unique drawbacks and with a dense semantic channel adding reach. Notably, a replication of DPR showed that hybrid retrieval yields better results than either alone \cite{ref16}. In practice, when combining BM25 and dense dual-encoders as a two-stage pipeline, results are strong for document and passage ranking---but the hybrid stage consistency requires careful calibration \cite{ref19}. At a coarser granularity, dense phrase retrieval has demonstrated that phrase-level retrieval can recreate strong passage retrieval accuracy because---with its supervision required---with filtering and vector quantization---the index size can be reduced by 4--10×, making phrase search versatile for multi-granular retrieval \cite{ref20}. \\
\textbf{Non-Parametric and Zero-Shot Retrieval Adaptation.} Beyond direct comparison of existing signals, retrieval can be adapted to new tasks without labeled data through synthetic training data generation---an idea complementary to llm-driven query expansion discussed in the pre-retrieval step. Embedding-based zero-shot retrieval through query generation has shown significant improvements over BM25 on recall, sometimes outperforming even models trained on real data \cite{ref21}. \\
\textbf{Retrieval Efficiency, Compression, and Context Thematically.} All retrieval methods must reside inside a larger module that handles encoding, routing, and execution in real-time. A deeper optimization relies on efficient encoding: CITADEL offers a 40x reduction in compute for multi-vector retrieval via dynamic routing \cite{ref17}. At the same time, single-vector methods remain very popular due to lightweight encoding and the simplicity of their training---and when coupled with quantization techniques allow for reducing memory and time expenditure with relatively little recall loss, making them suitable for resource constrained and industrial search systems \cite{ref20}. \\
Looking forward, no single retrieval strategy fits all tasks. Sparse, dense, single-vector, and multi-vector, growing hybrid, iterative, recursive, graph-based, or multi-hop logic, to answer multi-hop questions---each presents performance--deployment trade-offs. The optimal selection---methods for br/multi-exchange---rather than the will-be immediately retrievable for all queries---has no substitute in terms of domain and downstream generation models wire. Hybrid-retrieval-corpus-constraint are more often necessary to solve decision when all these considerations---accuracy, latency, index size, user query distribution irrelevance---are excruciatingly intermingled with the generation module. Thus, the choice of the optimal retrieval strategy cannot be made in isolation: it must be made with reference to the constraints and strengths of the target generation stage, the subsequent augmentation techniques (namely, reranking andcompression) that allow filtering and refining of the initial retrieval set, and the end task demands. In the next section, interaction of these methods with generation strategies---in particular, how the generation module handles the context these methods return---is essential for full utilization of the augmentation in the overall RAG pipeline. \\
\#\# Query Processing and Augmentation Techniques \\
Query processing and augmentation techniques form a critical front-end component of Retrieval-Augmented Generation (RAG) systems, addressing the fundamental semantic gap that exists between user-formulated natural language queries and the vocabulary used within target document collections. This mismatch, often termed the vocabulary mismatch problem, arises when users lack familiarity with the terminological landscape of relevant documents, which directly impedes retrieval effectiveness \cite{ref22}. Within the broader RAG pipeline, these techniques operate upstream of the retrieval strategies discussed in the previous subsection: however sophisticated the retrieval backbone may be, its effectiveness is bounded by the quality of the query it receives. Query-side augmentation techniques attempt to bridge this gap by reformulating or enriching the original query before it enters the retrieval phase, thereby enhancing the likelihood of surfacing relevant knowledge that will subsequently augment generation. Within the RAG architecture, the quality of the retrieved context is a dominant determinant of final generation quality; thus, refining queries before retrieval constitutes a vital strategy for reducing downstream hallucination and enhancing answer relevance---a concern that becomes even more critical when considering how the generation module later integrates the retrieved evidence. \\
Query expansion (QE) is among the most longstanding and extensively researched techniques in this domain, aiming to reformulate the initial query by appending semantically related or synonymous terms \cite{ref23}. A primary challenge for QE methods is ensuring that expansion terms balance two competitive properties: diversity and relevance. While sampling a greater number of expansion contexts from language models can boost diversity, it often compromises relevance, due to the well-documented tendency of these models to hallucinate incorrect or extraneous content \cite{ref24}. To address this, approaches such as the combination of filtering strategies along with fusion mechanisms based on the generation probabilities of the expansion context have been proposed, successfully achieving high retrieval recall while substantially reducing the index size \cite{ref24}. Similarly, researchers have recognized that query expansion methods are not uniformly beneficial across different query categories; some techniques improve performance for individual term queries but fail on phrase queries, and vice versa. This observation has motivated the development of taxonomy-driven query classification and targeted expansion strategies \cite{ref25}. \\
The sources of expansion terms are heavily varied. Beyond unigram and phrase-term extraction, external knowledge resources such as Wikipedia, WordNet, and other lexical databases have been found to supply rich, semantically coherent expansion terms \cite{ref26}. In this work, a hybrid approach---leveraging Wikipedia, TermNet, and WordNet---followed by a correlation-based scoring scheme for essential expansion terms yields substantial improvements over unexpanded within heterogeneous test collections \cite{ref26}. An additional line of research focuses on the linguistic attributes of queries to guide expansion. For instance, conventions akin to part-of-speech labels, grammatical roles, and syntactical dependencies have been integrated into expansion frameworks \cite{ref27}. A framework has also been proposed that encodes query intent by recognizing core (central) concepts, descriptive, relational, and structural concepts, leading to more precise concept weighting and robust retrieval improvements \cite{ref27}. A modular and comprehensive decomposition of expansion systems, covering corpus-based statistical statistics, syntactic dependencies, and ontology sources, has been systematically reviewed in the literature \cite{ref28}. \\
The emergence of large language models (LLMs) has triggered a new wave of expansion techniques. In this recent wave of generative query expansion, capitalizing on the model's implicit knowledge extraction capabilities. For instance, prompting LLMs, especially with chain-of-thought (CoT) prompts that decompose queries, can yield many expansion terms and outperform pseudo-relevance-feedback-based expansions on benchmarks like MS-MARCO and BEIR \cite{ref29}. Yet, such generative methods generally lack corpus-specific knowledge, a gap that has required calibration with corpus-specific cues. To address this, MILL (Mutual Verification with Large Language Models for Zero-Shot Query Expansion) leverages zero-shot reasoning from an LLM to generate diverse sub-queries and their convergent documents, then mutually verifies these generated outputs against retrieved documents to select the optimal results, with success demonstrated in benchmark datasets \cite{ref30}. Alternatively, the ensemble prompting approach of generating diversified query reformulations from paraphrased prompts, like GenQREnsemble, and incorporating pseudo-relevant feedback, yields improved retrieval and robustness, affirming the necessary hybridity of generation and verification to combat hallucinations \cite{ref31}. \\
Complementing expansion, reduction techniques specifically target verbose or noisy queries. Long queries, infested with extraneous term, dilute the semantic focus and drift from a precise user' intent. Contextualized Query Reduction (ConQueR) acknowledges this challenge of long verbose queries and implements term-level core extraction and sequence-level sub-query selection, through an ensemble trained and tested on search logs including; the reduction using pre-trained language model achieves better exact-match accuracy over the best physical baseline \cite{ref32}. In the domain of database query rewriting, GenRewrite is a large language-model-based rewriting system that represents rewriting as a natural-language rewriting rule generation problem, achieving higher coverage of old SQL rewrites, demonstrating the versatility of LLM-generalization beyond plain-text retrieval \cite{ref33}. \\
Beyond textual lexicon, knowledge has been integrated into query expansion by exploiting graph-based knowledge bases. Massive Query Expansion by Exploiting Graph Knowledge Bases conducts lexical and topological cues. By adding synonyms and semantically related terms and performing network analysis of concept relationships, research explores and considers that both local and topological features (graph paths and community structure) lead to a higher degree of multimodal diversity and to augmented retrieval outcomes \cite{ref34}. These techniques underline how leveraging structured external knowledge within query processing provides customization that can flexibly integrate diverse, computational, and broad context sources. \\
Another highly creative and interdisciplinary direction emerges at the intersection of neuroscience and language. Brain-Aug enhances enrichment of queries by generating semantic continuations from decoded brain signals, bypassing any dependence on content-language-based approximations, such using functional magnetic resonance imaging (fMRI) representations of concept before query words, to decode semantic conditions that augment effective ranking---especially remarkable for ambiguous queries \cite{ref35}. Though early stage, this shows the ambition and broad scope of query augmentation strategies. \\
Traditional query expansion based on retrieval feedback remains practical and attractive. However, standard pseudo-relevance-feedback-based expansion methods often suffer from being misled by noisy pseudo-retrieval. To mitigate sources, a proposal on division of the documents set via equi-frequency partitioning and the incorporation of document partition statistics (e.g., tf-idf) for constructing a term candidate set, combined with a fusion function, provides a new flavor of robust, low-cost, rule-based coordinated regarding expansion. Certain expansion terms highways exist which do correspond to established set-wise interdocumentation forms, thus leading to candidates considered \cite{ref23}. A learning-based strategy to select better candidate expansion terms relies on deep neural network classifiers in the semantic space. By training of scoring of a binary-like classifier, the filtering of the expansion terms over the embedding representation significantly improves the term selection when compared against simple baselines \cite{ref36}. \\
In summary, query processing and augmentation techniques in RAG systems have evolved from simple lexical additions, to knowledge-rich and graph-informed, to linguistically motivated, and to generative and mutually-verified LLM prompting strategies. These techniques aim to minimize the distance between user intended meaning and actual retrieval capability. The optimal approach strikes a balance among zero-shot generalization, source availability, test rigor, diversity and relevance of expansions, and minimization of the risk of hallucination in the pre-retrieval stages. Additionally, these improved query representations are evaluated and fed into various retrieval strategies---from sparse and dense to hybrid---, and ultimately fused with generated content through the integration mechanisms described in the next subsection, thereby closing the loop between the query-adaptive, retrieval-centric, and generation-oriented layers of the RAG pipeline. \\
\#\# Integration Mechanisms for Augmented Knowledge Fusion \\
The integration of retrieved knowledge into the generation process is a pivotal architectural decision that determines the quality, coherence, and factual grounding of RAG outputs. This subsection examines the diverse mechanisms through which external documents are fused with parametric knowledge within large language models, spanning from simple input composition to sophisticated attention-based strategies. The choice of integration mechanism profoundly impacts not only downstream task performance but also the model's latency, robustness, and explainability. \\
At the most fundamental level, knowledge integration often begins with \textbf{input composition}, where retrieved passages are concatenated with the original query to form an augmented prompt. This approach, popularized by early Retrieval-Augmented Generation (RAG) systems, contextualizes the generator by placing the retrieved evidence directly into its input sequence. For example, the Fusion-in-Decoder architecture employed in multi-task training settings conditions generation on a concatenation of the query and multiple retrieved passages, using a unified encoder-decoder framework to aggregate evidence \cite{ref37}. While simple and effective, the quality of this composition can be resource-intensive, especially when integrating retrieval knowledge in non-knowledge-intensive (NKI) tasks, where existing works that focus on concatenating retrievals to inputs as context require the language model to handle long texts and consume significant computational resources \cite{ref38}. To address this inefficiency, a more computation-efficient approach, ReFusion, directly fuses retrieval representations into the hidden states of a language transformer rather than at the raw token level. Instead of extending the input context, ReFusion fuses representations of similar sentences obtained via an online retrieval module---with a reranker-based scheme and an ordered-mask-based scheme---directly into the middle layers of the network \cite{ref38}. This highlights a shift away from simple concatenation towards a more intricate information routing and fusion process. Furthermore, the role of a precisely defined reranking strategy on top of the initial retrieval is an essential layer of composition; Re2G demonstrates that incorporating a dedicated reranker after neural initial retrieval and BM25 (whose scores are incomparable) allows for a principled ensemble. By fusing the reranked and merged results into a BART-based sequence-to-sequence generator, Re2G shows that integrating a re-ranking component in the augmentation path yields large gains in diverse tasks like question answering and fact-checking, evidencing the ``retrieve-rerank-generate'' paradigm as an effective composition strategy \cite{ref39}. \\
Moving beyond straightforward sequence concatenation, more fine-grained integration mechanisms like \textbf{attention-based sorting and fusion} have been proposed to guide the model's focus during generation. The challenge is that a standard transformer treats retrieved passages in a uniform, linear sequence, neglecting the singular hierarchy of information. To address this, \textbf{fine-grained approaches using self-attention} can be deployed. For instance, using ``fine- and coarse-granularity hybrid self-attention'' reduces computational cost by scoring the informativeness of tokens at each layer, where informative tokens are treated as fine-granularity computing units, while uninformative ones are clustered into coarse units \cite{ref40}. This type of integration mechanism collapses the passing of ``irrelevant'' context to the generation module to save FLOPs while still maintaining the flow of key evidence. Similarly, hierarchical reasoning in a fusion context can borrow principles from the attention mechanisms used in other knowledge retrieval tasks. For instance, \textbf{granularity-aware attention} in paraphrase generation extends multi-head self-attention by automatically inferring the hierarchical structure of a sentence and encoding it via ``granularity resonance'' and ``granularity scope'' masks \cite{ref41}. When adapted to RAG, such granularity-aware relative positions can help sort or re-weight tokens from retrieved passages based on their local relevance, ensuring that low-level details and high-level semantic relations are weighted correctly during answer generation. Furthermore, \textbf{multi-granularity perception networks} (MGPN), used in moment retrieval, allow the system to perceive the input multi-modally through a coarse-grained feature encoder and co-attention for preliminary perception, followed by a fine-grained encoder for conditioned interaction \cite{ref42}. If this idea is transferred to RAG, a similar two-level fusion can be designed that first filters retrieved evidence based on broad relevance and then selects the context conditional on fine-grained details. In a modern generative large language model, these types of mechanisms map onto \emph{bifurcated attention} systems that separate the attention for the prefill KV cache (context from retrieved documents) from the decoding process, enabling high batch sizes and longer contexts with minimal memory IO \cite{ref43}. \\
Retrieval is rarely a one-shot process; sophisticated sequential integration is necessary to ground a generative answer comprehensively. Indeed, an ideal integration module must engage in metacognition---to monitor, evaluate, and plan its response strategy based on the adequacy of retrieved information \cite{ref44}. The metacognitive regulation pipeline identifies the inadequacies of the initial cognitive response and leverages further retrieval or generation to correct them; in practice, this requires the RAG architecture not only to fuse evidence in a single forward pass but to support iterative refinement loops or multiple retrieval-generation rounds. Moreover, the idea of routing task-relevant information into a specialized, shared space, with a bottleneck on less-relevant knowledge, directly relates to retrieval integration: such an architecture suppresses irrelevant context noise while retaining task-relevant knowledge, a fundamental challenge in composition \cite{ref45}. Although designed for continual learning, the routing principle is transferable to RAG, ensuring that the knowledge flow maintains relevance with the surrounding context during generation. \\
\textbf{Context and representation compression} emerges as a crucial technique in the integration process. More accurate fusion mechanisms seek to reduce the extra sequence length introduced by retrieved documents, and \textbf{fine- and coarse granularity hybrid self-attention} does so by scoring informativeness per token, reducing computation while maintaining a compressed sequence of key units \cite{ref40}. When adapted to RAG, this technique acts as an attention-sorting module that compresses the retrieved context by filtering out non-essential tokens, preserving the passages' evidence while lowering the computational overhead of generation. \\
Beyond concatenating top-k documents, \textbf{strategies for query enhancement and dataset-level quality of retrieval} must be tightly coupled to the integration module. \textbf{Multi-Task Retrieval-Augmented Text Generation with Relevance Sampling} demonstrates that removing noise from the training set by filtering examples not answerable by the knowledge base reduces the burden on the integration mechanism, thereby improving the generator's ability to extract the correct evidence \cite{ref37}. This illustrates that integration is not solely a decision made during generation, but also depends on the quality of the retrieval evidence that lies upstream in the pipeline. \\
In summary, effective fusion of retrieved knowledge with parametric memory defines the core of a capable RAG system. It is a multi-faceted integration suite that spans from token-level fusion techniques (ReFusion and granularity-aware attentions) on the one side, to hierarchical attention computation for efficiency \cite{ref40}, the combination with a separate re-ranking component in its sequence-to-sequence pathway \cite{ref39}, and the use of semantic consolidation and compression in distributed memory components. Composition methods like simple concatenation remain a viable foundational method, yet novel reranking, compression, and attention-based routing improve its efficacy. This highlights the dynamic orchestration between parametric and non-parametric knowledge, which is a crucial dimension for the design of future RAG systems. Future implementations of iterative reasoning must embed these integration blocks with care, considering the entire taxonomy of retrieval representation and the precision of selection during each generation step---an open direction that merges the retrieval-centric insights from previous sections with the adaptive strategies for retrieval necessity and self-knowledge utilization that follow. \\
\#\# Adaptive Retrieval Mechanisms and Self-Knowledge Utilization \\
The architectural choices governing knowledge integration, as discussed in the previous subsection, often presuppose that retrieval is a beneficial and necessary first step. However, a growing body of research recognizes that this ``always-retrieve'' strategy is suboptimal. The necessity of external knowledge varies significantly across queries; some are answerable entirely from the model's parametric memory, while others demand fresh, specific, or up-to-date information. The indiscriminate retrieval of external documents can not only waste computational resources and increase latency but can also introduce noise, distract the generator, and paradoxically degrade performance on queries that the model could have handled correctly on its own. This realization has motivated the development of adaptive retrieval mechanisms---systems capable of making dynamic, context-aware decisions about \emph{whether}, \emph{when}, and \emph{how} to retrieve information. This is intrinsically linked to the utilization of a model's own knowledge, where the decision to retrieve is often guided by an assessment of the sufficiency and confidence of the model's parametric memory. \\
The most fundamental form of adaptability---and the natural complement to the integration techniques described earlier---lies in determining the necessity of retrieval itself. Early knowledge-based dialogue systems often overlooked this critical decision, assuming that all conversations require external knowledge. This leads to an over-dependence on external sources, even when they are not required, which can introduce irrelevant information \cite{ref46}. To directly address this, some models unify the tasks of retrieval decision, query generation, and response generation into a single framework, enabling the system to first deliberate on whether external knowledge is needed before generating a retrieval query. Such unified architectures reduce system complexity while allowing for the conditional replacement of retrieval operations \cite{ref46}. \\
A more sophisticated trigger for adaptive retrieval is the model's internal state and self-knowledge. This approach, often termed ``self-knowledge guided retrieval,'' relies on the LLM's ability to assess its own confidence in generating a response. The idea is to let the generative model answer the query from its parametric memory first and then assess the confidence level of this initial prediction. If the model is sufficiently confident in its prediction, retrieval is skipped entirely. Conversely, if it exhibits uncertainty---often measured through a semantic similarity threshold or conformal prediction techniques---the system initiates an active retrieval step to ground and correct the response \cite{ref47}. This can be seen as a robustness-based trigger for calibration that directly enhances the trustworthiness of RAG systems by determining when parametric knowledge is self-sufficient or needs external supplementation \cite{ref47}. This module acts as a dynamic pointer, connecting the adaptive retrieval process to the quality of the parametric knowledge stored within the model's memory. \\
Furthermore, adaptive retrieval can go beyond a simple binary decision and involve evaluating the complexity of the user's query. The system can utilize the LLM's capacity to reason about the intermediate steps of a complex query. For instance, a model might first generate several high-level search queries in previous factuality evaluation tasks, breaking down the task into sub-tasks that can be answered more effectively \cite{ref48}. Based on these intermediate evaluations, the model can determine the optimal retrieval depth and breadth. For simple, well-known facts, the model can directly locate the relevant passage and answer, while for complex, multi-step questions, it might generate new sub-queries and iterate to gather information from distinct sources \cite{ref48}. This integration of planning and retrieval execution is a key feature of advanced adaptive systems. It significantly enhances the quality of the evidence supplied to the integration mechanisms, ensuring that the retrieval strategy is not a uniform process but is dynamically tailored to the specific semantic requirements and complexity levels of each query. \\
Context-aware reflexive retrieval is another, even more granular, form of adaptability. Instead of making a single global decision about a query, these mechanisms monitor and evaluate the model's output during the entire generation process. For instance, the model can pause the decoding process at strategic points to iteratively consult the retrieval module and access critical information for the next segments of generation \cite{ref49,ref50}. This alternates the possibility of interleaving retrieval and generation single steps, enabling the model to perform multiple retrieval rounds per query at strategic positions, perfectly complementing the attention-based integration modules from the previous sections. This is particularly effective for long-form generated content, where the integration of newly retrieved evidence at later reasoning steps can correct error cascades from earlier stages and enhances the overall coherence of the final response. \\
The foundation of these adaptive strategies is a shift from a \emph{contextual extraction} to a \emph{dynamically controlled generation} paradigm. These adaptive strategies respond to the need for dynamic knowledge access: by integrating decision steps that evaluate the model's parametric confidence, the query's complexity, and the contextual requirements of the generation process itself, these systems ensure retrieval is a valuable and planned intervention rather than an automatic action. This adaptive control is crucial for improving the accuracy and robustness of a model that must also handle the dynamic nature of its ever-growing knowledge sources, and it sets the stage for the more explicit, persistent memory frameworks discussed in the next subsection---where active experience not only triggers retrieval but reshapes the environment itself for seamless and efficient acquisition of knowledge. \\
\#\# Memory and Knowledge Augmentation Techniques \\
The integration of memory and knowledge augmentation layers has emerged as a transformative direction within Retrieval-Augmented Generation (RAG), moving beyond static one-shot retrieval paradigms toward systems capable of dynamic, lifelong learning and adaptation. This subsection examines how additional memory mechanisms---ranging from key-value structures to sequential frameworks and feedback-driven learning---extend the core RAG tripartite architecture by providing persistent, evolving knowledge repositories that work in concert with the primary datastore. These techniques address a fundamental limitation of conventional RAG, i.e., the static nature of the knowledge base, thereby advancing the field from the foundational Naive and Advanced RAG paradigms toward a more fluid and adaptive form of Modular RAG \cite{ref1,ref2}. \\
A central theme in memory augmentation is the implementation of key-value versus vector memory structures that function as an intermediate layer between user queries and the generation model. Within the broader exploration of RAG's modular components, the retrieval stage often relies on vector embeddings to fetch relevant document chunks \cite{ref1}. Key-based memory augmentation extends this concept by allowing the system to store pairs of certain queries (keys) or concepts tied to specific response formulations (values), enabling more nuanced association-level retrieval. This is particularly impactful in domain-specific setups where the embedding model might be suboptimal for extracting latent semantics, and a key-based layer can funnel the model toward more specific, previously contextualized snippets before generating a response. Furthermore, decoupling query-answer pairs ensures that the pre-retrieval and retrieval stages are better evaluated, addressing the issue where generic retrieval via dense passage or BM25 retrieval often underperforms on update tasks, as shown in benchmarks covering the CRUD operation categories \cite{ref51}. These methods create an organizational semantic scaffold that helps LLMs navigate complex knowledge structures, bridging the gap between a basic retriever and the iterative refinement in the generation module. \\
Deeper into the stack, sequential memory frameworks have been crucial in establishing persistent generative context. Standard practice in RAG pipelines often relies on inserting top-k retrieved documents merged with the entire prompt \cite{ref7}. Sequential memory posits a more deliberate approach: the system encodes retrieved knowledge into a sequential cache. When faced with reasoning-heavy tasks that require an association between knowledge and the prompt---including multi-hop reasoning---the pattern of stored memory allows the model to internally process and structure how external data should be used \cite{ref1}. However, advanced memory variants that reinforce useful knowledge via backpropagation face fragility due to many models' reliance on pre-encoded cognitive associations, which critically affects how they consume external knowledge and results in a mix of context and accuracy \cite{ref7}. These methods also consider a middle ground in the broader retrieval trajectories, addressing how the system handles retrieved irrelevant documents through memory reconfiguration, leveraging the 30\% accuracy difficulty on facts highlighted in relevant foundational work \cite{ref10}. As this diversification of input re-configures what counts as the `knowledge source' of the generation stage, sequential models can be trained to use the correct solutions through memory reconfiguration \cite{ref52}. \\
Moreover, the induction of an explicit memory offers resilience when the evaluated model is used in an unsupervised setting. Many base models rely heavily on internal knowledge and provide less attention to the external context, especially when top-k documents are computed with less efficient retrievers \cite{ref51}. External memory induction methods combat this by introducing a module that refines the retrieved context into concise, accurate responses, effectively forcing the reasoning module to act as an information refiner \cite{ref52}. This is particularly important when there are knowledge gaps, as the model outputs plausible but introduces distractors. In domain-specific scenarios such as the financial domain (where high-quality data is highly specialized), memory augmentation takes shape within the embedding space. By utilizing memory layers alongside adapters, the system can optimize the model's embedding attention to enhance retrieval utility, thereby reducing the reliance on general-knowledge models and forcing the generator to focus on the facts available in the evidence \cite{ref3,ref53}. This becomes critical to ensure that the memory is ranked based on utility, not only on relevance or lexical ranking, thereby shifting learning toward high-value data and filtering out noise. Memory here is not only activative but also a form of continuous tuning; a Mafin-like adaptation across embedding models and memory cells yields better retrieval bridges for black-box embeddings into knowledge-augmented scenarios \cite{ref54}. \\
Perhaps most significantly, the latest advancements attempt to build memory systems that learn from user feedback, creating ever-improving research loops. Traditional RAG experiences feedback, albeit largely superficially, through lexical changes in the knowledge base. New frameworks such as RAM echo this direction, which build upon the human pedagogical process in multi-turn interactions by using user feedback signals to continually update the memory embeddings and determining when and how to retrieve. It orchestrates reflection based on experience to ground the model, showing improvements on the robust query types, including false premises and multi-hop queries, without requiring explicit retraining of the original datastore \cite{ref55}. This transformation is complementary to our architectural taxonomies; it signals that the memory module can benefit from direct learning, enabling dynamic shifts in query semantics without re-querying the entire knowledge base. For example, the memory can encapsulate the knowledge that RAG finds, the answer knowledge stored in model parameters, and the process of ever-improving the answer \cite{ref52}, a promising feature that decouples fine-tuning from semantic row structures. \\
Despite these advances, there still exist tradeoffs. The overarching challenge of retrieval accuracy, especially in continuous multi-turn settings, requires the utility-aware storage and calibration of retrieval signals with confidence \cite{ref53}. To prevent hallucination, memory can also be inspired by conformal prediction in RAG: the knowledge that meets an implicit confidence threshold and stored typically ensures that the system provides a guarantee \texttt{(1-alpha)} that the answer is available in the retrieved context given similarity scores \cite{ref47}. For conversational systems that embed user profiles and use memory stores as semantic nodes rather than flat documents, user context is effectively preferable to providing only flat documents, thus reinforcing the memory modules with user-specific dynamics \cite{ref56}. \\
In summary, memory and knowledge augmentation techniques represent a critical pivot in the RAG architecture. They create a proactive and interactive system that evolves with its environment, moving beyond static repositories to integrate dynamic embeddings, nuanced feedback from users, and previous interactions. This allows the system to deftly navigate mixed-knowledge sources and tailor responses to the context of predictable and lifelong learning scenarios, advancing towards the potential of lifelong learning and synchronization in RAG systems \cite{ref1}. \\
\#\# System-Level Optimizations and Efficiency Innovations \\
System-level engineering innovations play a pivotal role in translating the conceptual promise of Retrieval-Augmented Generation (RAG) into practical, scalable, and production-ready systems. While the preceding discussion established the architectural frameworks and memory augmentation techniques that enable dynamic knowledge integration, the operational viability of these systems hinges on addressing substantial computational, memory, and latency bottlenecks inherent to their deployment. These challenges stem not only from the large language models (LLMs) themselves but also from the complex interplay between retrieval processes, long-sequence augmentation, and iterative generation loops. This subsection examines the key system-level optimizations---including caching strategies, pipeline parallelism, speculative retrieval, computational acceleration, and distributed inference frameworks---that are critical for enhancing the efficiency, throughput, and end-to-end latency of modern RAG systems, serving as the essential bridge between the conceptual designs discussed earlier and the effective context utilization strategies that follow. \\
A fundamental performance bottleneck in RAG arises from the sheer volume of the external knowledge base and the computational cost of searching it for every query. As datasets grow into the hundreds of millions of vectors, the retrieval phase can constitute a substantial portion of the overall generation time, especially when periodic retrievals are performed to align newly generated tokens with the latest state of the output \cite{ref57,ref58}. To directly address this, algorithm-system co-design strategies such as the one proposed in PipeRAG advocate for overlapping retrieval and generation steps through pre-retrieval and retrieval pipeline parallelism \cite{ref57}. This approach enables concurrent retrieval operations, where the system fetches documents for the next query while the generating processor concurrently executes the current step, thus improving resource utilization. However, this introduces the challenge of balancing the frequency of retrieval. A crucial insight is that the efficiency of pipelining is tied to the retrieval interval---how often the system pauses generation to fetch more data (e.g., for iterative or multi-step RAG). Too frequent retrieval can lead to serialization of the pipeline, whereas too infrequent retrieval could degrade context freshness in subsequent generation steps. PipeRAG leverages a performance model to automatically balance this trade-off, dynamically setting optimal retrieval intervals based on the current state of generation and the underlying hardware environment, showcasing the power of co-designing hardware-aware algorithms with storage systems \cite{ref57}. \\
This pipeline-level optimization directly informs how retrieved information is subsequently managed within the system, leading to complementary strategies that target the memory side of the equation. The stateful nature of LLM inference means that every time a retrieval corpus is fetched, the intermediate states shared with the personalized RAG query triangle are recomputed, leading to severe memory bottlenecks and low throughput. The challenging problem of caching and reuse of intermediate states is addressed by the multilevel dynamic caching system called RAGCache \cite{ref58}. This work analyses the performance of current RAG systems and identifies the optimization opportunities in eviction strategies and storage in GPU memory and host memory. To that end, RAGCache organizes the retrieved documents and their intermediate states in a specialized data structure called a `knowledge tree', caching them across a hierarchy spanning the fast GPU memory and the larger host memory. An informed memory management policy based on the analysis then dynamically decides which documents to cache at which cached level based on the compiled length of the documents, and this adaptive method compresses the computation of the same document when multiple queries share the same context. This approach significantly reduces the computational expense of recomputing the same documents and therefore improves the latency and throughput. The performance analysis of RAGCache claims to reduce the time to first token generation and increases throughput relative to the baseline using the vLLM and Faiss, indicating how tailored caching mechanisms can be used to directly address the memory side of the problem \cite{ref58,ref59}. \\
Beyond optimizing for generating the same documents repeatedly, the efficiency of the entire generation workload can be approached from the angle of computational acceleration and distributed fine-tuning. The efficiency improvements of computational acceleration and the efficient fine-tuning of small models to perform RAG directly can reduce both cost and carbon footprint. The system-level optimizations for distributed training that leverage just-in-time compilation and tensor sharding, such as in the case of JORA, demonstrate the potential of reducing the resources needed for fine-tuning LLMs for specific RAG tasks. However, the broader lesson from the work is the importance of directly mapping the workload onto the distributed system resources \cite{ref60}. These methods lower the barrier for deploying RAG models on systems with limited resources, by efficiently executing parts of the LLM across multiple GPUs and reducing memory overhead, achieving fine-tuning that is significantly faster than fully-connected systems and enabling wider adoption of RAG strategies \cite{ref60}. \\
While these caching and pipeline strategies directly optimize for latency, the broader computational memory footprint of scaling RAG with massive combinations of data and model requires efficient partitioning across available hardware. In the context of `accelerator-level parallelism' (ALP), the deployment of a heterogeneous set of specialized accelerators to simultaneously execute parts of a computationally intensive pipeline can become a central theme for RAG workloads \cite{ref61}. For RAG systems, this can mean that the embedding computation for retrieved data is executed on one dedicated accelerator while the core LLM inference is handled by a high-throughput processor on another, or that embedding models and generation ensembles are spread across available units. This line of work has shifted focus towards the orchestration and pragmatic use of the available computational units, emphasizing the need to co-design the algorithm and the hardware to maximize throughput. \\
Despite the theoretical merits of these optimizations, the reality of productionization reveals that are often evaluated on single a\emph{components} and that end-to-end behavior must be holistically validated. One of the key \textbf{Seven Failure Points} in engineering RAG systems is the grave failure to account for the dynamic and slowly evolving nature of the underlying data, and the fact that system-that would only be ``optimal'' at design-time may not hold throughout the system's lifecycle \cite{ref59}. In this regard, the design of robust systems should account for shifts in the data distribution, model update constraints, and hardware heterogeneity. The issue is not simply optimizing a single component, but rather building a system that can dynamically discover the correct configuration. By having a flexible system that can select between different retrieval and generation strategies in real time---similar to a \emph{portfolio-based architecture} seen in parallel solvers---the RAG system can adapt and maintain performance under a diverse set of workloads \cite{ref62}. When we consider the need to operate inside context, we observe that simply maximizing the number of documents provided does not always improve performance for all models; indeed, decoder-only models have been shown to not fully utilize long-context windows and saturate in performance after a certain number of input documents, a factor to consider when tuning the system \cite{ref7}. \\
In a larger scale, techniques from distributed systems, such as those used in SMURF, reveal that optimizing for metadata access---not just the data itself---is key to enabling fast distributed tasks, as the index metadata can become the bottleneck as the knowledge graph grows \cite{ref63}. However, an important caveat emerges: the physical implementation of the memory system can reduce the benefits of purely logical caching decisions, as the performance of the cache is coupled to the distribution of resources, as seen in cache coherence traffic and the inherent difficulty in controlling latency in many-core systems \cite{ref63,ref64}. Looking at the challenges from a distributed systems perspective and the problem of skill specialization, the success of RAG in a production setting requires that the system design considers the mitigated interaction between the data, the model, and the hardware. \\
In conclusion, system-level optimization for RAG is about radically deconstructing the RAG pipeline to map its computational stages onto the underlying capabilities of specialized hardware and orchestration schedulers. The key innovations that are providing breakthroughs in performance include efficient caching and reuse of intermediate computations, algorithm-system co-design with intentional scheduling of retrieval and inference by the model, and the distribution of model replicas and data across heterogeneous devices. As RAG applications scale in text length, in the complexity and scale of the knowledge base, and in user volume, hardware-aware reasoning over not just the task at hand---such as the probability of a retrieval task driving the scheduling of next steps---becomes more crucial. By treating retrieval, generation, and the ingestion pipeline as parts of a single, unified system view, it is possible to minimize the end-to-end latency and compute, enabling the practical and responsible use of a RAG architecture. This foundation sets the stage for the subsequent discussion on how effectively leveraging the retrieved contexts---determining optimal context granularity, implementing relevance-aware reranking, and controlling the generator---can build upon these performance-critical optimizations to achieve both efficient and high-quality RAG systems. \\
\#\# Context Utilization and Relevance-Aware Retrieval \\
The effective utilization of retrieved contexts is a cornerstone of Retrieval-Augmented Generation (RAG), determining whether augmented knowledge translates into superior generation or introduces noise and confusion. Building upon the system-level optimizations for efficient retrieval and inference discussed previously, this subsection shifts focus from \emph{how} to deliver context efficiently to \emph{what} context is delivered and \emph{how} the generator can best exploit it. We examine research that clarifies contextual influence, the number of documents a generator should consider, strategies for relevance-aware retrieval and reranking, and the mechanisms through which generators can be finely controlled to leverage retrieved evidence effectively. \\
A fundamental yet often underappreciated aspect of RAG is determining the optimal amount and granularity of context to supply to the language model. Feeding too many or too few documents can degrade performance, making the study of document count effects crucial. The relationship is frequently non-monotonic: while some tasks benefit from additional evidence, others suffer from attention dilution or the introduction of distracting noise. Insights into this phenomenon can be gleaned from research on how models handle varying amounts of textual information. For instance, when grappling with long documents, traditional Transformers face prohibitive computational costs, and a common remedy involves truncation to the first few tokens, a strategy that often introduces bias and underretrieves longer documents. As demonstrated in literature exploring local self-attention mechanisms, models can process thousands of tokens with a fraction of the memory and compute cost by attending to moving windows, without disproportionately favoring shorter texts \cite{ref65}. This efficiency consideration has direct implications for RAG system design, suggesting that the chunking strategy and maximal input length of retrieved documents are critical design choices that influence the comprehensiveness and balance of the context the generator can digest. \\
Building upon the efficiency-driven context selection, RAG systems must also proactively improve the relevance and fidelity of what they retrieve---an area generally referred to as relevance-aware retrieval and ranking. When naive retrieval results are insufficient, a common approach is to leverage reranking pipelines to enhance precision and filter out irrelevant information. Traditional lexical retrievers can efficiently fetch an initial candidate set, but their effectiveness is limited. To bridge this gap, research has introduced sophisticated pseudo-relevance feedback mechanisms that enhance dense retrieval. One notable method, ColBERT-PRF, extracts representative feedback embeddings from an initial set of pseudo-relevant documents using KMeans clustering, ensuring the clusters are discriminative via IDF. This technique improves MAP by up to 26\% and 10\% on TREC 2019 and TREC 2020 collections, respectively, demonstrating that enriching the query representation can substantially correct initial retrieval errors \cite{ref66}. However, building or retrieving from a knowledge base is only part of the issue; the system must also adapt itself to the unique task of locating source citations for a final text. \\
When reranking for the purpose of source attribution (e.g., locating Wikipedia articles used to write an essay), the type and amount of supervision can affect the quality of the reranking. Effective reranking is not solely a function of the model size, but also the nature of supervision, where semi- and fully-supervised methods can be elicited from coarse document-level labels rather than expensive span-level annotation \cite{ref67}. Similarly, an approach called COntextual Document Embedding Reranking (CODER) demonstrates that incorporating the ranking context---such as using query-specific negatives and list-wise losses---can lead to substantial retrieval improvements \cite{ref68}. These tests point to a general principle that models which consider the ranking context, rather than only isolated documents, exhibit stronger zero-shot generalization and effectiveness. Interestingly, the benefits of scaling model parameters on reranking versus dense retrieval are modest in-domain, but rerankers significantly outperform dense retrievers in zero-shot settings \cite{ref69}. This suggests that early query-document integration is essential for adapting to novel inputs, even with compact representations. \\
To bridge the gap between user needs and the explicit query---especially when operating in an interactive setting, relevance feedback methods provide a way to refine the system's understanding. In information-seeking scenarios, feedback signals (clicks or explicit ratings) can be integrated into the document-candidate rerankers via few-shot techniques. A kNN-based reranker that leverages similarity to positive feedback documents can be fused with a lexical ranker, showing consistent improvements \cite{ref70}. This highlights that queries alone are often ambiguous, and the ability to integrate interactive signals is a promising direction for future RAG systems. \\
Finally, the logic of relevance-aware retrieval can be extended upstream into the generation controller itself. A common strategy is to use generative retrieval models that encode pointers to information in their parameters and then decode page titles as a retrieval step. However, this pipeline does not fully account for the contextual relevance of the retrieved information to the user's query. This limitation is addressed by the model called Re3val, which utilizes Dense Passage Retrieval to acquire related contexts, then applies a generative relevance encoder and, through reinforcement learning, optimizes the ranking of relevant page titles at the token level \cite{ref71}. Beyond parameterized controllers, the model can also better leverage the context itself. Through hierarchical chunking methods, a retriever can break long documents into smaller sub-chunks and then allow the query to attend to all the chunk encodings simultaneously, enabling a global interaction across the entire source document even when indexing is local \cite{ref65}. This mirrors the principle that users require a holistic understanding of the context when the retrieved evidence is deconstructed into smaller units. A further critical dynamic is the mismatch between what a snippet provides and what the model or human actually needs. A study shows that long, diverse, and information-rich queries benefit more from full in-context documents, whereas navigational and short queries may not, highlighting that the optimal context format depends on the intent and ambiguity of the user's request \cite{ref72}. \\
In summary, effectively leveraging retrieved contexts in RAG moves far beyond a simple BM25 lookup followed by a prompt. It involves dynamically determining the optimal context granularity and size, reweighting evidence to discriminate against irrelevant signals, and conditionally guiding the generator. Promising future directions include integrating these methods into robust architectures that leverage mixed human input and adaptive, quality-aware reranking strategies to prevent noise from dominating the final synthesized, high-fidelity response to the user. \\
\# Training, Optimization, and Augmentation Strategies for RAG Systems \\
\#\# End-to-End and Component-Specific Training Paradigms \\
The development of effective Retrieval-Augmented Generation (RAG) systems hinges critically on the training paradigms employed to optimize their dual-component architecture. Unlike traditional language models trained solely on static corpora, RAG frameworks require coordinated optimization of the retriever---which must locate relevant information---and the generator---which must synthesize that information with its parametric knowledge. While the preceding discussions have established the architectural foundations and component-wise designs of RAG systems, this subsection shifts focus to the pivotal question of how these systems are trained, exploring the full spectrum of strategies that span from holistic end-to-end fine-tuning to component-specific and even black-box adaptations. In doing so, it sets the stage for the subsequent examination of how reinforcement learning and contrastive methods refine the retrieval--generation alignment, ultimately revealing a progressive transition from static, decoupled training workflows toward dynamic, co-adaptive optimization loops. \\
At the foundational level, the original RAG architecture introduced a paradigm where a pre-trained seq2seq model acts as the parametric memory and a dense vector index of Wikipedia serves as non-parametric memory \cite{ref73}. While this formulation proved superior to parametric-only baselines on knowledge-intensive tasks, early implementations frequently froze the retriever and generator components, training only a subset of parameters or none at all. Recognizing this limitation, subsequent research demonstrated that fine-tuning the entire RAG architecture in an end-to-end manner---jointly optimizing the retriever and generator---yields significant improvements over the frozen base model for question-answering tasks. This end-to-end approach, while encountering notable engineering challenges related to gradient propagation and memory constraints, allows the retriever to learn query representations that better align with what the generator actually optimizes for, creating a more cohesive system \cite{ref6}. The joint training paradigm, however, is computationally intensive, requiring careful attention to loss combination, and often necessitates the use of scale instruments like JAX's tensor-sharding to distribute parameters and data across multiple GPUs, as demonstrated by libraries such as the JAX Tensor-Parallel LoRA library for retrieval-aided fine-tuning \cite{ref60}. \\
Complementing complete fine-tuning, the broader literature has evaluated a comparison of two common methods for injecting knowledge into LLMs: unsupervised fine-tuning (via next-token prediction) and retrieval-augmented generation. Empirical findings reveal a decisive advantage for RAG across a variety of knowledge-intensive tasks. Fine-tuning, particularly unsupervised, struggles to effectively integrate entirely new factual information into the model's frozen weights and exposes a tendency to overfit or fail to generalize synthesized variations of the same fact \cite{ref74}. This finding reinforces the value of RAG as a low-cost solution for keeping knowledge up-to-date without the risk of catastrophic forgetting. As the prohibitive cost of updating model parameters in large-scale, high-capacity models becomes increasingly apparent, a parallel strand of research has pivoted toward black-box RAG, where a large language model (LLM) is treated as a fixed, frozen component, with optimization placed on the retrieval side or on lightweight, trainable add-ons \cite{ref75}. For instance, an approach called Mafin targets scenarios where embedding models are themselves trained and served via a black-box API. Mafin enhances the black-box embedding by augmenting it with a small trainable embedding model, which is trained on a downstream task, effectively fine-tuning the information access mechanism without directly modifying the inaccessible, high-capacity retrieval backbone \cite{ref54}. This model augmented fine-tuning offers a practical path for domain adaptation of the retrieval stage when the embedding service is closed-sourced. \\
Beyond the choice between full end-to-end tuning and black-box adaptation, a critical distinction emerges in how parameter-update strategies are decomposed across the retriever and generator. Component-specific fine-tuning provides a practical and often more efficient middle ground. For many knowledge-intensive applications, a high-performance retriever is the single most important factor in improving generation quality; RAG systems that fine-tune only the retriever to align with a canned LLM's preferences can yield large gains, even when the generator remains frozen \cite{ref75}. Similarly, parameter-efficient fine-tuning technique---particularly Low-Rank Adaptation (LoRA)---has emerged as a standard for deploying RAG pipelines, since it dramatically reduces memory footprint while enabling multi-GPU parallel training, thereby making the fine-tuning of LLMs with long retrieved contexts tractable \cite{ref60}. This includes library development that uses just-in-time compilation mechanisms to improve training speed while lowering VRAM consumption, emphasizing the importance of co-designing system-level considerations with algorithmic progress. Complementing these adaptation methods, Informational Refinement Training (InFO-RAG) offers an unsupervised approach that trains the LLM as an information refiner to process and compress augmented input documents into more concise and complete formats suitable for generation, regardless of the initial retrieval purity \cite{ref52}. \\
An equally important axis of the training landscape is the trade-off between fine-tuning and RAG, which becomes particularly pronounced when considering data popularity and user access. Baseline studies with minimal programming overhead (e.g., unmodified GPT-3.5) show that fine-tuning models for general-purpose QA tasks underperforms compared with using a vectorized RAG database, particularly when answers are current and post-training. Combining RAG with a simple soft prompt boosts performance further \cite{ref76}. However, this comparative advantage of RAG is heavily dependent on the popularity of the knowledge. When applied to popular or well-represented topics, fine-tuned models may show more significant improvements, whereas RAG typically proves superior for less popular knowledge coverage \cite{ref77}. These observations point to a flexible design space: system designers can and should mix-and-match tuning options to align with the target domain and the underlying knowledge distribution. This idea is further amplified by domain-adapted RAG systems such as those that jointly train the retriever and generator on domain-specific corpora like COVID-19, News, or conversational data, demonstrating performance gains in the face of external knowledge and auxiliary losses \cite{ref78}. \\
In summary, the landscape of RAG training is an exciting, multi-faceted spectrum where the choice among full parameter fine-tuning, component-specific tuning, or black-box augmentation is dictated by the use-case's individual requirements---particularly the demand for continuously updated, low-frequency knowledge. While black-box augmentation techniques like Mafin \cite{ref54} show effective performance without extensive computational overhead, there is no universal solution. Rather, successful deployment often involves mixing and matching fine-tuning strategies to robustly align the retriever and generator in the presence of inherent noise and uncertain retrieval quality. This co-adaptation philosophy directly foreshadows the more advanced techniques discussed in the following subsection, in which reinforcement learning and contrastive learning are applied to close the loop between retrieval and generation feedback, further transforming the optimizer from a one-time fine-tuning stage into a continuous, interactive process. \\
\#\# Reinforcement Learning and Contrastive Learning for Alignment and Optimization \\
Reinforcement learning (RL) and contrastive learning have emerged as two powerful paradigms for optimizing retrieval-augmented generation (RAG) systems, particularly for aligning the retrieval component with the ultimate generation objectives. Building upon the training strategies discussed in the previous subsection---where end-to-end fine-tuning, component-specific adaptation, and black-box augmentation were presented as distinct optimization philosophies---this subsection introduces the next logical evolution: closed-loop, feedback-driven optimization that dynamically co-adapts the retriever and generator. While traditional RAG pipelines treat the retriever as a relatively static module that can be trained with supervised signal or left frozen, the more sophisticated frameworks recognize that the retriever's behavior must be directly optimized according to downstream task requirements, which often hinge on the idiosyncrasies of the large language model (LLM) employed. This subsection explores the specific methodologies---spanning fine-grained RL feedback, supervised and RL chained bridge models, and contrastive alignment strategies---that researchers have developed to close the optimization loop between retrieving and generating, transforming the one-time training procedures described earlier into a continuous, interactive process. \\
A central challenge in leveraging RL for retrieval is the black-box nature of modern LLMs; since many frontier models are only accessible via APIs and do not expose their internal parameters or token probabilities, the reward signal for training the retriever cannot be derived through a straightforward end-to-end backpropagation. The approach of Reinforcement Retrieval Leveraging Fine-grained Feedback for Fact Checking News Claims with Black-Box LLM (FFRR) directly addresses this problem by introducing a two-level, fine-grained feedback strategy \cite{ref79}. In the FFRR methodology, the black-box LLM is treated as an evaluator, generating ratings for each retrieved document based on the non-retrieval ground truth of the fact-checking task. These ratings serve as a dense reward signal that is used to optimize the retrieval policy through RL. This approach demonstrates a notable departure from merely optimizing for lexical or semantic similarity---it explicitly aligns the retriever's preferences with the downstream judgment behavior of the LLM, aligning with the black-box adaptation philosophy introduced earlier while extending it from a one-time fine-tuning step to a continuously optimized policy. \\
Critically, the relationship between the retriever and the LLM is not mechanical---modern work has shown that a retrieved document that is human-friendly may not be in the format, order, or granularity that the LLM finds most useful. The bridge model paradigm addresses this by chaining supervised and reinforcement learning methods to directly connect the retrieve and generate components \cite{ref80}. In this approach, the retrievers are validated by their ranking assumptions while the LLM assembles the augmented context. The gap between them is bridged by training a separate---but jointly optimized---model. Supervised learning is first used to teach the bridge model to mirror the retriever's rankings, while reinforcement learning is used to optimize the bridge to produce a context sequence that leads to a higher reward from the generation task. This explicit chaining means that the bridge model transforms the raw retriever output into an LLM-favorable context, a mechanism that conceptually complements the end-to-end and component-specific fine-tuning strategies discussed earlier, yet goes beyond them by introducing a learned intermediary that is explicitly optimized to bridge the retrieval--generation preference gap. This approach results in substantial gains in both QA and personalized generation settings. \\
A different but complementary line of work proposes the use of generative adversarial feedback to improve alignment between retriever and generator. The study in Reinforcement Learning with Generative Adversarial Feedback (RLGAF) uses this adversarial style to enable LLMs to learn useful expert demonstrations without explicit exposure to expert-labelled data \cite{ref81}. Through a discriminator distinguishing between model outputs and expert-like preferences, the approach provides a promising context for RAG optimization, reinforcing the theme of dynamic co-adaptation by introducing an iterative, adversarial feedback loop that continuously sharpens the retriever and generator. \\
While RL provides a direct optimization route for retrieval, it also introduces complications, especially around training stability and credit assignment. This complexity motivates the careful consideration of contrastive learning as a complementary method. Contrastive methods offer a structure to directly match the distributions of retrieved states to the generation expectations without relying solely on policy gradient methods. Work on Reinforcement Learning with Contrastive Feedback (RLCF) demonstrates such a structure: it constructs unsupervised contrastive feedback based on similar document groups to train reward functions that encourage the LLM to generate highly relevant, distinct responses \cite{ref82}. The reward function---denoted by Group-wise Reciprocal Rank---is used in the reinforcement learning with Proximal Policy Optimization, delivering substantial improvements in performance and value over previous baselines. \\
Another example, the Contrastive Learning Framework for Human Alignment (CLHA), adopts pair-wise learning with adaptive supervised fine-tuning to dynamically adjust likelihoods and enhance alignment with human preferences, using contrastive signals and data-driven rescoring for a stable and effective quality \cite{ref83}. Although less directly tied to the retrieval step, this approach improves model alignment overall. In a RAG setting, a model aligned through such contrastive signals can integrate retrieved contexts into generation more robustly, thereby complementing the retriever-focused RL methods and contributing to the broader co-adaptive design of the entire pipeline. \\
The idea of contrastive penalties has also been explored to counter the inherent noise and uncertainty in reward modeling. The work Improved Robustness in RLHF Using Contrastive Rewards explicitly adds a penalty to the reward signal to increase error tolerance and variance control \cite{ref84}. This adjustment improves reinforcement learning convergence and strengthens the overall RAG foundation, particularly in the noisy, real-world retrieval environments anticipated in the following subsection. Furthermore, graded Reinforcement Learning from Reflective Feedback (RLRF) shows how fine-grained criteria and self-reflective processes can move RL-based alignment beyond surface-level adjustments and affect core decisions on information retrieval and focus \cite{ref85}. This implies that the retrieval module in RAG pathways can benefit from reflective-loop optimization, enhancing the robustness of the co-adaptive design in the presence of retrieval noise and providing the stable foundation for the efficient deployment strategies examined next. \\
In sum, the recent research in RL and contrastive methods in RAG spaces marks a decisive progression from static retrieval toward a dynamic, co-adaptive RAG optimization paradigm. FFRR demonstrates how black-box LLM feedback can refine retrieval policy \cite{ref79}. Bridge models explicitly close the retriever--generator gap \cite{ref80}. RLCF and CLHA contribute via contrastive alignment strategies to strengthen the context of the language model outputs. Together, these approaches resonate with the earlier training paradigms by reinforcing the core message: a unified RAG system requires joint optimization of its components rather than independent design. Yet, these RL-driven methods add a crucial new dimension---the introduction of an interactive reward against which a system continuously monitors itself. For dynamic, real-world applications, feedback-rich, closed-loop optimization becomes a necessity rather than a nice-to-have. However, these computational overheads introduced by such iterative, reward-driven tuning further underscore the need for the parameter-efficient fine-tuning and system-level co-design discussed next, ensuring that the co-adaptive, RL-driven RAG systems are both powerful and operationally viable. \\
\#\# Parameter-Efficient Fine-Tuning and System-Level Co-Design \\
The deployment and customization of Retrieval-Augmented Generation (RAG) systems at scale are critically dependent on the computational and memory efficiency of the underlying large language models (LLMs). While the previous subsection examined how reinforcement learning and contrastive methods align retrievers with generators, such co-adaptive optimization is only practical if the underlying model can be efficiently updated and served. Full fine-tuning of these models---often containing billions of parameters---becomes prohibitively expensive in terms of GPU memory, training time, and storage, especially when multiple models must be adapted for different tasks or domains. As a solution, parameter-efficient fine-tuning (PEFT) has emerged as a cornerstone methodology, and its integration with system-level co-design has become essential for building practical and scalable RAG pipelines. This subsection explores the landscape of PEFT methods---particularly low-rank adaptation (LoRA) and its variants---and how system-level innovations, including caching, compression, and parallelism, are designed to overcome the computational challenges inherent in RAG. These efficiency gains are not merely operational conveniences; they provide the computational substrate that enables the dynamic, co-adaptive retrieval--generation loops described earlier, while also ensuring that RAG systems remain reliable and deployable in the noisy, resource-constrained environments anticipated in the following subsection. \\
The premise of PEFT is to achieve performance comparable to full fine-tuning by updating only a minimal subset of parameters while freezing the bulk of the pre-trained model \cite{ref86}. This approach offers significant reductions in GPU memory consumption during training and often mitigates the issue of catastrophic forgetting, thereby better preserving the general-purpose knowledge of the base model. The resulting adapter weights are also considerably smaller to store and communicate. For instance, rather than storing a full copy of a model for each task, PEFT can produce per-task adapter weights that are orders of magnitude smaller. The compression and communication of these expert models can be further optimized through techniques like sparsification and quantization, enabling more efficient retrieval over high-latency networks and facilitating on-the-fly model composition \cite{ref87}. ComPEFT, which projects PEFT updates into a sparse and ternary form, can achieve high compression ratios while preserving or even enhancing performance, especially in large models \cite{ref88}. \\
At the heart of PEFT approaches lies Low-Rank Adaptation (LoRA), which introduces trainable low-rank matrices into the layers of a frozen backbone model. The central hypothesis is that the weight updates during adaptation have a low intrinsic rank. However, this assumption creates a low-rank bottleneck that can limit effectiveness. Research has sought to transcend this limitation. For instance, PeriodicLoRA accumulates update matrices over multiple training stages by unloading the accumulated weights back into the frozen model at the end of each stage. This tactic effectively increases the rank of the update lifting the rank bottleneck without increasing memory usage during training, yielding a stronger learning capacity \cite{ref89}. The standard LoRA framework also uses a scaling factor inverse to the rank. A theoretical analysis of the learning dynamics reveals that a more appropriate scaling factor for low-rank adapters should be inversely proportional to the square root of the rank instead. Using this corrected scaling factor with rank-stabilized LoRA maintains the learning of high-rank adapters stable, enabling a fine-tuning compute performance trade-off where performance can improve with rank without a corresponding increase in inference cost \cite{ref90}. \\
In multi-task and multi-domain scenarios, a static adaptation rank is often suboptimal, effectively fixing the complexity of the fine-tuning task to a predetermined rank. ALoRA addresses this by dynamically allocating the low-rank budget across the model. It uses an importance estimation method to prune less influential low-rank updates and then reallocates the freed-up rank to transformer layers that require a higher rank \cite{ref91}. This dynamic and adaptive allocation allows the model to flexibly adapt to each task, often matching or improving upon the performance of static-rank baselines. Similarly, the concept of Mixture of Experts (MoE) can be applied to LoRA as well. By treating each LoRA adapter as an ``expert'', novel architectures have been developed to serve multiple LoRA adapters, as explored in Mixture-of-Experts frameworks like MixLoRA and MoELoRA. MixLoRA creates a resource-efficient sparse MoE model by inserting multiple LoRA-based experts into a frozen dense model and a top-k router. This approach enhances multi-task performance while requiring fewer parameters than an equivalent MoE model, making it possible to train multiple experts on a small-scale GPU with 24GB of VRAM \cite{ref92}. The introduction of routing functions into the low-rank layers has also proven crucial, as it efficiently aligns multimodal representations without introducing trainable parameters \cite{ref93}. Moreover, allocating more experts to higher layers tends to be more crucial for model capacity, yielding better parameter efficiency than a uniform distribution of experts \cite{ref94}. This MoE paradigm enhances the performance and generalization of adapters, especially in multi-task RAG pipelines. \\
While these algorithmic innovations improve adaptation quality, they also introduce system-level complexities, particularly around memory management, scheduling, and serving. The efficiency of PEFT in large-scale system design becomes paramount. Beyond algorithmic advancements, system co-design is essential. State-of-the-art serving libraries for large models need to be extended to simultaneously support PEFT methods. For instance, S-LoRA addresses the challenge of serving thousands of LoRA adapters at scale by introducing a unified memory pool that efficiently stores adapter parameters and KV cache tensors, using a paging mechanism to avoid memory fragmentation. The introduction of customized CUDA kernels and a novel tensor parallelism strategy for heterogeneous batching of LoRA computation allows S-LoRA to efficiently scale the serving of many adapters \cite{ref95}. Methods like MixLoRA similarly demonstrate that system-level efficiency in terms of GPU memory consumption and latency can be substantially reduced while still serving multiple models \cite{ref92}. \\
System-level co-design also involves cost-effective compression techniques that allow PEFT methods to be delivered efficiently across networks. For example, ComPEFT compresses the task vectors (fine-tuning residuals) using sparsification and ternary quantization \cite{ref88}. It has demonstrated that compressed models can maintain strong performance and are easier to store, making them particularly relevant for fine-tuning and few-shot learning in dynamic RAG applications that require rapid, on-demand model updates \cite{ref88}. Likewise, algorithmic selection between LoRA and alternative methods such as Representation Editing---which uses scaling and biasing instead of low-rank decomposition---can reduce the number of trainable parameters by 25,700 times compared to full fine-tuning and 32 times compared to LoRA, while achieving strong performance across a variety of tasks \cite{ref96}. \\
In conclusion, the success of PEFT at scale is not the result of a singular algorithmic breakthrough but emerges from the close interaction between novel neural network architectures and underlying systems. The evolution from LoRA to its diverse variants (PeriodicLoRA, ALoRA, MixLoRA) demonstrates the continuous push toward breaking performance bottlenecks while maintaining parameter efficiency. Meanwhile, system-level design across the continuum---from memory organization \cite{ref95} to weight compression \cite{ref88} and efficient adapter serving \cite{ref92}---ensures that these techniques can be leveraged at the scale demanded by RAG systems. This holistic approach to efficiency is crucial for deploying RAG in domains requiring lower latency and higher throughput, providing the operational foundation that makes the co-adaptive and RL-driven optimisation discussed in earlier RL-based methods feasible in practice, and laying the groundwork for the robust, noise-tolerant systems required in the next subsection. \\
\#\# Robustness to Noise, Counterfactuals, and Unsupervised Adaptation \\
The robustness of Retrieval-Augmented Generation (RAG) systems is a critical determinant of their reliability in real-world deployments, where retrieved corpora are rarely clean, perfectly relevant, or free from adversarial interference. The documented fragility of these systems---where even minor textual errors can ``break the RAG's back''---underscores the need for dedicated strategies that fortify the entire pipeline against noise, counterfactuals, and distribution shifts \cite{ref97}. While the previous subsection established parameter-efficient fine-tuning (PEFT) as the computational substrate for scalable, co-adaptive RAG systems, this subsection addresses a complementary imperative: ensuring that these systems remain trustworthy when the retrieved context itself is imperfect, a challenge that becomes even more acute as we transition to the dynamic, memory-augmented architectures discussed in the following subsection, where feedback loops and iterative retrieval inherently amplify the impact of noisy inputs. This subsection synthesizes approaches that enhance RAG robustness along three interlocking dimensions: noise-resistant training, counterfactual data augmentation, and unsupervised adaptation. \\
\textbf{Noise-Resistant Training.} The first line of defense involves imbuing the retriever--generator pipeline with inherent tolerance to imperfect context. Early work in adversarial robustness provides foundational intuitions for RAG: models trained exclusively on clean data tend to be overconfident and catastrophic failures when faced with subtle perturbations, often exhibiting a pathological ``catastrophic overfitting'' where rapid gains in robustness are undone by severe degradation on clean samples \cite{ref98}. To counteract this, adversarial training and its variants have been repurposed as building blocks for robust RAG. For example, \emph{ensemble adversarial training} and parameter-efficient methods like \emph{parameter-saving adversarial training} could be adapted at the embedding level to harden retrieval encoders against typo-less and low-level perturbations \cite{ref99,ref100}. Critically, the inherently multi-component nature of RAG---retrieval, augmentation, and generation---mirrors the multi-perturbation scenario well-known in adversarial learning, and hypernetwork-based approaches that effectively handle diverse attack surfaces suggest a blueprint for jointly protecting retrievers and LLMs \cite{ref100}. \\
The concept of instance-specific difficulty is equally relevant to RAG. Models trained with uniform margins for all training samples struggle to generalize to unperturbed data, and the others directly affect the read of a RAG's generator when confronted with an intermingled set of relevant and irrelevant documents \cite{ref101}. By adopting the principle of \emph{instance-adaptive} margins, we can advocate for RAG-specific training that uses dynamically allocated noise budgets---this is further reinforced by \emph{blind adversarial training}, where a cutoff-scale strategy avoids the extremes of overestimation and underestimation of adversarial noise \cite{ref102}. Using these insights, robust RAG systems can avoid both overfitting to noisy tokens and over-reliance on improbable clean contexts. \\
\textbf{Counterfactual Data Augmentation.} Counterfactuals are crucial for teaching RAG models to reason about what is not bound; the cascade of AI failures caused by small perturbations can be interpreted as an over-reliance on lexical and syntactic cues \cite{ref97}. The \emph{learning to learn from mistakes} framework introduces a meta-optimizer that learns robust loss landscapes in low-data regimes, which can be applied to generating adversarial training examples for RAG without requiring a large budget of labeled data \cite{ref103}. Data augmentation itself---swapping, cropping, composition---when combined with weight averaging is shown to provide significant gains in robustness against the sorts of realistic document errors that create distribution shifts in practice \cite{ref104}. For RAG, this translates into constructing counterfactual pairs where a plausible but wrong passage is inserted into the retrieval set. These counterfactual documents must be semantically close to the correct answer but entail opposite semantics, forcing the generator to sharpen its preference for verified, entailment-consistent passages. \\
One particularly effective counterfactual leveraging the multi-attack robustness paradigm comes from \emph{learning to generate noise for multi-attack robustness}---a meta-learning framework that learns to generate a range of perturbations that jointly equate to realistic document noise. This approach suggests that efficient RAG-specific adversarial training can be extended to generate a spectrum of low-level, multi-modal noise that simulates diverse distribution shifts in real-world corpora, aligning well with the adaptive, multi-modal robustness requirements of the previously discussed PEFT \cite{ref105}. \\
\textbf{Unsupervised Adaptation.} The third dimension of robustness is that the system is able to operate robustly under distribution shift with little to no labeled data. RAG deployments consistently encounter new domains and emerging knowledge bases where no supervision is available. We have learned from domain adaptation that the \emph{Fourier adversarial attacking} method can introduce large-magnitude perturbations with semantic minimal modification to improve internal representations, and this approach can be applied to the retrieval and generation stages of RAG to maintain semantic faithfulness under noisy or out-of-domain inputs \cite{ref106}. In a RAG context, Fourier-style perturbations can be used to augment a sentence to produce retrievable passages that remain consistent with the original semantic, while the generator is forced to manage for semantic invariance, thereby improving domain robustness. \\
Another important unsupervised avenue borrows from contrastive and speech-based methods: \emph{multi-task continual pre-training} that jointly denoises via auxiliary reconstruction and learns robust representation, which can be used to make the embedding space of a retriever invariant to noise \cite{ref107}. An analogue for RAG would be a system where the retriever is trained using a reconstruction task for raw passages ; the reconstructive task enforces noise-invariance in the embedding space, and the resulting robust embeddings are retrieved for effective nearest-neighbor search---All without explicit domain labels. Finally, the addition of \emph{blinding adversarial training} heuristics as an unsupervised wrapper ensures that the magnitude of the perturbation is adaptively adjusted during the fine-tuning, so that the system maintains a robustness point that aligns with the actual attack distribution at runtime, which complements the dynamic resource constraints inherent to PEFT-based systems \cite{ref102}. \\
In summary, true robustness for RAG systems is not a single methodology but a composite discipline; it requires removing token-level/typo-level flaws in the retriever and generator (noise-resistant training), builds in exposure to diminish context misalignment (counterfactual augmentation), and retains reliability across unlabeled shifts (unsupervised adaptation). From architecture-aware to low-level counterfactual augmentation to unsupervised adaptation in noisy domains, these collective strategies prevent the cascade of failures observed in empirical settings---where even a single typo or malformed prefix can trigger systemic failures---but are overcome when explicitly prepared for those disturbances by the training scheme \cite{ref97}. Crucially, these robustness measures are not isolated patches; they must work in harmony with the efficiency techniques described in the previous subsection, ensuring that performance gains do not come at the cost of fragile retrieval capabilities. As we now turn to dynamic, feedback-driven memory and self-correction mechanisms in the next subsection, the robust foundations established here become the prerequisites for systems that can incrementally learn and adapt from their own mistakes---a necessity often more so, when the external environment itself is evolving. \\
\#\# Integrating Memory, Feedback Loops, and Self-Improving Systems \\
The integration of external memory, feedback loops, and self-improving mechanisms represents a critical frontier in the evolution of Retrieval-Augmented Generation (RAG) systems. While foundational RAG architectures rely on static, one-time retrieval from external corpora, contemporary research increasingly focuses on endowing these systems with dynamic capabilities that emulate human learning processes---allowing models to accumulate knowledge over time, adapt to user preferences, and iteratively refine their own retrieval and generation strategies. This subsection examines how training paradigms and system-level designs are being reshaped by three interconnected dimensions: external memory induction, reinforcement from interactive feedback, and self-improving architectures. These dimensions not only bridge the gap between retrieval quality and generation performance but also extend naturally from the robustness considerations discussed previously---where systems must withstand noisy inputs---to the domain-specific adaptations that follow, which demand continuous learning from specialized corpora and user interactions. \\
The incorporation of external memory layers into RAG frameworks represents a fundamental shift from the single-pass retrieval paradigm. Traditional retrieval methods that augment LLMs with relevant knowledge on a per-query basis are often limited in their ability to learn from past successes or failures. To address this, systems like Auxiliary Rationale Memory for Retrieval Augmented Generation (ARM-RAG) introduce a persistent, external memory that stores reasoning chains from previous problem-solving attempts \cite{ref108}. This approach instantiates a dynamic memory layer that is independent of the LLM's parametric weights, demonstrating that storing and subsequently retrieving coherent chains of reasoning---rather than just isolated documents or facts---can positively influence performance on tasks such as grade-school math problems \cite{ref108}. The implication of such memory modules is that they enable frozen or black-box LLMs to effectively ``learn from experience'' without the substantial computational costs associated with fine-tuning. Thus, the retrieval process transcends a simple search over a static knowledge base, becoming a targeted retrieval from an accumulated knowledge base that grows and evolves with every interaction. \\
Expanding on the concept of memory, as highlighted in the robust adaptations discussed previously where noise resilience is key, frameworks such as RAM (Retrieval-Augmented Memory) push this idea further by simulating a pedagogical process in which the system continually updates its memory through interactive communications \cite{ref55}. RAM's innovative approach trains a RAG-based framework to use recursive reasoning-based retrieval and experiential reflection. The system uses communications---often in the form of clarifications or corrections---to actively refine its stored knowledge. This is a significant departure from traditional static corpora, as the system is no longer purely dependent on a fixed retrieval pipeline. By actively learning from user instructions and communicative feedback, the framework outperforms traditional RAG and self-knowledge generation methods, especially in multi-hop reasoning tasks that require iterative knowledge accumulation \cite{ref55}. The success of such a framework highlights the concept of memory, where the system's ability to store and reason over feedback is as crucial as the initial retrieval stage. \\
Beyond simple memory storage and retrieval, integrating feedback loops is essential for creating self-correcting and self-reinforcing RAG systems. This concept used in the field is prominently embodied in the Self-Reflective Retrieval-Augmented Generation (Self-RAG) framework, which introduces an adaptive retrieval mechanism controlled by self-reflection tokens \cite{ref109}. Rather than indiscriminately retrieving and using a fixed number of passages, Self-RAG learns to evaluate its internal state and the task requirements to choose when to retrieve and how to judge the relevance of retrieved passages and its own generation. The framework trains a single language model to generate special reflection tokens, which act as a self-feedback mechanism. These tokens enable the model to explicitly critique its own outputs, altering it for confidence, relevance, and factual correctness. This dynamic feedback loop ensures that retrieval and generation are mutually reinforcing, aligning with the robustness principles of adaptive noise budgets and instance-specific difficulty discussed earlier, allowing the model to adapt its behavior to diverse tasks, from open-domain QA to fact verification \cite{ref109}. \\
Reinforcement Learning (RL) offers another crucial mechanism for aligning the preferences of the retriever and the generator in a RAG pipeline. Many current RAG frameworks treat the retriever and the LLM as independent components, which leaves a behavioral gap: retrievers output results that are ``human-friendly'' while LLMs generate more ``LLM-friendly'' contexts. To mitigate this, a bridge mechanism has been proposed. This framework designs a ``bridge model'' by chaining supervised learning to construct an initial system and then applying reinforcement learning to optimize the bridge, maximizing the end-to-end generation performance \cite{ref80}. This is an exemplary instance of a self-improving system, because the retrieved context quality is evaluated downstream, and the retriever's preferences are continuously updated (via the bridge model) based on the LLM's feedback---a form of preference-based RL that dynamically adapts the system's behavior. The integration of RL is not solely about retrieval optimization, but also about understanding communication patterns. In the RAM framework, learning from feedback is executed via a communicative learning process, where the system learns from user feedback through communicated signals, in contrast to traditional real-time feedback into a static database \cite{ref55}. Thus, RL components ensures that the system can handle false-premise and multi-hop questions. \\
Beyond the adaptation of learned retrieval, iterative self-feedback mechanisms are designed to create robust systems. The Retrieval Augmented Iterative Self-Feedback (RA-ISF) framework proposes an iterative process that decomposes a task into sub-tasks and processes them through three self-supervised modules for understanding and reasoning, thereby validating and refining its own output \cite{ref5}. By applying multiple iterations of retrieval and self-correcting steps, the system prevents irrelevant text retrieval from degrading the model's performance. This iterative feedback loop ensures that augmentation is robust---the model is able to verify and re-align its understanding with the external texts, rather than solely depending on the initial retrieval attempt \cite{ref5}. It aligns with a novel concept of unsupervised refinement called ``Information Refiner,'' which posits that LLMs should be trained to refine retrieved texts of varied quality into cleaner, more complete answers \cite{ref52}. This perspective is implemented through an unsupervised training objective called InFO-RAG, which itself serves as a self-learning mechanism facilitating the model's ability to filter out irrelevant information and generate concise and accurate outputs, without external supervision, improving the overall robustness of RAG \cite{ref52}. \\
The synergy of memory and self-feedback is precisely that the system can overcome common RAG limitations. Using feedback to update its memory, the RAM system tests the existing information and updates it accordingly demonstrating significant improvements in handling false premises and multi-hop questions \cite{ref55}. Similarly, self-improving systems depend on a proper evaluation of utility. As LLMs are increasingly used as tools, there has been a focus on developing advanced utility evaluation methodologies that can accurately judge the potential of generated information to help provide answers, not just assess lexical relevance \cite{ref53}. This type of judgment is the baseline for the built-in feedback loop: if a self-adaptive RAG can internally judge the quality and usefulness of retrieved content, it can decide if the retrieval is insufficient or misleading and \emph{act on it} by generating or retrieving additional, more suitable content. This reinforces the iterative retrieval approach where the generated information is not treated as final. Direct applications of specialized memory for query expansion without any training costs are demonstrated in ARM-RAG \cite{ref108}. \\
However, the design of such synergistic systems involves critical trade-offs, echoing the ``noise-resistant training'' and ``counterfactual data augmentation'' from the robustness discussion. The integration of memory and feedback loops often introduces more demands on computational capabilities \cite{ref55}. For example, approaches like Self-RAG, with specialized reflection tokens, require large-scale and meticulous training on self-critique that can be costly and might not transfer to all tasks equally \cite{ref109}. The crucial challenge is balancing self-correction with internal parametric knowledge. However, by narrowing the interval between past and future interactions, the next generation of RAG systems incorporates the capability to retain and leverage feedback and self-improvement as a core principle. This paves the way for more autonomous, resilient, and context-aware RAG systems, facilitating dynamic QA, personalized assistants, and robust fact-checking. Moreover, as we transition to the following subsection on domain-specific fine-tuning, this adaptive principle ensures that the accumulation of feedback and the iterative refinement loops can be effectively specialized and transferred to industries such as legal, medical, and code generation, where precision in retrieval and continuous adjustment based on domain-specific data are paramount. \\
\#\# Domain-Specific Fine-Tuning and Multi-Modal/Semi-Structured Adaptations \\
Domain-specific fine-tuning and adaptation are critical for deploying Retrieval-Augmented Generation (RAG) systems in specialized fields where general-purpose models struggle with unique vocabularies, reasoning patterns, and data structures. While RAG's retrieval component can source relevant domain documents, the retriever and generator often require targeted optimization to fully exploit domain knowledge. Building on the adaptive principles of memory and self-feedback discussed previously, this subsection explores techniques for adapting RAG to domain-specific text (e.g., medical, legal, code) and semi-structured/multi-modal content, focusing on aligning retrievers and generators with specialized linguistic and reasoning patterns. These strategies ensure that the iterative refinement loops and preference-based learning mechanisms introduced earlier can be effectively specialized and transferred to high-stakes verticals. \\
A foundational challenge in domain adaptation is the scarcity of high-quality, labeled in-domain data. While fine-tuning remains a primary tool, its dependence on human-curated data can be impractical for new domains, languages, or tasks, especially when training data is nonexistent. Zero-shot transfer offers an alternative but often suffers an effectiveness cost, particularly for first-stage retrievers \cite{ref110}. Recent work has shown that simple domain adaptation methods can be effective even for these retrievers---a family of models often considered to generalize well. By leveraging domain-specific knowledge through pre-training-style objectives on target-domain corpora (without requiring annotated labels), these retrievers can be adapted to bridge cross-topic and cross-domain discrepancies. This approach---transposing a method designed for language adaptation to topic/domain adaptation---demonstrates that even traditionally ``simple'' models benefit from domain-aware pre-training rather than relying solely on task-specific fine-tuning, reinforcing the memory-induction principles from the previous subsection by making retrieval itself a continuous learning process \cite{ref110}. \\
The nature of the domain itself significantly impacts the adaptation strategy. For instance, in the patent domain, texts use a highly technical sublanguage that differs from general-purpose or even scientific language. A naive fine-tuning approach on such corpora may not focus on the truly domain-specific linguistic patterns. The Linguistically Informed Masking (LIM) method was designed to quantify differences across patent, scientific, and general-purpose language, showing that domain-adaptive pre-training using linguistically informed masking leads to systematically better representations compared to standard masking. Using LIM for the domain-adaptive pre-training of BERT and SciBERT consistently improved performance on patent-specific downstream tasks such as IPC classification and similarity matching \cite{ref111}. This highlights a broader insight for RAG in verticals like law, medicine, and code: retrieval components should be adapted using masking or modeling strategies that emphasize the domain-specific syntactic and semantic recurring patterns---much like the experiential memory from communication discussed earlier, where the system trains on repeated structures---rather than just factual knowledge. \\
Beyond masking strategy, the modular nature of RAG allows targeted optimization of its components. In a domain-specific RAG pipeline, encoders can be adapted to specific data structures such as expert texts, clinical guidelines, or peer-reviewed literature to improve semantic matching. However, when RAG is built on strong, pre-trained encoders, the question arises whether to fine-tune the encoder jointly with retrieval or through a decoupled post-hoc adaptation. Learning invariant representations with consistency and diversity in semi-supervised source hypothesis transfer has shown that when only a source-trained model is available, and a few labeled target examples---a setting akin to adapting a RAG retriever built on a general corpus to a specialized domain---maintaining prediction diversity and consistency regularization can enhance generalization ability \cite{ref112}. Furthermore, leveraging the low-confidence or ``hard'' examples during adaptation is critical; low-confidence samples should be used to encourage the model to explore the target's data structure, a crucial nuance for RAG where retrieval quality depends largely on the embedding model's fidelity to unseen medical or legal terms---directly mirroring the iterative self-verification mechanisms used in self-correcting retrieval loops for domain-specific queries \cite{ref113}. \\
The target domain in RAG rarely has cleanly separated label distributions for the generator or the retriever, but for semi-structured content (e.g., knowledge graphs, tables, and heterogeneous documents), the distinction between source and target becomes more pronounced, affecting adaptation at a structural level. The universal semi-supervised setting emphasizes the mitigation of common-class bias during pseudo-labeling and refinement, which directly impacts retrieval quality in RAG's feedback loops; when the system self-critiques its outputs, it depends on unbiased, context-aware retrieval \cite{ref114}. Similarly, the ``Pred\&Guide'' paradigm demonstrates the effectiveness of using a few labeled in-domain examples to refine sampling or class-wise weighting in an unlabeled target set, a direct analog to how domain-specific fine-tuning can re-calibrate dense vectors of a RAG generator to select relevant sources over generic ones \cite{ref115}. Such guided weighting improves retrieval precision and avoids ``common-class'' bias in the knowledge store, ensuring that domain-specific terms are not overshadowed by generic entities. \\
Multi-modal and semi-structured adaptation in RAG necessitates not only semantic alignment but also structural alignment across the source and target modalities. An extension of domain adaptation in a multi-modal context (e.g., referring expression grounding in the text-image setting) treats it as a knowledge transfer between textual and visual sources. Enriching inter-domain relations and transferring them across domains improves retrieval effectiveness, suggesting that RAG modules should learn cross-modal alignment functions to generate informative context across modalities and thus propagate multi-source signals in its external memory \cite{ref116}. Additionally, when RAG processes structured or semi-structured domains, the underlying adaptation should use those structural relations; multi-source domain adaptation can use domain disentanglement to generalize without relying on densely labeled target data \cite{ref117}. Given that heterogeneous RAG inputs (text, tables, graphs) correspond to multiple source domains, domain disentanglement is critical for selective augmentation, which integrates with the earlier feedback loop that favors instance-level customization. \\
For true domain alignment---especially in high-stakes domains like patient data---strong empirical performance remains a challenge. Strategies from domain adaptation literature show robustness against data imbalance and unknown label distributions; methods like ``Minimum Class Confusion'' incorporate a penalty that reduces entropy and can be used as a regularizer orthogonal to other adaptation methods, handling retrieval on relevant and irrelevant domain examples \cite{ref118}. Last, when source data are not available---as in privacy-preserving domains, which is common in decentralized RAG for finance or health---the RAG system can adopt a ``semi-supervised hypothesis transfer'' model, where the generator is fed with only the source hypothesis and a handful of in-domain pairs, allowing self-adaptive feedback loops to refine results without requiring full source corpus \cite{ref119}. Combining such transfer with consistency regularization improves robustness; contrastive and consistency regularization can be used to learn domain-invariant features even in the image segmentation setting, a paradigm for embedding formal-consistency of cross-domain entities within an RAG---reinforcing the need to align underlying latent representations from semi-structured data streams \cite{ref120}. \\
In summary, the overarching approach to domain-specific fine-tuning in RAG systems blends: (1) simple data-free domain adaptation to a target data corpus \cite{ref110}, (2) linguistically informed pre-training that captures the domain's unique syntactic and semantic patterns (echoing how the previous subsection trained on reiterative patterns) \cite{ref111}, (3) label-prediction guided class-weigh balancing to avoid common-class bias in retrieval \cite{ref114,ref121}, (4) contrastive disentanglement for multi-modal and semi-structured input \cite{ref122,ref117}, and (5) use of transfer from a source-only retriever/generator with no required access to source data \cite{ref119}. These strategies are especially important in legal, medical, and code domains where the cost of mis-retrieval is high, and where leveraging semi-structured, multimodal data is paramount to knowledge-intensive RAG systems. In this way, the adaptive memory and self-feedback capabilities from the previous section are systematically extended to specialized, high-vertical applications, preparing the ground for fully deployable RAG systems that continuously refine domain expertise through interaction and structured data. \\
\# Evaluation, Trustworthiness, and Risks in RAG Systems \\
\#\# Evaluation Frameworks and Metrics for RAG Systems \\
Evaluating Retrieval-Augmented Generation (RAG) systems is a multifaceted challenge that extends beyond traditional metrics for standalone large language models (LLMs). Because RAG systems integrate parametric knowledge stored in LLM weights with non-parametric knowledge retrieved from external corpora, their evaluation must measure not only the fluency and coherence of the final output, but also the relevance of retrieved evidence, the factual consistency of the generated text, and the correct attribution of claims to the provided sources. A robust evaluation framework is therefore essential for systematically measuring progress, diagnosing failure modes, and ensuring the trustworthiness of RAG-based applications. This subsection establishes the foundational evaluation paradigms, covering widely-adopted benchmarks, key quality metrics, and comprehensive evaluation frameworks that support multifaceted and multidimensional analysis. \\
The cornerstone of RAG evaluation is a suite of benchmarks designed to probe specific failure points, with hallucination---the generation of content that is unfaithful to the retrieved context or factually incorrect---being a primary concern. Among these, the HaluEval benchmark provides a large-scale collection of generated and human-annotated samples for evaluating LLMs' ability to identify hallucinated content \cite{ref123}. This benchmark, utilising a ChatGPT-based sampling-then-filtering framework, revealed that ChatGPT generates approximately 19.5\% hallucinated responses in open-ended tasks, and showed that existing models at the time struggled to consistently recognize these falsehoods. While HaluEval offers a controlled evaluation setting, the gap between such static NLP tasks and dynamic real-world user-LLM interactions is addressed by HaluEval-Wild, which collects challenging user queries from real-world interaction datasets like ShareGPT, categorizing query types to enable fine-grained analysis, and uses GPT-4 with retrieval-augmented generation to synthesize reference answers \cite{ref124}. Beyond user-centric interaction scenarios, robustness in multilingual settings constitutes another distinct concern. To this end, NoMIRACL provides a human-annotated benchmark for evaluating LLM robustness in RAG across 18 typologically diverse languages, introducing a non-relevant and a relevant subset to measure two capacities: the hallucination rate (tendency to hallucinate when the answer is not in passages) and the error rate (inability to recognize relevant passages) \cite{ref125}. These benchmarks differ in their focus but complement each other in revealing LLM failure modes across domains, languages, and interaction paradigms. \\
Measuring key evaluation criteria requires a suite of dedicated metrics, which again underscores the need for diverse specialized measures. Generation accuracy and factual consistency with the source form the majority of these metrics, with HaluEval providing binary annotations for identifying hallucination, accompanied by detailed explanations to distinguish factual from faithfulness hallucination \cite{ref123}. For dialogue-specific contexts, DiaHalu introduces a dialogue-level benchmark that covers four multi-turn dialogue domains and extends the factuality/faithfulness dichotomy into five subtypes \cite{ref126}, whereas for topic-focused dialogue summarization, TofuEval provides binary sentence-level hallucination annotations \cite{ref127}. In addition to these content-focused metrics, attribution quality deserves special attention. The question of how precisely each claim in a generated response is supported by its cited evidence is not automatically solved. In this area, AttributionBench stands out as a comprehensive benchmark compiled from various attribution datasets, with results showing that even fine-tuned GPT-3.5 models only achieve around 80\% macro-F1 on binary attribution formulations, highlighting the difficulty of automatic attribution evaluation \cite{ref128}. Combined, these metrics and benchmarks help assess correctness, but they often rely on gold standard references. \\
The production of gold-standard answers and human evaluation for high-quality assessment is a core component of the standardized evaluation process but can be costly and prone to error. To address these limitations, several frameworks focus on minimizing human annotation effort. For instance, the FEWL metric is specifically designed for contexts where human-written gold-standard answers are unavailable \cite{ref129}. FEWL uses answers from LLMs as a proxy for gold-standard answers and weights the expertise of these reference models, establishing sound theoretical grounding with empirically validated accuracy \cite{ref129}. Similarly, the general-purpose framework UFO introduces a unified evaluation pipeline that compares fact claims extracted from the generated text against four types of fact sources (human-written evidence, reference documents, search engine results, and LLM knowledge) \cite{ref130}. This plug-and-play approach allows for both flexible and scalable evaluation, with results suggesting that for most question-answering (QA) tasks, human-written evidence and reference documents are key and can sometimes substitute for one another in RAG settings. \\
Extending on top of the core evaluation metrics for correctness and factual consistency, we must consider how to ensure an evaluation framework captures the reliability and dimensional variety of RAG outputs. Several additional benchmarks and evaluation pipelines introduce advanced hallucination detection perspectives. Whereas HaluEval uses warm-up and filtering steps to generate hallucinated samples and mixes them with correct samples \cite{ref123}, the ChainPoll framework introduces an innovation in LLM hallucination detection \cite{ref131}, presenting the metrics Adherence---quantifying a model's capability to reason from retrieved documents---and Correctness---identifying logical fallacies---along with a refined set of benchmarks, RealHall. It is shown that ChainPoll outperforms state- of-the-art theoretical methods by 11\% and industry standards by 23\% in terms of AUROC \cite{ref131}. Although ChainPoll shows strong performance, there remains a gap in resource-constrained scenarios, underscoring the need for cost-effective detection methods. Holistic evaluation should not be limited to text-only generation. For vision-language models, benchmarks such as MMBench introduce evaluation paradigms like CircularEval that use ChatGPT to turn free-form model predictions into objective choices for robust evaluation \cite{ref132}, while VALOR-EVAL generalizes the CHAIR divergence metric to evaluate both faithfulness and coverage simultaneously, addressing the trade-off between hallucination control and informativeness that is often neglected when using metrics individually \cite{ref133}. To exploit pretrained models internally, INSIDE builds on this using dense semantic information and the EigenScore metric to evaluate self-consistency within the model's internal states, offering an efficient in contrast to traditional token-level uncertainty estimation for hallucination detection \cite{ref134}. \\
Given the importance of accuracy in the final generated text on the one hand and its groundedness and factuality on the other, evaluation must also reflect the stand-alone performance of the retrieval components of RAG. It is not sufficient to only measure the final generative output---evaluation should also examine retrieval quality, for instance precision/recall or end-to-end task success rates. With this integration, the interplay between coverage and faithfulness emerges; for example, retrieving more results can heighten coverage but may increase the likelihood of including contradictory sources, thereby affecting factuality. As with attribution, a comprehensive evaluation framework can attach penalties that are specific to whether the failure stems from the retriever or the generator. This dual-component assessment also underlines the importance of evaluating system capabilities in a broader context, as is done in benchmarks such as eRAG (evaluation required retrieval) or ClapNQ (curated, long-form, preserve-answer data), where contextual adaptation, temporal robustness, and seamless integration of evaluation mechanism are considered across different RAG pipelines. \\
In summary, the evaluation of RAG systems is increasingly realized through a conjunction of multi-scenario benchmark suites, a variety of quantitative metrics for correctness and hallucination, and unified flexible evaluation frameworks. The field is moving beyond isolated metrics toward holistic strategies that capture attribution, retrieval quality, faithfulness, and hallucination as interconnected pillars. Some evaluation methods stress provenance and evidence linking \cite{ref128}, while others improve on how external knowledge should be used for evaluation itself \cite{ref125,ref124}. While diverse in approach, they all converge on a common goal: to provide robust, reliable, and human-aligned assessment of augmented generation for real-world systems. Future efforts to build on these foundations must focus on realistically evaluating RAG with minimal human effort, standardizing good references when they are missing, and consolidating unsupervised, unsupervised, and contextual-awareness evaluation, so that RAG systems can be systematically compared and confidently deployed. \\
\#\# Theoretical and Empirical Analyses of Generation Risks \\
The theoretical and empirical understanding of generation risks in Retrieval-Augmented Generation (RAG) systems is critical for establishing trust and reliability. While RAG models are widely adopted to ground LLM outputs in external knowledge, rigorous frameworks for quantifying their residual risk and for characterizing the behavioral dynamics when parametric and non-parametric knowledge interact remain active areas of investigation \cite{ref135}. The evaluation strategies discussed above provide the measurement tools for diagnosing these risks; this subsection delves into the theoretical underpinnings and the empirical manifestations of generation risk---a necessary antecedent to the hallucination challenges examined next, essentially focusing on \emph{whether} and \emph{how} RAG systems fail before addressing the specific forms such failures take. \\
At the theoretical forefront, the C-RAG framework introduces a formal methodology to certify generation risks in RAG architectures. C-RAG addresses a fundamental question: can RAG reliably lead to lower generation risk compared to a single LLM, and if so, how can this risk be guaranteed? Specifically, C-RAG develops conformal risk analysis tailored for retrieval-augmented language models, deriving an upper confidence bound for generation risk---termed ``conformal generation risk.'' This provides provable guarantees under test distribution shifts for any bounded risk function. A key theoretical insight is the proof that RAG achieves a lower conformal generation risk than a single LLM, conditioned on the quality of the retrieval model and the transformer being non-trivial \cite{ref135}. \\
The technical foundations underlying such guarantees draw on broader advances in conformal prediction (CP), which offers a robust framework for uncertainty quantification that yields statistically effective prediction regions without strong assumptions on data exchangeability in its basic formulation \cite{ref136}. Since violations of exchangeability are common due to distribution shifts in real-world applications, subsequent work addresses this issue by exploring upper bounds on coverage differences across all confidence levels under conditional distribution shift, leveraging cumulative density functions of calibration and test scores. This becomes tractable when one assumes structural causal models with physical invariance \cite{ref137}. Similarly, extending conformal predictive systems to covariate shift settings through likelihood ratios demonstrates a persistent effort to establish non-parametric predictive distributions with calibration guarantees under non-IID conditions \cite{ref138}. Beyond these, researchers handle data efficiency and missingness challenges: cross-validation-based conformal risk control produces valid prediction sets when data are scarce \cite{ref139}, while marginal coverage guarantees can be upheld even when covariates contain missing values via imputation \cite{ref136}. \\
From a risk-analysis viewpoint, evaluating a RAG model's retrieval quality can itself be framed directly. The CONFLARE framework proposes a conformity-based method to link retrieval similarity scores to the correctness of the LLM's context. Its procedural setup---calibrating retrieval scores on questions answerable from the knowledge base, determining a similarity score cutoff, and retrieving all chunks above that threshold---provides control over the retrieval-sufficient rate at a user-specified error level \cite{ref47}. A failed retriever solution yields the result risk; CONFLARE thus implements robust uncertainty analysis at the retrieval stage as a crucial component in the overall chain of reasoning, linking formal conformal guarantees directly to the upstream retrieval component. \\
Empirically, the interaction between parametric knowledge and retrieved, non-parametric knowledge sets the generation risk in RAG systems. A key scenario is the knowledge conflict, where a document contradicts the parametric answer stored in the LLM, limits---this is where RAG's assumptions are tested. Recent work presents a framework for studying knowledge conflicts using realistic conflicting documents that update incorrect parametric answers, in contrast with prior work relying on synthetic contradictions \cite{ref9}. Under this realistic scenario, the rate of a model failing to update to correct knowledge in context turns out significantly lower than in synthetic cases. However, when failures do occur, the presence of an incorrectly parametric answer in the retrieved context reveals a ``parametric bias'': the model is \emph{more} likely to fail to update from the context if the wrong answer appears in it \cite{ref9}. This highlights a deep interaction between the model's fixed, parametric memory and retrieved evidence, confirming that internal beliefs can override contextual content. This provides a critical source for failure modes where the retrieved context is correct and yet the generation remains unanswered---a core component of empirical generation risk in RAG. \\
In the broader frame, the theoretical foundations from conformal and risk control are combined with empirical observations on parametric--non-parametric knowledge conflicts, highlighting the multifaceted nature of generation risk in RAG. This conceptual synthesis---that the intrinsic quality of retrieval is provably interwoven with the model's internal biases, and that these two objectives can be severely misaligned under realistic conflicts---lays the groundwork for designing systems that not only ground outputs but also remain well-calibrated and robust to component failures. This understanding is an essential precursor to the next section on hallucination mitigation; by appreciating when and why RAG systems fail, one can better design the detection and mitigation strategies that specifically address hallucinatory outputs which arise from these very risk dynamics. \\
\#\# Hallucination: Detection, Mitigation, and Trade-offs \\
Hallucination remains one of the most persistent and consequential challenges in Retrieval-Augmented Generation (RAG) systems, directly undermining the reliability and trustworthiness of generated content. While retrieval augments parametric knowledge with external evidence, the fusion of these two knowledge sources can introduce new failure modes, including conflicts between the model's internal priors and the retrieved passages, or over-reliance on noisy contexts. This challenge is intimately connected to the generation risk discussed in the previous subsection, as hallucinations represent a primary manifestation of the model's failure to correctly integrate parametric and non-parametric knowledge; indeed, the knowledge conflict scenarios that expose parametric bias directly parallel the conditions under which hallucination is most likely to occur. Building on this conceptual foundation, this subsection examines the landscape of hallucination detection and mitigation strategies within RAG, and the critical trade-offs that emerge, particularly between attribution and fluency. \\
\textbf{Detection of Hallucination} \\
Effective hallucination mitigation begins with reliable detection. A significant body of research focuses on identifying hallucinations at the token level during decoding, leveraging the model's own uncertainty. Since fluent and coherent generations can obscure factual inaccuracies, methods that analyze the model's internal state or output probabilities offer promise for early detection. One line of work proposes a token-level uncertainty quantification pipeline that fact-checks atomic claims within the output, introducing Claim Conditioned Probability (CCP), which isolates the uncertainty attributed to a specific claim value by removing the influence of surface form choices, showing improvements over standard uncertainty baselines for biographic generation \cite{ref140}. This is extended by research exploring sequence-level certainty for knowledge-grounded dialogue generation, demonstrating that sequence-level certainty---both probabilistic and semantic---correlates with lower hallucination rates, leading to a Certainty-based Response Ranking (CRR) decoding strategy \cite{ref141}. \\
Complementing these logit-based approaches, recent work probes the \emph{internal states} of LLMs since hidden layers retain dense semantic information that contributes to the final output distribution. The EigenScore metric and a test-time feature clipping method analyze LLM internal states to measure semantic consistency and diversity in the latent space, where higher consistency among responses suggests lower hallucination risk \cite{ref134}. More recently, an inner representation perspective reveals that correct generations exhibit significantly sharper context activations in the hidden states of in-context tokens compared with incorrect ones \cite{ref142}. In addition, a knowledge-constrained tree search decoding method is introduced to guide generation at the token level, coupled with a token-level hallucination detection mechanism that uses reward inflection approximation to convert sequence-level classifier scores into token-level signals \cite{ref143}. \\
Beyond the final token, heavier detection approaches often rely on analyzing the chain of reasoning, either via consistent cross-checking of context or by evaluating the effect of perturbations. For black-box LMs, a semantic-aware cross-check consistency detection method (SAC3) identifies question-level and model-level hallucinations by comparing responses across semantically equivalent perturbations of the same query, which self-check consistency methods alone cannot handle \cite{ref144}. Also, a method that focuses on the most informative and unreliable tokens, as well as token properties such as type and frequency, has been proposed to detect hallucinations without external knowledge or multiple samples \cite{ref145}. Alternatively, a stronger uncertainty-like approach evaluates how well the model can distinguish information from a potentially unreliable retrieval source, mimicking human attention to focus on the most valuable elements of the text---a critical marker of retrieval--input mismatch during RAG. The classic principle of \emph{self-consistency} is, in fact, weakened by model-specific biases and question-level ambiguities, making detection of induced wrong answers unreliable unless a cross-check is performed \cite{ref144}. \\
\textbf{Mitigation Strategies and Trade-offs} \\
Once hallucinations are detected, mitigation strategies aim to restore faithfulness. One effective and lightweight approach applies post-hoc citation generation---not merely penalizing hallucinations but also explaining the model's evidence. In particular, a new paradigm called Factual Entailment (FE) detects hallucinations by evaluating whether the generation is entailed by a specific portion of the retrieved document; combined with multi-task learning frameworks, this approach can localize the exact piece of text that contradicts reality \cite{ref146}. This finding can also be used to generate post-hoc retrieval, though one must note its interaction with top-k vs.~top-1 evidence. Furthermore, checking factual consistency via entailment probability between the input context and the model-generated beams using Natural Language Inference (NLI) shows that \textbf{beam search re-ranking} improves fluency while simultaneously reducing hallucinations \cite{ref147}, reinforcing the value of targeted decoding strategies. \\
In terms of model-intrinsic methods---those that do not rely on external sources---mitigation can proceed by reducing the model's over-reliance on retrieved tokens. An ``in-context sharpness'' decoding strategy reduces over-trust in hallucination-free intervals and applies penalties to the tuned in-context tokens based on hidden state entropy \cite{ref142}, a line of work closely tied to \emph{token-level and sequence-level uncertainty} \cite{ref140,ref141}. Another angle modifies model architecture through \textbf{elastic weight removal} of parameters such as those responsible for hallucinations, with extensions also targeting extractive hallucinations in dialogue grounding \cite{ref148}. In multimodal domains, hallucination is reduced through a joint ``over-trust penalty and retrospection'' decoding approach based on an attention matrix; the authors find that most hallucinations are linked to the model focusing too much on summary tokens while neglecting image tokens \cite{ref149}. Moreover, improved retrieval designs---including certainty-based tuning and adaptive retrieval strategies using semantic-aware modules that check cross-language consistency to activate retrieval \emph{only when needed}---demonstrate that selective retrieval can reduce both parametric and retrieval-driven hallucinations \cite{ref150}. \\
However, the central challenge in RAG design is that reliability comes with trade-offs. A particularly well-documented tension exists between \emph{attribution} and \emph{fluency}. Empirical analysis reports that naive top-k retrieval improves attribution relative to top-1 but can hurt fluency, while increasing model size improves both consistently. The authors also propose a simple heuristic---date-strengthened token selection---that retains the benefits of top-k planning without destabilizing the system \cite{ref151}. The same tension extends to the longstanding emphasis on \emph{fluency} at the expense of \emph{faithfulness}. Emerging ``selective generation'' methods aim to \emph{calibrate trust-aware decoding} using miscalibration in output probabilities, thus pushing toward spectrum-aware and uncertainty-calibrated systems that provide more nuanced, confidence-informed outputs but at some computational overhead. \\
In summary, hallucination mitigation in RAG is an interplay between advanced detection signals (token-level uncertainty, internal-state analysis, and entailment-aware cross-checking) and mitigation methods (post-hoc reasoning, re-ranking, and in-context decoding). Yet, as with the generation-risk guarantees and the security threats discussed above and below, the overarching challenge is navigating unavoidable trade-offs---among faithfulness, fluency, and computational efficiency---while remaining robust to adversarial perturbations. This points to an ongoing effort toward adaptive and well-calibrated robustness that eqepipes RAG systems for increasingly complex, real-world deployment. \\
\#\# Vulnerability to Adversarial Attacks and Security Risks \\
Retrieval-Augmented Generation (RAG) systems, while powerful, inherit significant security vulnerabilities from both their constituent large language models (LLMs) and their external knowledge components. The reliance on retrieved, often unstructured, data creates a broad attack surface that adversaries can exploit to manipulate outputs, exfiltrate sensitive information, or degrade model performance. These adversarial threats directly undermine the trustworthiness and reliability of RAG systems, intersecting with the hallucination challenges previously discussed---as poisoned or manipulated contexts can exacerbate the knowledge conflicts that trigger hallucination---and laying the groundwork for the broader safety and responsible generation concerns that follow, where security failures can cascade into misinformation, privacy breaches, and unethical outcomes. This subsection systematically examines the landscape of adversarial threats facing RAG systems, categorizing them by the attacker's knowledge and capabilities, and subsequently details both specific attack methodologies and emerging defense strategies. \\
\textbf{Threat Models and Attack Surfaces} \\
Adversarial threats to RAG systems can be broadly classified into two major categories based on the attacker's access, echoing the classic cyber-security division: white-box and black-box attacks. In a white-box setting, the attacker has full access to the model architecture, parameters, and gradients. This allows for sophisticated gradient-based optimization strategies. For instance, interpretable gradient-based attacks like AutoDAN can generate adversarial prompts by leveraging gradient information, producing readable but malicious prompts that can bypass perplexity-based filters---a common defense mechanism that flags fluent but nonsensical adversarial text \cite{ref152}. These crafted prompts can be used to jailbreak the LLM's safety alignment, leading to harmful outputs, which in the RAG context could co-opt the model into producing hallucinated or intentionally false content from otherwise benign retrieved knowledge. \\
However, a more practical and prevalent threat model is the black-box scenario, where the attacker has no knowledge of the model's internals and can only interact with it through queries. This paradigm has been extensively studied in computer vision \cite{ref153,ref154} and adapts readily to RAG. Black-box attacks in the NLP domain come in several forms. Query-inefficient attacks can be optimized using meta-learning frameworks, where a meta-generator learns from historical attacks to quickly adapt to a new target, thereby improving query efficiency \cite{ref155}. To further improve query efficiency, a spanning attack confines adversarial perturbations to a low-dimensional subspace derived from an auxiliary unlabeled dataset, making the attack more effective and harder to detect \cite{ref156}. Similarly, a margin-approaching zero-confidence attack framework like MARGINATTACK is designed to find the smallest perturbation causing a misclassification, reducing computational cost and improving accuracy \cite{ref157}. Local black-box attacks that constrain perturbations to the discriminative parts of an input, and also use model interpretability to guide the perturbation, have shown to significantly improve query efficiency \cite{ref158}. \\
\textbf{Prompt Injection and Multi-Turn Leakage} \\
A potent vector for black-box attacks is prompt injection, where malicious instructions are embedded in the input text to override the model's intended behavior. In RAG, this could be through malicious text injected into a retrieved source, making the LLM's distinguishable outputs particularly susceptible to adversarial manipulation that amplifies the overconfidence and misattribution noted in the hallucination discussion. A vital concern in RAG systems is prompt leakage, which can expose the RAG system's internal instruction structure, allowing the adversary to tailor future attacks. This risk is amplified in multi-turn interactions due to the LLM's sycophancy effect, which makes the model more compliant when a user presents themselves as vulnerable. A defined multi-turn threat model that shows significantly higher attack success rates (ASR) of 86.2\% in attacking a RAG system, leading to high leakage of this meta-knowledge \cite{ref159}. To counter prompt leakage and protect against these attack threats, defense strategies include the use of black-box query rewriting models and multi-tier defense mechanisms, where a combination of techniques is applied sequentially \cite{ref159}; these defenses inherently aim to render the retrieval and generation process more robust, similar to the retrieval reliability-enhancing techniques discussed in the previous context. \\
\textbf{Retrieval and Datastore Poisoning (Polluting)} \\
A principal attack vector in RAG is the corruption of the underlying non-parametric memory, a.k.a., the `\texttt{datastore\textquotesingle{}\textquotesingle{}.\ An\ attacker\ can\ manipulate\ the\ retrieved\ documents\ themselves\ to\ induce\ targeted\ misbehavior.\ This}poisoning' of the knowledge base can occur when the RAG system ingests content from previously untrusted or unverified public sources. By injecting a unique domain-specific token into a few malicious documents, the attacker can make the retriever fail and retrieve them with high confidence, even for queries that do not contain the token, and thereby controlling the external knowledge the LLM must incorporate into its output. The attack, demonstrated across various datasets and LLMs, highlights that if the poisoned documents are retrieved, the LLM's output is manipulated even if it maintains aligned behavior \cite{ref160}. The risk is exacerbated by the tendency of the LLM to be highly receptive to any factual context provided, even when that context conflicts with its parametric memory. This over-reliance is essentially a failure point that form the hallucination challenges in the previous subsection. A defense mechanism involves using a framework that aligns the LLM to handle inter-context conflicts and context-memory conflicts in a dialectical manner, and improves the defense against poisoned context attacks by 20\% \cite{ref160}. \\
\textbf{Low-Level Perturbation and Robustness to Corruption} \\
Low-level perturbation attacks can also compromise RAG systems. These attacks involve crafting deliberate, subtle perturbations to the retrieved documents or the input query that are imperceptible to a human but cause significant performance degradation in the generator, which we note directly contributes to the retrieval--input mismatch and hallucination risks of the previous subsection. This attack can bypass the retriever, often because the retriever's decision boundary is sensitive to numeric similarities. In common scenarios, the adversary can alter the target knowledge dataset, thereby manipulating the entire RAG output without passing explicit instructions. This shows that the retrieval module is not only a tool for the defense that retrieves accurate and contextual information, but also a key vulnerability, reinforcing the need to for RAG systems to calibrate their trust in the retrieval signals and to correlate that with the quality of evidence. \\
\textbf{End-to-End and Component-Level Defense Strategies} \\
The defense mechanisms against such adversarial threats span a multi-layered spectrum, closely aligning with the theme of the following subsection of providing credibility through reliability but extended is fully here due to their primary focus on safety rather than fluency. First, rewriting the user query before passing it to the RAG pipeline acts as a preemptive defense against attack and can be used to rewrite malicious content \cite{ref159}. Second, a post-retrieval hard sanitization step, such as a fact-checking process executed using NLI, can be used to remove malicious information before it reaches the LLM. Third, guardrails within the LLM's output, such as a classifier-based output filter, been proposed to block malicious or inconsistent outputs (e.g., attacks to the agent memory). The multi-tier combination of defenses still shows, within those bounds, an ASR of 5.3\%, indicating that current defenses under attack are not fully foolproof \cite{ref159}. Finally, robustness evaluation frameworks are being developed to simulate \emph{variants} of attacks and with which to test the defensive performance, showing that the defenses are robust to different attack vectors \cite{ref161}. By improving the ability to simulate variants through this method, the defense can be made more robust against generalizable attack patterns, making it less brittle and more effective for safe deployments. \\
In summary, security in RAG is a multi-faceted problem that spans the entire pipeline. Together with the hallucinations---is central to the foundational problem deeper risk that arise in the next subsection: trustworthiness, safety, and responsible generation. The the security attacks discussed here also include a need to \textbf{maintain fluency} which, given their core, ties to ensuring that attacks are detected without degrading actual user experience---and thus having a overarching insight into the trade-off landscape of future generation, and that the responsibility and safety of subsection 3.5 naturally follows. From this perspective, the distrust, privacy vulnerability, and possibility of new malicious input deserve the attention laid out next in this survey. \\
\#\# Trustworthiness, Safety, and Responsible Generation in RAG \\
Building upon the adversarial threats and robustness considerations discussed in the previous subsection, the deployment of Retrieval-Augmented Generation (RAG) systems in real-world applications introduces a broader set of challenges related to trustworthiness, safety, and responsible AI principles. While the prior section focused on defending against malicious external attacks, this subsection examines the equally critical task of ensuring that RAG systems are reliable, credible, and ethically sound in their everyday operation. Even in the absence of an active adversary, these systems are susceptible to failure modes that can produce harmful, misleading, or privacy-violating outputs. This subsection explores strategies for mitigating these risks, focusing on overconfidence calibration, misinformation mitigation, privacy preservation, and the development of mechanisms that ensure credible and responsible generation. \\
This discussion naturally connects to the evaluation frameworks presented in the following subsection, which provide the quantitative tools needed to assess these trustworthiness properties rigorously. \\
A fundamental challenge in RAG systems is their propensity for overconfidence, where the model produces predictions with high confidence that are nonetheless inaccurate. This issue is particularly acute when the retrieved context contains misleading, distracting, or contradictory information. Overconfidence directly undermines the utility of RAG for tasks like misinformation mitigation, where the system must provide not just an answer but a calibrated signal of its reliability. In this context, the model's confidence score becomes a crucial piece of information for the decision-maker, determining when to trust the model's output and when to escalate for human review. If miscalibrated---being overly confident when wrong, and uncertain when the model is right---the system becomes a poor tool for high-stakes decision-making. For applications such as fact-checking, medical advice, and legal guidance, the quality of the evidence retrieval pipeline alone does not guarantee safety. Rather, when the model's internal priors conflict with retrieved evidence, calibration ensures that we know whether to trust the source or the model. To address this, calibration methods are essential, aligning the model's internal confidence with its actual accuracy, especially in the context of retrieved evidence that may range from high-quality to misleading \cite{ref162}. \\
To build systems that are accurate and trustworthy, it is essential to measure and calibrate confidence for the specific task of mitigating misinformation. Existing approaches often struggle with a lack of reliable confidence measures. Combining confidence elicitation with sample-based consistency methods provides a more holistic approach to uncertainty quantification, improving calibration for misinformation mitigation tasks. For example, consistency among multiple LLM samples, combined with the model's self-proclaimed confidence, yields a more reliable measure of uncertainty than either method alone, enabling practitioners to know when to overrule the model \cite{ref162}. These uncertainty metrics also act as a trigger for human-in-the-loop inspection, complementing the evaluation metrics detailed in the following subsection, which measure the frequency of context hallucination and epistemic uncertainty at the claim level. \\
Beyond model confidence, the credibility of the generated content itself must be ensured. This entails implementing mechanisms for machine-generated text to be verifiable and traceable to reliable sources. In high-stakes contexts where decisions carry significant consequences, the ability to attribute knowledge correctly is a direct research area of active exploration \cite{ref163}. In environments where information changes rapidly, the RAG system's ability to correctly attribute knowledge is critical for building and maintaining trust. This requires not only identifying sources, but also ensuring the model's fluency does not overshadow its need for attribution, as humans may over-rely on a confidently generated, fluent narrative even when it is wrong, which directly ties again to overconfidence. A further factor that complicates trustworthiness is the presence of biases. Since models are trained on human-generated data, they are prone to inheriting and amplifying biases present in those data. RAG models can even exacerbate these biases if retrieved sources are skewed, as the system will preferentially retrieve and generate content that reflects tainted distributions. Developing systems that filter out irrelevant, opinionated, or biased information improves objectivity and is crucial for tasks involving intrinsically controversial topics, such as political or heath advice. \\
A significant area of research focuses on privacy preservation in RAG systems. Unlike traditional LLMs, RAG's retrieval component often needs to access, store, and process user's personal information to provide contextually relevant answers---such as in personalized conversational agents, medical assistants, or recommendation systems. The workflow processes user queries, retrieves relevant content from the datastore, and then generates a response, raising significant concerns about how these queries and retrieved data are stored and transmitted. Notably, when viewed through a game-theoretic lens, the act of data obfuscation by the user---in response to privacy concerns---triggers an interaction that can lead to degradation of model performance. This necessitates proactive measures from system designers, forcing a balance between privacy leak and utility, as the user becomes an active player in the system, not just a passive receiver of outputs \cite{ref164}. Furthermore, the entire RAG retrieval and generation pipeline must be privacy-preserving in itself, since privacy-preserving methods have their own limitations, and we therefore need to aim to integrate differential privacy into the output generation while guaranteeing a certain utility tradeoff \cite{ref165}. \\
This privacy concern is particularly acute for fact-checking and misinformation detection systems, where users may query about sensitive topics that reveal their beliefs, political leanings, or health information. For instance, a user might query a system about a medical condition, revealing personal health details. While the system may use the abstraction to protect self-disclosure, the process of generating a personalized response can inadvertently leak other sensitive information. RAG systems must therefore create an abstraction of a user's disclosure to a less specific term while still retaining the utility of that piece of information for the ongoing task \cite{ref166}. This creates a complex trade-off: without a clear boundary between data utility and its private information content, the abstraction must preserve the right signals for the model to generate contextually relevant outputs while simultaneously limiting leakage of sensitive attributes. \\
Just as the system must protect user privacy, it must also defend the integrity and quality of the information it relies on. The entire RAG pipeline is dependent on a high-quality datastore, and malicious manipulation can compromise its output. Thus, managing the datastore requires designing robust mechanisms to reduce the risk of misinformation propagation. One of the primary emerging threats is output that arises from a data store containing factually incorrect or maliciously placed context. Capabilities for detecting backdoor or irrelevant influences in the knowledge store or prompt itself are crucial. Instead of directly addressing poisoning attacks after retrieval, the system can implement a pre-retrieval filter that first assesses the content of the retrieved contexts and checks for misinformative or inconsistent pieces of information, blocking them before they reach the generation stage. This anti-poisoning lens connects back to the retrieval poisoning attacks described in the preceding subsection. \\
For RAG to be a reliable source of information, it must not only be accurate but also explicitly communicate its reliability to the user. Measuring the uncertainty of the system and detecting uncertain claims as a potential trigger for a human-in-the-loop is vital. Tests that leverage calibration of models are transferable, but the parameter updates, and even more so when RAG selects an output, the model across all its internal parameters should be well-calibrated to the context that is currently injected, even if distribution shifts have occurred \cite{ref167}. A human-in-the-loop could step in to filter or correct outputs that fall below a confidence threshold. By integrating such calibrated signal with the more detailed evaluations from the next section (e.g., claim-level uncertainty), the RAG system can be `self-aware' of its own limits---a prerequisite for trustworthy use. \\
Moreover, research must examine the issue of user interaction and the feedback mechanisms embedded in the system. Using users to provide feedback on the credibility and relevance of generated content creates a \texttt{knowledge\ feedback\ loop,} allowing the RAG system to adapt over time, learning to be more conservative or more exploratory based on the specific user's context. This point is also related to personalization: users with different backgrounds, cultures, or expertise levels will have different requirements and expectations from a RAG system. If developed in a privacy-preserving manner, personalized systems can improve trust by tailoring outputs to the user's level of context and understanding. \\
From a broader ethical standpoint, trust is not a binary property (true/false), but a multidimensional construct. To create RAG systems that are truly acceptable and trustworthy, we need a unified approach that combines calibrated model confidence, source provenance (credibility), and a controlled trade-off between privacy and utility that accounts for the strategic behavior of actors. Additionally, such systems should incorporate features such as selective generation and abstention--- where the model can say ``this falls outside my reach'' when it lacks relevant information---rather than forcing an overconfident or fabricated answer. In doing so, RAG system demonstrates a measure of honesty about the limits of its own knowledge, a key principle of responsible AI. \\
In summary, trust in RAG needs a multi-level, integrated strategy. Future work should focus off on hybrid designs that align uncertainty estimation, source verification, user feedback, and calibration into a unified framework, guided by the real-world demands for responsible and reliable AI systems. Such outcomes will not only produce better-functioning RAG models, but also broader, and systemic trust in LLMs. The metrics in the following subsection offer the necessary quantifying tools to measure and validate these dimensions. \\
\#\# Evaluation Metrics for Specific Quality Criteria and Edge Cases \\
The evaluation of Retrieval-Augmented Generation (RAG) systems demands metrics that extend beyond generic quality scores to capture nuanced failure modes and domain-specific challenges. This subsection examines targeted evaluation approaches addressing context hallucination, coverage errors, entity-level factual consistency, answer-conditioned uncertainty, resource efficiency considerations, and cognitive-contextual frameworks that attribute hallucination causes while controlling for confounding variables. \\
Building directly on the trustworthiness and safety concerns outlined in the previous subsection, evaluation must first provide the quantitative tools to measure whether a RAG system is genuinely reliable---not merely fluent. A central failure mode in this regard is context hallucination, where generated text contradicts or fabricates information relative to retrieved passages, a phenomenon that directly undermines the credibility and calibration principles discussed earlier. Specialized detection approaches are required to identify such hallucinations. One method leverages token-level uncertainty quantification to fact-check atomic claims within generated outputs \cite{ref140}. This approach recognizes that different components of an answer carry varying uncertainty levels, which can be harnessed to pinpoint hallucinations at the claim level. In this vein, the Claim Conditioned Probability (CCP) method isolates claim-specific uncertainty by removing the impact of surface-form choice and content-selection uncertainty, providing a more precise measurement of factual reliability \cite{ref140}. Conversely, sequence-level certainty metrics---both probabilistic (mean log-probability) and semantic (entailment-based agreement scores)---offer useful ensemble-level assessments of factuality. These sequence-level metrics are particularly well-suited for RAG systems where multiple retrieval contexts contribute to a single generation, and correlations between higher certainty and lower hallucination rates have been empirically validated \cite{ref141}. \\
Beyond detecting outright fabrication, evaluation must also capture more subtle trustworthiness failures, such as coverage errors. Coverage errors refer to omissions or incomplete statements of required information, often resulting from insufficient retrieval or inadequate synthesis, which can erode the system's reliability for high-stakes decision-making. To this end, entity-level factual consistency metrics provide a complementary lens, focusing on the accuracy of named entities and their relations within the generated content. This approach ties directly into research on entity-based knowledge conflicts, which formalizes the problematic scenario of contextual information contradicting learned parametric knowledge \cite{ref168}. Entity-specific evaluation of this kind can detect over-reliance on memorized parametric knowledge versus contextual retrieval---an important indicator of whether the RAG framework is genuinely augmenting the model or merely defaulting to its priors \cite{ref168}. Such metrics also support measuring the degree to which RAG frameworks mitigate hallucination and adapt to evolving time-dependent knowledge \cite{ref168}. Collectively, these findings highlight that entity-level metrics are crucial for evaluating RAG systems, since entity-level errors often expose fundamental weaknesses in source selection or attribution. \\
To provide a more complete diagnostic view, evaluation can also be conditioned on the internal conditions under which an answer is generated---an approach known as answer-conditioned uncertainty quantification. This method captures both epistemic and aleatoric uncertainties---where epistemic (knowledge-related) uncertainty is more indicative of hallucination than aleatoric (inherent randomness) \cite{ref169}. By separating these two sources of uncertainty, answer-conditioned metrics serve as a diagnostic tool for understanding whether hallucination stems from flawed retrieval, in line with the adversarial and datastore concerns mentioned in prior subsections, or from noisy generation. For a robust evaluation, answer-conditioned uncertainty should be complemented by internal representation analysis---such as measuring the semantic consistency and sharpness of in-context token activations \cite{ref142} and truncated extreme activations on the internal states at test time \cite{ref134}. These internal assessments shed light on how the model processes retrieved information and provide a richer, representation-grounded perspective than surface-level evaluation alone. \\
Beyond quality-based metrics, the practical feasibility of RAG systems also hinges on their resource footprint---another critical dimension that directly impacts real-world trust. However, it is worth noting that the mechanisms for improving efficiency---such as dynamic caching of intermediate states or pipeline parallelism---are not directly substantiated by the hallucination-detection papers cited in this survey given their evaluative focus. Unless such efficiency improvements are explicitly measured in a given study, they should be treated as an area for future research. Nonetheless, resource-aware metrics---accounting for computational costs like token consumption, latency, and GPU memory, as well as energy usage---should be measured alongside quality metrics, and sometimes traded off with them, in order to offer a fuller picture of practical capacity for deployment at scale. \\
A further layer of diagnostic depth comes from cognitive-contextual frameworks that aim to attribute hallucinations to specific model capabilities rather than surface-level errors. A system's simple hallucination ratio can mask crucial differences between in-context and parametric sources of hallucination. Advanced evaluation frameworks should account for confounders and attribute hallucination tendencies to specific risk factors such as commonsense memorization, relational capability, and instructional consistency \cite{ref170}. By using the effects of confounders, association analyses can clarify the underlying mechanisms and risk factors that drive hallucinations \cite{ref170}. Such frameworks allow evaluators to identify aspects of the model driving systemic errors, moving beyond aggregate scores to targeted improvements---a perspective that supports both robust diagnosis and transparent safety practices. \\
Finally, standard evaluation benchmarks are necessary but not always sufficient, particularly in edge-case and specialized scenarios where hallucination is most harmful. Resources like HypoTermQA demonstrate a use case for synthetic datasets of hypothetical claims, effectively pushing LLMs into unfamiliar or counterfactual territory to expose their propensity to hallucinate \cite{ref171}. This approach offers an automated and scalable proxy for testing RAG in high-stakes domains such as law, health, or finance. A complete evaluation strategy should, therefore, complement quantitative benchmarks with qualitative and worst-case assessment frameworks, along with user-centric assessments that measure actual perception of reliability \cite{ref172}. Such user-focused tools frame RAG evaluation as a means to foster trustworthy human-machine interaction, ensuring the metrics themselves are aligned with how real users perceive and rely upon the system. \\
In summary, robust RAG evaluation metrics go beyond standard quality measures, integrating targeted specifications such as context hallucination, coverage errors, entity-level deficits, answer-conditioned uncertainties, resource efficiency, and hallucination attribution towards causes. By adopting a multifaceted and cognitively aware approach, researchers can both compare RAG systems more effectively and diagnose them more precisely, thereby contributing to the responsible deployment demanded by the trustworthiness principles previously discussed. These evaluation frameworks thus serve as the necessary empirical counterpart to the safety and responsibility concerns raised earlier, ensuring the system's performance is not only measurable but also aligned with the higher principles of responsible AI, as they evolve in the concluding sections of this work. \\
\# Applications, Benchmarks, and Domain-Specific Adaptations \\
\#\# Domain-Specific RAG Systems in Medicine and Healthcare \\
The application of Retrieval-Augmented Generation (RAG) in the medical and healthcare domain represents one of the most active and consequential areas of research, driven by the critical need for factual accuracy, up-to-date knowledge, and the reduction of harmful hallucinations in clinical settings. Unlike general-purpose applications, medical RAG systems must contend with specialized vocabularies, complex reasoning over biomedical entities, high-stakes decision-making, and strict requirements for reliability and safety. This subsection examines key advances in this area, focusing on frameworks for clinical question answering, medication safety, and biomedical reasoning, alongside critical evaluations of their performance and the strategies employed to ground models in verifiable knowledge. \\
A fundamental challenge in medical RAG is bridging the gap between raw textual information and structured, relational knowledge. A widely adopted approach is the integration of knowledge graphs (KGs) to provide a structured, factual backbone for generation. For instance, the KG-RAG framework leverages the massive biomedical knowledge graph SPOKE to generate meaningful text that is rooted in established knowledge. This task-agnostic framework enhances the performance of large language models (LLMs) such as Llama-2, GPT-3.5, and GPT-4 across various prompt types, from simple true/false queries to complex drug repurposing, with a notable 71\% performance boost for the smaller Llama-2 model on a challenging multiple-choice dataset \cite{ref173}. This demonstrates that explicit, structured knowledge can help compensate for the parameter limitations of open-source models. Similarly grounded in knowledge-graph principles but addressing multimodal data, the REALM framework applies RAG to multimodal Electronic Health Records (EHR) for improved clinical predictive tasks. It uses an LLM to extract task-relevant medical entities and match them with a professionally labeled external knowledge graph (PrimeKG), pulling in corresponding medical knowledge. By aligning with clinical standards, this framework aims to eliminate hallucinations and consistency issues in a multimodal setting, integrating the extracted knowledge with time-series and clinical note data through a custom adaptive fusion network \cite{ref174}. Across both frameworks, the success of these approaches highlights that the integration of curated knowledge graphs is crucial for imposing machine-readable logic and constraints on raw text, reinforcing answer reliability. \\
Beyond static knowledge graphs, significant progress has been made in retrieving richer and more relevant textual evidence at scale. The MIRAGE benchmark, which introduces a first-of-its-kind evaluation framework including 7,663 questions from five medical QA datasets, provides a rigorous methodology to validate such systems. Through the associated MedRAG toolkit, large-scale experiments were run on 40 combinations of corpora, retrievers, and LLMs. This influential framework demonstrated that a robust RAG system can improve the accuracy of six different LLMs by up to 18\% over chain-of-thought prompting, potentially elevating the performance of models like GPT-3.5 and Mixtral to GPT-4-level with the right combinations of medical corpora and retrievers \cite{ref175}. The ``lost-in-the-middling'' effect and a log-linear scaling property between retrieval token count and system predictability were also identified, providing critical, actionable insights for implementation. \\
To further enhance the reasoning capacity of these systems, frameworks have been developed that introduce mechanisms of self-reflection and domain-specific instruction. Traditional RAG approaches often treat LLMs as passive consumers of retrieved information. The Self-BioRAG framework addresses this limitation by introducing a component for domain-specific retrieval. Through 84k filtered biomedical instruction sets, it trains the model not only to retrieve and generate but also to issue specific reflective tokens and self-ensure the quality of its explanations and generated responses. This makes the reasoning process more explicit and transparent, improving generalization across domain-specific problems \cite{ref176}. Correspondingly, the ActiveRAG framework focuses on transitioning from passive knowledge acquisition to active learning mechanisms. It utilizes a ``Knowledge Construction'' mechanism to associate newly retrieved external knowledge with previously learned generative memory, and a ``Cognitive Nexus'' mechanism to calibrate the model's intrinsic cognition with contexts generalizing for chains of thought and reasoning process of reasoning \cite{ref177}. \\
However, despite the promise of RAG, the retrieval process itself introduces inevitable uncertainty. When retrieval fails to identify the necessary information, there is no guarantee of response validity. The CONFLARE (CONFormal LArge language PK) framework addresses this challenge by introducing techniques from conformal prediction. Through statistical analysis on a development set, CONFLARE sets a confidence-based threshold of the selected retrieval pool, ensuring that the true answer is captured in the context with a user-specified high-probability guarantee (IQR) -- potentially critical for targeting a pre-specified error rate \cite{ref47}. In practice, this translates into continuously checking if the retrieved indices align to the clinical query guarantee that if the actual information resides in the corpus, then the chance that it is missed by the retrieval satisfies defined control levels \cite{ref47}. Thus, it explicitly ties statistical calibration to retrieval recall, ensuring that none of the critical stages silently fails. \\
In clinical practice, RAG has shown significant potential in decision support systems and safety-critical applications. A study investigating a novel RAG framework as a Clinical Decision Support System (CDSS) for unsafe medication prescriptions showed that a ``co-pilot'' mode in tandem with junior pharmacists significantly improved the detection of severe drug-related problems (DRPs), while optimizing accuracy, recall, and F1 scores, especially in identifying life-threatening conditions \cite{ref178}. A pre-operative medicine case study further illustrated the direct clinical utility: a RAG-enhanced LLM automatically generated responses for surgical procedures, achieving performance non-inferior human-generated instructions (91.4\% for GPT4.0 + RAG vs 86.3\% for junior doctors) with substantially faster turnaround times \cite{ref179}. \\
Efficient and well-scoped retrieval supports play a critical role in achieving high overall performance. The release of open-source corpora, such as the MedWiki dataset within the clinical error detection task, provides a rich resource for improving retriever methods \cite{ref180}. For specialized, zero-shot retrieval, medical knowledge graphs containing rich semantics can be leveraged to train deep models without any annotation. Such architectures significantly outperform common zero-shot baselines, such as BM25 and Clinical BERT, by 7\% to 30\% higher recall across major medical ontologies including UMLS, SNOMED, and ICD-10 \cite{ref181} *. These methodologies ensure that RAG systems can reliably access both broad and specialized medical knowledge---critical when working with both high and low-resourced domains. In addition, the multi-view framework called MVRAG has been proposed for knowledge-dense domains, where specialized environmental constraints generated through query reformulation into various different domain vantages have shown significant gains in recall precision using unique contexts \cite{ref182}. \\
Finally, integrating retrieval and generation directly during training is crucial for optimizing the full RAG pipeline. The JMLR (Jointly Medical LLM and Retrieval) framework trains the retriever and generator jointly. This synchronized training not only aligns the retriever better with the clinical reasoning tasks, but also reduces overall computational cost. The resulting model surpasses larger benchmarks like Meditron-70B and Llama2 with RAG on medical QA tasks, all with significantly lower training time budget \cite{ref183}. \\
In summary, the landscape of medical RAG is evolving toward more sophisticated, dynamic architectures that can reason over different types of medical knowledge---clinical, multimodal, knowledge-graph-based, and textual evidence---while actively reflecting on retrieval uncertainty and output accuracy. Future advances will likely focus on building enhanced robustness for real-world clinical deployment, large-scale integration with health informatics systems, and novel approaches to the overarching challenges of privacy, interpretability, and evidence traceability in personalized medicine and streamlined clinical trial applications. \\
\end{longtable}
}

\subsection{RAG in Legal, Multilingual, and Enterprise Knowledge-Intensive Domains}

Building upon the structured knowledge integration and uncertainty-aware mechanisms developed in medical RAG, the deployment of Retrieval-Augmented Generation in legal, multilingual, and enterprise environments reveals both the generalization of these architectural principles and the emergence of domain-specific challenges. These domains share with medicine a demand for high precision and factual grounding, but introduce additional complications: fragmented language coverage, multi-perspective reasoning over dense corpora, cultural adaptation, and the integration of heterogeneous knowledge sources into production-grade systems. The progression of RAG research in these settings demonstrates how the core paradigms of structured retrieval, self-reflection, and statistical calibration---introduced in clinical contexts---must be adapted to operate across linguistic and jurisdictional boundaries.

\subsubsection{Legal Domain Implementations: Precision, Multi-View Retrieval, and Low-Resource Challenges}

The legal domain presents unique challenges for RAG, paralleling the knowledge-graph grounding strategies discussed in medical contexts. Legal reasoning requires not only the retrieval of relevant statutes, case law, and legal articles but also an understanding of distinct domain perspectives---a limitation of conventional retrieval methods that lack multi-perspective views, which are essential for improving interpretability and reliability \cite{ref182}. To address this, the \textbf{multi-view retrieval framework (MVRAG)} was introduced, utilizing intention-aware query rewriting from multiple domain viewpoints to enhance retrieval precision in legal and medical case retrieval \cite{ref182}. Rather than simply rewriting queries into different semantic forms, MVRAG generates query variants that express specific domain knowledge perspectives, unleashing the potential of multi-view information to accelerate the application of LLMs in legal knowledge-intensive fields, demonstrating the same multi-view strategies that proved successful in medical retrieval \cite{ref182}.

Low-resource languages represent a significant barrier in law, echoing the resource-scarcity challenges encountered in specialized medical subdomains. The solution to the Vietnamese Automated Legal Question Answering Competition 2023 (ALQAC 2023) combines similarity ranking with deep learning models for legal document retrieval, and employs a range of adaptive techniques to handle different question types for answer extraction \cite{ref184}. This work achieved outstanding results, demonstrating the potential of RAG-based systems in legal QA for lower-resource languages, while also highlighting the value of data enrichment strategies \cite{ref184}. To support robust development in this area, large-scale curated resources are being assembled. For instance, a 689GB corpus spanning languages and jurisdictions has been made available for pretraining legal NLP models \cite{ref185}. The multilingual models trained on this corpus achieved state-of-the-art performance on LEXTREME, while English models surpassed existing benchmarks on LexGLUE, paving the way for stronger baseline retrieval and generation models in the multilingual legal domain \cite{ref185}. Another dataset further contributes to the legal NLP landscape, containing 65k EU laws officially translated into 23 languages and annotated with the EUROVOC taxonomy \cite{ref186}. This work critically highlights the effect of temporal concept drift, showing that chronological splits are essential for realistic evaluation \cite{ref186}. More significantly, the experiments reveal both that fine-tuning a multilingually pretrained model in a single source language can lead to challenges in preserving multilingual knowledge, overall harming zero-shot cross-lingual transfer, and that strategies like partial fine-tuning and adapters can help retain transferability---lessons that directly parallel the calibration techniques deployed for retrieval uncertainty in clinical settings \cite{ref186}.

\textbf{Multilingual RAG: Persistent Dual Challenges of Low-Resource Languages and Knowledge Conflicts}

The broad challenges of cross-lingual knowledge grounding found in legal settings extend naturally to multilingual RAG at scale. A significant hurdle is the evaluation of LLMs' ability to abstain or refuse when retrieved contexts do not contain correct answers in a typologically broad set of languages. The NoMIRACL effort addresses this head-on, providing a human-annotated dataset that includes both ``non-relevant'' subsets (where the passage contains no correct answer) and ``relevant'' subsets across 18 typologically diverse languages \cite{ref125}. Evaluation of multilingual-focused LLMs reveals a stark trade-off: most models struggle with balancing the two necessary capacities. For example, LLAMA-2, Orca-2, and FLAN-T5, achieve hallucination rates above 88\% on the non-relevant subset (hallucinating the answer when it is not in the passages), while Mistral overall reduces hallucinations but still maintains an error rate of up to 74.9\% on relevant subsets \cite{ref125}. The best model---while demonstrating some compromise---lacks the precision for safety-critical deployment, reinforcing the need for explicit confidence thresholding mechanisms akin to the conformal prediction strategies used in the clinical domain \cite{ref125}.

Pushing the frontier forward, the Prompts Augmented by Retrieval Cross-lingually (PARC) pipeline improves zero-shot performance on low-resource languages (LRLs) by augmenting the context with semantically similar sentences retrieved from a high-resource language (HRL) as either labeled or unlabeled prompts \cite{ref187}. Their experiments verify a consistent improvement in zero-shot performance --- reaching +5.1\% in the unlabeled and +16.3\% in the labeled setup across three downstream tasks across 10 LRLs \cite{ref187}. They additionally found a significant positive correlation between transfer performance and the degree of language relatedness, underlining the importance of language closeness in RAG \cite{ref187}. Work on cross-lingual question answering over knowledge bases (xKBQA) has also gained insights from cross-lingual understanding. One approach frames xKBQA from a reading comprehension perspective, converting KB subgraphs into a passage format to reduce the gap between structured KB schemas and natural language questions, effectively applying the same knowledge-graph-to-text bridging principle that proved effective in biomedical retrieval \cite{ref188}. This method leverages multilingual pretrained language models and existing high-quality cross-lingual machine reading data to fine-tune and strongly mitigate data scarcity \cite{ref188}.

However, robust multilingual RAG systems must also grapple with cultural adaptation and induced knowledge conflicts. Cross-lingual transfer for Named Entity Recognition depends largely on the overlap of entity chunks between languages, and despite positive transfer, subtle robustness issues arise under input perturbation \cite{ref189}. The low-resource and non-Latin script languages have been shown to be a persistent benchmark bottleneck: when evaluating LLMs on multilingual, multimodal (text and image), and multilevel exams, top performers such as GPT-4 still see significant performance dips for these languages \cite{ref190}. This is important assessment, because the fact that current models, even strong ones, struggle with low-resource scripts directly translates to RAG, as the quality of low-resource retrieval---and therefore the reliability of the knowledge grounding---is correspondingly lower.

\textbf{From Research to Deployment: Enterprise, Knowledge-Intensive, and Multi-Source Systems}

Operating at the intersection of the technical sophistication found in medical and legal applications and broader multilingual concerns, enterprise deployment of RAG systems demands both functional integration and an ability to handle long-tail, personalized, and multilingual data. A consistent design philosophy---from domain-specific architectures to general enterprise assistants---has been to build multi-source, multi-view RAG pipelines. For instance, the MVRAG framework supports a multi-view abstraction not just from the retrieval side, but as a higher-level, composite method with broadly applicable to legal and medical use cases, demonstrating commonalities across knowledge-dense verticals \cite{ref182}. Analogously, the ALQAC solution learns across a multi-step architecture that includes similarity ranking and deep learning for retrieval and a sequence of adaptive answer extraction techniques, simulating how a production legal assistant might function \cite{ref184}. The final generation step for a mix of question types in that paper was significantly boosted by iterative, across-level QA strategy.

A crucial enterprise capability---cross-lingual entity-centric retrieval and personalization---has also seen advances. Training retrieval models on artificially code-switched data yields large gains simultaneously in multilingual and cross-lingual scenarios \cite{ref191}. The ``MultiEURLEX'' and ``MULTI3NLU++'' datasets further stress the demand on modeling for realistic enterprise use-cases: unified systems must handle multiple intents simultaneously within a single utterance across multiple languages \cite{ref186,ref192}. The MIA 2022 Shared Task similarly handles open-retrieval QA across 16 diverse languages, indicating continuous development of multilingual benchmarks that push retrieval and generation systems closer to handling realistic, multi-language scenarios \cite{ref193}.

In summary, RAG systems in legal, multilingual, and enterprise environments are capable of delivering targeted answers with precision when the associated corpora are well-curated and aligned with the retrieval and generation components. However, the dominant dangers---sparse knowledge in low-resource languages, high hallucination rates when retrieval fails, linguistic and cultural interference, and the enormous complexity of deploying these models in time-critical, privacy-sensitive production settings---remain, and their remediation requires the same deep coupling between robust retrieval, statistical calibration, and multi-view reasoning that proved essential in both the clinical and multimodal domains\cite{ref125,ref190}. The continuity across these settings is architecture parallel and methodological---the next generation of RAG must not only become technically better at integrating multi-view retrieval \cite{ref182} but must also become robust in the face of missing knowledge and ambiguous queries, consolidating knowledge across diverse data modalities---a requirement that remains equally pressing in the technical, scientific, and multimodal domains that follow, where retrieval scaffolding becomes the crucial interface between the user's intent and heterogeneous data sources. The same retrieval-robustness principles---abstention when knowledge is absent, confidence calibration, and multimodal grounding---that enable clinical safety and multilingual reliability const to shape the design of technical and scientific RAG pipelines, ensuring that the integration of heterogeneous knowledge remains a source of accountability rather than overconfidence in high-stakes settings.

\subsection{Multimodal, Video, and Speech RAG Systems}

The extension of Retrieval-Augmented Generation (RAG) into multimodal domains represents a significant frontier in artificial intelligence, moving beyond purely text-based knowledge retrieval to incorporate visual, auditory, and spatiotemporal information. This shift is driven by the recognition that real-world data is inherently multimodal, and systems that can seamlessly integrate and retrieve across these modalities are better positioned to tackle complex, open-ended tasks. In video understanding, speech processing, and cross-modal retrieval, RAG architectures are being adapted to overcome the limitations of purely parametric models, providing a dynamic interface to external, non-parametric knowledge stores. The convergence of these multimodal capabilities with the technical and scientific domains discussed in the following subsection highlights a broader architectural evolution: retrieval is no longer just about text passages but about structured, heterogeneous knowledge that can be queried on demand.

A core challenge in applying RAG to multimodal data, particularly video, lies in the inherent difficulty of representing heterogeneous content in a unified and searchable format. Traditional RAG systems often rely on pre-processing data into text descriptions, which can result in significant information loss and computational overhead. For instance, converting a long video into a single textual summary sacrifices the fine-grained visual details, conversational dynamics, and audio cues that may be crucial for answering specific user queries. To address these limitations, an incremental workflow for RAG, iRAG, has been proposed \cite{ref194}. Instead of upfront full-video transcription, iRAG quickly indexes large repositories of multimodal data, using an incremental workflow to opportunistically extract more details from select portions of the video when an interactive user query is presented. This on-demand, query-driven fine-grained extraction not only resolves the information loss issue but also avoids the lengthy conversion times of traditional systems, achieving 23x to 25x faster video-to-text ingestion while maintaining high-quality responses. This highlights a key trend in multimodal RAG: moving away from a ``transcribe everything first'' approach and towards a more adaptive, query-aware retrieval of relevant knowledge---mirroring the adaptive and hierarchical retrieval mechanisms observed in high-stakes technical and scientific domains.

The construction of large-scale, diverse multimodal datasets is another cornerstone of this domain. To train foundational models capable of processing video, audio, and subtitle tracks in conjunction with text, creating high-quality aligned data is paramount. VAST-27M, is a large-scale omni-modality video corpus that addresses this challenge \cite{ref195}. The corpus was created by automatically collecting 27 million video clips, using a vision captioner and an audio captioner to generate corresponding captions, and then employing a Large Language Model (LLM) to integrate previously generated captions with subtitles into omni-modality captions. This approach uses the LLM to structure raw, multi-source data into cleaner and richer text that serves as high-quality training examples. Building on such a corpus, the VAST architecture itself is another key example; it has been designed to perceive and process vision, audio, and subtitle modalities to support various tasks including video-text retrieval, captioning, and QA, demonstrating the effectiveness of training models with curated and retrieved cross-modal data.

Interaction between modalities is not merely a in retrieval problem; it also involves intentional multimodal representation learning. Inter-modality alignment and contrastive learning are primary techniques for mapping different data types into a shared semantic space. The CLIP4VLA architecture, for instance, extends the vision-language contrastive paradigm of CLIP to include audio by applying contrastive learning to explore the correlation between audio and other modalities. Further, CLIP4VLA introduces a type token to dynamically learn audio information (verbal vs.~nonverbal), demonstrating the need to model diverse sub-information when processing a single modality \cite{ref196}. This approach helps establish a viable cross-modal joint embedding space from pre-trained backbones, a critical requirement for any retrieval system that attempts to locate and retrieve relevant content. In this same spirit, the model CLIP4CMR leverages vision-language pre-trained models as a backbone while adding a novel learning objective to computationally efficiently perform supervised cross-modal retrieval, thereby using and also studying their robustness to language availability \cite{ref197}.

Achieving retrieval across modalities brings unique complexity to representation learning---and understanding how these representation to interact with each other is not trivial. A breakdown on transferability can be dissected through ``Video-Teller'', which is built on the image-based BLIP-2 model, incorporating a cascaded Q-former with a fine-grained alignment objective to enhance video-to-text generation \cite{ref198}. This cascaded structure helps in fusing parallel visual and auditory information with text, while the fine-grained modality alignment objective shows that with modest training costs, the generated language representations can be enhanced by capitalizing on pre-trained vision and language modules. This cascaded approach is conceptually similar to the hierarchical knowledge graphs and skill graphs used in engineering and robotics, where multi-scale, structured organization enables efficient retrieval for downstream generation. Similarly, work on universal multimodal retrieval, such as MuMUR, demonstrates how knowledge can be transferred from a multilingual model to boost performance when no explicit labeled video data is provided, extending unified retrieval to both image and video data with multilingual text queries \cite{ref199}. This shows the impact of aligning embedding spaces across languages and modalities, potentially using state-of-the-art machine translation to generate pseudo-ground-truth multilingual visual-text pairs to improve model adaptability.

Yet, integrating text with non-linguistic inputs has its own barriers, especially at the intersection of fine-grained visual object detection and missing object representation. In addressing visual hallucination in multimodal LLMs, the Lyrics approach solves it by injecting local visual features extracted by a visual refiner (for image tagging, detection, and semantic segmentation) into the Querying Transformer \cite{ref200}. This system also integrates boxes and tags into the input text, providing access to concrete object location not well captured by global features. Such an architecture pushes multimodal RAG beyond retrieval-passes that rely on coarse captions, enabling more precise visual grounding and reasoning. Another emergent line of research---coherent paragraph captioning from untrimmed video---has been built via ``VLCap'' \cite{ref201}. VLCap leverages human-like perception of an activity in untrimmed videos, by breaking into both ``vision'' to perceive global scene semantics and ``language'' to describe sub-actions of visual and non-visual elements. Training with a Vision-Language contrastive loss keeps global-local alignment, leading to qualitatively accurate and text-coherent paragraphs.

To leverage the strong foundational capabilities of image LLMs rather than building from scratch, a resource-efficient pipeline anchors from Image to Video LLMs provides a useful blueprint---by inserting a ``temporal-token'' to facilitate model compression. In this ``RED-VILLM'' pipeline \cite{ref202}, the model inserts a temporal-query module into the query of an Image LLM, enabling the extended video understanding abilities with minimal training data. It highlights that limited additional data and modules can be used in architecture scaling, and such adaptation pathways often align with routing-level augmentation for retrieving relevant video events.

However, the integration of multimodal information for common grounding remains an open question. One study probing a pretrained multimodal video transformer against functional activity from participants viewing a popular show found that those models do not completely align embeddings to brain-derived representation---showing new complexities in internal semantic hierarchies are still missing \cite{ref203}. These qualitative results, while showing design-driven patterns in certain setups, additionally caution against relying solely on model representations for comprehensive reasoning across modalities.

Collectively, this research trail shifts the RAG paradigm from the text-based retrieval of explicit textual knowledge to the combining of latent knowledge across multiple sensory channels. These efforts emphasize that solving the multimodal-RAG problem is not only ``what to retrieve'' but also ``how to represent and \emph{fuse}'', whether the retrieval works in query-specific timing, streaming detail extraction, or actively selecting complex from audio, visual, and text. This enables the design of future systems for real-time fused retrieval---via aligned foundation modules.

This transition in generation is also linked to scheduling and feature fusion, which is synchronous with the technical and scientific domains discussed in the next subsection, where retrieval scaffolds reasoning for code, robotics, and science. There, retrieval becomes a structured and operational memory that prioritizes --- demonstrated invariant retrieval methodologies important to any high-performance generation.

\subsection{RAG for Code Generation, Robotics, and Scientific Discovery}

The application of Retrieval-Augmented Generation extends far beyond general-purpose question answering, proving particularly transformative in technical and safety-critical domains such as software engineering, robotics, and scientific discovery. In these high-stakes environments, the need for factual precision, domain-specific vocabulary, and traceable reasoning is paramount, which makes the integration of retrieval mechanisms with large language models an indispensable architecture. This subsection explores how RAG frameworks address the unique challenges of code generation, autonomous systems, and knowledge-intensive scientific research, with a focus on domain adaptation, the handling of rare tokens, and the integration of structured knowledge representations. This discussion builds naturally on the multimodal foundations established earlier---where retrieval increasingly draws on structured, heterogeneous knowledge---while also setting the stage for the benchmarking challenges examined next, which assess how well such domain-specific systems meet real-world requirements.

In software engineering, RAG systems are increasingly engineered to provide developers with contextually relevant code snippets, API documentation, and repository knowledge. A fundamental challenge in this domain is the semantic discrepancy between a developer's high-level, natural language intent and the specific functional implementation details available across open-source packages. Traditional keyword-based approaches, which often rely on general-purpose search engines, frequently fail to associate task-related language with the concrete functionalities of a target package. To address this, a novel approach constructs a semantic-level ROS Package Knowledge Graph distilled from package text descriptions, leveraging multi-dimensional feature extraction and a domain-specific fine-tuned BERT model to bridge the lexical gap between user queries and package semantics \cite{ref204}. This shift towards embedding domain-specific structures within retrieval itself illustrates not only a technical enabler but also a design philosophy---tightly coupling retrieval to the demands of the target domain---which will later be scrutinized from an evaluation perspective in the next subsection.

Once relevant components are identified, retrieval serves another important purpose: ensuring that LLMs ground their output in validated technical facts. A critical step is the extraction of explicit, structured facts of the form (head entity, relationship, tail entity) from unstructured patents and technical descriptions to create a searchable knowledge base \cite{ref205}. By training sequence classification models to identify entities and relationships within technical text, one can populate a knowledge base with millions of facts, supporting fine-grained retrieval that contextualizes an LLM's output, ensuring that generated design justifications or suggestions adhere to patented prior art and engineering know-how \cite{ref205}. This population of domain-specific knowledge graphs not only provides the non-parametric memory needed to give models awareness of constraints and norms but also replaces shallow similarity retrieval with a more principled representation---a theme that will be echoed in the benchmarking subsection, where evaluating knowledge-grounded generation depends on the structure and fidelity of the underlying retrieval source.

The fields of robotics and autonomous driving raise further challenges, particularly the need for human-understandable explanations in safety-critical contexts. The RAG-Driver model addresses this by incorporating retrieval-augmented in-context learning into a multi-modal large language model to produce driving action explanations and control signal predictions, grounding generation in retrieved expert demonstrations to offer justifications for decisions \cite{ref206}. This approach also tackles data scarcity and domain gaps in autonomous driving datasets by leveraging non-parametric memory, enabling effective in-context learning without additional fine-tuning \cite{ref206}. A parallel concept is found in the Robot Skill Graph, which organizes an agent's behavioral skills as nodes and their relations as edges, supporting a quadruped robot in inferring how to recombine and adapt fundamental skills to new tasks \cite{ref207}. Here, retrieval becomes closest not to text but to actionable knowledge---enabling flexible and interpretable decision-making in unstructured scenarios, which aligns with evaluating how robustly a system can navigate novel queries, an assessment increasingly covered by benchmark tasks that measure adaptability in unseen or multi-step settings.

Scientific discovery is another area where RAG is essential, given the rapid growth of data across publications, code repositories, and experimental results. Systems must integrate heterogeneous information sources into structured memory for knowledge-intensive use cases, often going beyond simple passage lookup. For example, pipelines aimed at creating computable scientific models from code and natural language attempt to use domain knowledge and data-driven extraction to identify concepts from code, equations from documents, and semantic annotations \cite{ref208}. This creates a knowledge graph of executable models, enabling users to gain context and query the underlying information. Similarly, a pipeline applied to NASA Science Mission Directorate data constructs knowledge graphs that capture abstract relations and semantics across datasets from many disciplines \cite{ref209}, supporting dataset search engines and allowing researchers to trace indirect connections that are invisible to ordinary search. This mirrors a broader trend from text retrieval to structured knowledge retrieval, where the act of finding relevant output is inseparable from determining the right source---an important issue that benchmarking must later quantify using composite tasks and grounded generation metrics.

In highly specialized software engineering and robotics domains, RAG further operates under unique constraints. For instance, in low-code programming for mechatronic systems, a knowledge-graph-driven approach can capture software artifacts and the technical and physical constraints of the environment, reducing the burden on human domain experts for creating reliable programs \cite{ref210}. While this line of work focuses on complexity management rather than retrieval-augmented generation output, it embodies the underlying philosophy of using structured representations to support problem-specific computation. Effective RAG here relies not just on constructing a graph, but on retrieving the relevant nodes in an operationally efficient and scalable manner, providing a resource for future generative tasks or downstream decision support.

Taken together, the technical and scientific settings illustrate how RAG is evolving from simple passage-level search toward an architecture built around structured knowledge, skill graphs, and heterogeneous corpora. Rather than retrieving whole documents, these systems selectively interact with verified facts, executable components, and behavioral routines, emphasizing a deep coupling between retrieval representation, domain constraints, and generative robustness. This integrated perspective underscores the need for controlled benchmarking that can measure such multi-faceted performance, especially when systems must produce faithful, well-grounded responses under safety-critical and resource-aware constraints. These diagnostic requirements are precisely what the next subsection examines, using tailored benchmarks to delineate failure modes and headrooms across retrieval, synthesis, and generation.

\subsection{Benchmarking RAG Systems: Task-Specific and General-Purpose Evaluation}

Benchmarking is foundational to advancing Retrieval-Augmented Generation (RAG), as it enables objective comparisons of systems, drives methodological improvements, and exposes persistent limitations across the retrieval, augmentation, and generation pipeline. While the previous discussion highlighted how domain adaptation can specialize RAG components, benchmarking provides the essential mechanism for verifying that such specialization aligns with real-world requirements, rather than merely optimizing for generic metrics. Similarly, the evaluation strategies discussed here lay the groundwork for the subsequent analysis of domain adaptation by clarifying the specific failure modes---such as grounding errors, hallucination, and multi-hop synthesis gaps---that training and fine-tuning techniques are designed to mitigate. The evaluation strategies adopted in the RAG community range from general-purpose frameworks designed to assess a complete RAG pipeline to domain-specific benchmarks that target particular task families such as multi-hop reasoning, long-form grounded generation, and Chinese-language CRUD operations. A common theme is the realization that conventional question answering (QA) benchmarks, which often evaluate retrieval and generation in isolation or assume simple extractive answers, are insufficient to capture the complexities inherent to modern RAG deployment---including context grounding, faithfulness, robustness to irrelevant passages, and seamless interaction between retriever and generator. As a result, new benchmarks have been intentionally crafted to expose weaknesses, align evaluation metrics more closely with downstream requirements, and inform the kinds of module-level interventions discussed in both preceding and following subsections.

Among the most prominent general-purpose benchmarks is ClapNQ, which addresses a central gap in RAG evaluation by covering the full pipeline, including retrieval, generation, and their integration for long-form question answering \cite{ref211}. It builds upon the Natural Questions corpus but introduces a distinctive twist: the target responses are cohesive long-form answers that are roughly one-third the size of the original passage and that combine multiple non-contiguous text segments. This design moves beyond verdict benchmarks that only require locating evidence, emphasizing a style of grounded generation where synthesis and cohesion are necessary---echoing the multi-functional knowledge requirements explored earlier in this survey, and anticipating the domain-adaptive training discussed next. ClapNQ thus exposes critical space for improvement in systems intended to both retrieve and synthesize passages into a non-verbatim answer, providing a rigorous setting where extractive copying or superficial segmentation are insufficient \cite{ref211}.

Complementing ClapNQ is the MultiHop-RAG benchmark, which targets the more challenging setup of multi-hop questions requiring retrieval and reasoning over several distinct pieces of evidence \cite{ref8}. The creators identify a key deficiency in existing RAG datasets: they rely on simple factoid-style queries requiring only one evidence segment. By constructing a knowledge base from English news articles and generating questions that require multiple supporting evidence, MultiHop-RAG exposes weaknesses in standard embedding models when scaled to multi-document reasoning. In their empirical evaluations, they reveal that state-of-the-art embedding models often fail to retrieve all necessary evidence in a single pipeline, and benchmark questions whether even large LLMs like GPT-4 and Llama2-70B, when provided with genuinely relevant independent chunks can reliably synthesize multi-step answers. This inevitably consolidates this benchmark's role in measuring both retrieval effectiveness and the multi-step reasoning ceiling of LLMs deployed for RAG.

Another notable benchmark is CRUD-RAG, which expands the assessment of RAG systems beyond simple question-answering to four core content maintenance operations---Create, Read, Update, and Delete \cite{ref51}. This comprehensive Chinese benchmark represents a diversity of compounding tasks by preserving knowledge and assessing diverse use cases---creating original content, responding to complex queries, correcting inaccuracies, and condensing lengthy texts. Critically, CRUD-RAG evaluates not only the LLM's answer generation but also the retrieval component, the length of context provided, and the overall budget of the external knowledge base, providing fine-grained failure attribution often neglected in English-centric benchmarks. This holistic, component-aware perspective reflects the architectural concerns faced by industrial RAG systems, thereby setting a scene for the domain-adaptation strategies that later subsections will examine.

The emphasis on multilingual robustness is captured by NoMIRACL, a human-annotated dataset spanning 18 languages designed to test a model's ability to gauge when it lacks necessary information \cite{ref125}. The dataset is organized into two subsets: non-relevant and relevant. It introduces two key metrics: hallucination rate when a model becomes adversely critically influenced by non-relevant text, and error rate when relevant context is provided but not used. This setup reveals a fundamental tension within current models---those too eager to answer excessively hallucinate heavily, while overly conservative models struggle to provide correct responses. Their experiments with multilingual LLMs including Llama-2, Orca-2, and Mistral demonstrate that most models balance the ability to recognize relevant passages poorly, with hallucination rates exceeding 88\% in some models, while GPT-4 offers the best tradeoff but still leaves room for improvement. Failure analysis in these cases indicates that when less popular or domain-specific knowledge is present, context grounding failures dominate---directly motivating the RAG fine-tuning and hybrid-training solutions examined next.

Complementing these application-oriented benchmarks, a wave of analysis has emerged to understand the configurational capacities and utilization habits of RAG systems. For instance, the RAGGED framework analyses how often retrieval can help. It finds that decoder-only models often fail to robustly operate beyond 5 documents and tend to rely on memorized knowledge, whereas encoder-decoder models are more context-dependent and more sensitive to retrieval quality \cite{ref7}. Such findings speak to the utility of calibrating how the retrieval is configured, deciding the passage budget as a hyperparameter, and highlight the interdependent design choices that influence both retrieval and generation, as well as make obvious the need for refinements that later subsections will explore.

On the other hand, recent work also stresses the importance of utility---not just relevance---when judging passages \cite{ref53}. Utility, in this context, signifies the ability of a passage to directly support the answer, requiring a distinction between information that merely topically overlaps and information that actually satisfies the task, often considering careful ranking. The study suggests a listwise approach to mitigate the LLM's sensitivity to the order of input passages. This kind of meta-evaluation is essential, as improved retrieval fidelity can fail to map to payoff in final answer faithfulness, necessitating more sophisticated domain-oriented metric design.

Regarding general architectures, the evaluation of retrieval and generation often reinforce and expose the interconnection between components. MultiHop-RAG, for example, reveals that retrieval and generation components can combine poorly for multi-step queries, even when the right evidence is selected \cite{ref8}. This type of failure attribution can guide researchers to better design these systems rather than over-optimizing one component over the other, while also illustrating gaps that domain-specialized training often seeks to address.

Overall, the design of these RAG benchmarks underscores a shift beyond measuring retrieval accuracy and answer plausibility towards evaluating more granular yet integral attributes: faithfulness, robustness to intermediate reasoning, and the ability to reliably combine multiple relevant reasoning chains. Metrics tailored to these concerns often require custom datasets, layer-wise debugging, and alert analysis---highlighting exactly what domain adapted approaches must do to advocate for their methods. By clearly mapping failure points to retrieval ing, generation, and integration, these benchmarks provide critical evidence that the remaining headroom in RAG lies both in domain-appropriate adaptation (discussed next) and in carefully aligning retrieval to generation in robust and transparent ways across scales and data distributions.

\subsection{Fine-Tuning, Synthetic Data, and Domain-Adaptive Training for RAG}

Domain adaptation remains one of the most critical practical challenges in deploying Retrieval-Augmented Generation (RAG) systems. While the core RAG framework demonstrates strong performance on general knowledge bases such as Wikipedia, its effectiveness degrades significantly when applied to specialized domains characterized by unique vocabularies, document structures, and reasoning patterns. Building on the benchmarking insights from the previous subsection---which identified grounding errors, hallucination, and multi-hop synthesis gaps as persistent failure modes---this subsection examines the training and fine-tuning strategies designed to directly mitigate these deficiencies in specialized contexts. Addressing this gap requires sophisticated training strategies that extend beyond generic fine-tuning, encompassing domain-adaptive retriever training, generator refinement, and careful orchestration of both components. These approaches lay the groundwork for the entity-centric, long-tail, and personalization challenges discussed next, as many of the same techniques---such as synthetic data generation and parameter-efficient adaptation---are extended to handle rare entities and user-specific knowledge.

A foundational observation in this area is the critical role of joint retriever-generator training during domain adaptation, directly addressing the retriever-generator misalignment that benchmarks such as ClapNQ and MultiHop-RAG exposed. The original RAG framework was developed and optimized primarily with Wikipedia-based knowledge, leaving its dual components unaligned when transferred to domains like healthcare or news. To address this, an extension known as RAG-end2end has been introduced, which updates all components of the external knowledge base during training and injects a domain-specific auxiliary reconstruction signal. This auxiliary signal forces the model to access and reconstruct relevant information from the external KB, thereby enhancing its domain awareness. Experimentation across COVID-19, news, and conversational datasets demonstrates significant improvements over the standard RAG baseline, highlighting the value of synchronized joint training in specialized domains \cite{ref78}. This perspective underscores that static retriever-generator alignment is suboptimal when domain shift is prominent.

When considering the target task of adapting LLMs to less-popular or low-frequency knowledge---a key concern anticipated by the long-tail discussion in the following subsection---a comparative analysis of RAG and fine-tuning reveals important trade-offs. In knowledge-intensive question answering tasks, the performance of LLMs typically degrades when confronted with unpopular entities or domain-specific concepts. A systematic study comparing RAG and supervised fine-tuning over synthetic data across entities of varying popularity reveals that fine-tuning significantly boosts performance across all frequency groups, particularly in both the most and least popular segments. However, RAG still surpasses all other methods in accuracy overall \cite{ref77}. This indicates that for long-tail and rare-entity knowledge, retrieval augmentation provides a distinct advantage, while fine-tuning can solidify a baseline. Another relevant study demonstrates that RAG consistently outperforms unsupervised fine-tuning, not only for existing knowledge but also when incorporating entirely new facts, with the caveat that models struggle to learn new factual information via unsupervised fine-tuning \cite{ref74}. In low-resource and zero-shot settings, the choice between fine-tuning and RAG can also be influenced by the practicality of black-box models, as demonstrated by retrieval-based alternatives that interpolate outputs from a frozen LLM with a datastore from the target domain, showing that such retrieval-centric methods can outperform fine-tuning when labeled data is scarce \cite{ref212}.

The proliferation of embedding models as black-box services adds significant complexity to fine-tuning during RAG domain adaptation. Since fine-tuning the internal parameters of an embedding model is often infeasible in these specific contexts, model-augmented fine-tuning (Mafin) emerges as a novel solution. Mafin incorporates a trainable embedding model on top of the black-box variant, enhancing retrieval performance without requiring access to the black-box model's internals, demonstrating a clear pathway for domain adaptation using API access and trainable parameters \cite{ref54}. This approach is a practical answer for industries that deploy RAG systems as services and do not expose internal model weights---representing a level of adaptation that aligns with the resource constraints and robustness concerns reflected in the benchmarks discussed earlier.

Hybrid approaches and pipeline construction are increasingly important in modern RAG systems. Prior research comparing fine-tuning and RAG pipelines using an agricultural dataset shows that fine-tuning and RAG provide cumulative benefits, with fine-tuning producing a significant accuracy increase and RAG notable further gains, ultimately increasing answer composite for questions that require geographic-specific knowledge \cite{ref213}. Similarly, the FinanceBench Q\&A advocating system highlights that the combination of retriever fine-tuning and LLM fine-tuning yields higher accuracy on answering financial filings, and incorporating iterative reasoning on top of the pipeline yields even larger performance gains, thereby approaching expert-level quality \cite{ref214}. This points to an emerging consensus that the most effective specialization methods are hybrid---not selecting between retrieval or tuning---but orchestrating both to complement their respective strengths, while also feeding that aggregated increase to the type of personalized and scarce-knowledge settings discussed next.

However, fine-tuning and RAG are not always intrinsically compatible. An analysis across multiple generative language models such as GPT-J-6B, OPT-6.7B, LLaMA, and LLaMA-2 using several metrics (ROUGE, BLEU, METEOR, and cosine of similarity) reveals that when fine-tuned models are directly combined with RAG, there is notable performance degradation in some cases \cite{ref215}. This suggests that when robust retrieval is present, it often outperforms fine-tuning in grounding deal and reducing hallucination---outscoring fine-tuned models by up to 16\% ROUGE and 15\% BLEU---but also highlights that any hybrid combination should not be treated as a one-size-fits-all solution. More recent parameter-efficient adjustments emerge as a more effective domain adaptation as alternative.

Parameter-efficient fine-tuning techniques have become a key building block for domain adaptation in RAG, as they emphasize memory efficiency and modularity. Techniques such as LoRA and QLoRA allow RAG systems to achieve domain specialization by fine-tuning only low-rank adapters while the base model remains unchanged. Moreover, using a Quantized Influence Measure (QIM) as an AI-Judge mechanism for the retrieval selection and incorporating user feedback directly during training leads to a fine-tuning enhanced RAG system that adapts continuously to user expectations while remaining globally efficient on the order of few parameters \cite{ref216}. These parameter-efficient methods drastically reduce the computational cost of domain adaptation and, crucially, enable efficient adaptation without full model knowledge, which relates directly to determining the kinds of personal and real-time updates required by RAG systems.

Synthetic data generation has proven to be a necessary component of domain-adaptive training, particularly when there is no labeled data. Multiple pipelines construct QA pairs by extracting parts of a document, generating questions from segment content, and then using them for fine-tuning distinct stages of retrieval or generation \cite{ref213}. In domains where annotations are few, supervised synthetic QA datasets allow unsupervised fine-tuning to achieve performance increases as substantial as full-model adaptation but with fewer computational or constraints constraints on privacy. This is also the approach that the survey integrates with the next subsection, where synthetic data and fine-tuning are used.

Taken together, these findings indicate that domain-adaptive training for RAG systems has evolved from simple full-model fine-tuning into a modular ecosystem: joint retriever-generator updates \cite{ref78}, fine-tuning/RAG hybrids \cite{ref217} and \cite{ref74}, black-box embedding augmentation \cite{ref54}, parameter-efficient fine-tuning \cite{ref216}, and robustness comparisons across multiple Language models via precise metrics \cite{ref215}. The development of successful domain-specialized RAG systems require careful balance of retrieval grounding, parametric knowledge, and efficient adaptation loops---without harming the complementary strengths of each component. As more work moves toward real-world deployment, the joint fine-tuning and retrieval augmentation will be central, just as the following subsection serves as a bridge to the critical next steps for personalized and long-tail knowledge integration.

\subsection{Entity-Centric, Long-Tail, and Personalization in RAG Systems}

Entity-centric, long-tail, and personalization challenges represent a critical frontier in the deployment of Retrieval-Augmented Generation (RAG) systems, addressing the gap between generic question-answering capabilities and the nuanced requirements of real-world, user-specific, and rare-knowledge applications. While the preceding discussion on domain adaptation focused on aligning retrievers and generators for specialized fields, this subsection shifts attention to a closely related yet distinct set of obstacles: how RAG systems can be tailored to individual users, hierarchical entities, and uncommon knowledge---all of which demand more than generic fine-tuning or retrieval pipelines. This subsection synthesizes the key innovations designed to tackle these intertwined challenges, focusing on structural representations for entity awareness, multi-source augmentation for personalization, and retrieval fine-tuning for rare knowledge.

To begin with, a primary challenge for standard RAG systems is the effective handling of queries centered on specific, often hierarchical, entities within private or enterprise knowledge bases. Generic retrieval mechanisms may fail to correctly associate a query with the appropriate level of an entity hierarchy, leading to imprecise or incomplete context. To address this, the T-RAG system proposes a `Tree-RAG' architecture, which incorporates a tree structure to explicitly represent entity hierarchies within an organization \cite{ref218}. This method generates a textual description of the relevant entity's position and relationships within the hierarchy, which is then appended to the context provided to the generator. Designed for real-world deployment on private enterprise documents with constraints like on-prem hardware, T-RAG demonstrates that combining a fine-tuned open-source LLM with a RAG framework augmented by explicit entity hierarchy awareness yields superior performance compared to naive RAG or fine-tuning and that this holds particularly true when the retrieval process must navigate interdependent organizational data \cite{ref218}. This strategy for entity-aware augmentation emphasizes that addressing the nuances of organizational data often requires more than just optimized text chunking and retrieve; it requires a structural representation that mirrors the domain's fundamental logic, thereby enabling the LLM to reason over the correct subset of knowledge without being overloaded by irrelevant but superficially similar data.

Building on the need for richer and more contextual inputs, personalization and multi-source knowledge integration pose a significant next-level challenge. A powerful approach to personalization involves leveraging user profiles and historical interactions in tandem with general knowledge bases---but these sources are rarely all equally relevant for a given query. To address the complexity of incorporating multiple data sources, generated responses in personalized dialogue systems can be improved by processing compositionally. For instance, UniMS-RAG unifies Knowledge Source Selection, Knowledge Retrieval, and Response Generation into a single sequence-to-sequence framework \cite{ref219}. This method uses acting tokens to plan and evaluate the response process \cite{ref219}, moving beyond simple concatenation to a more dynamic and adaptive use of multi-source augmentations \cite{ref219}. Such an approach reflects a paradigm shift from static, one-time augmentation to adaptive evaluation of retrieved evidence, where the system controls the retrieval process and evaluates its own outputs in a self-aware loop \cite{ref219}. The system's modularity provides an intuitive pathway to personalization, demonstrating a marked improvement over prior methods \cite{ref219}.

The challenge of long-tail and rare entities further exposes the limitations of LLMs, which tend to perform sub-optimally when handling uncommon concepts. This problem can be addressed through systematic fine-tuning on synthetic datasets---a theme introduced in the previous subsection on domain adaptation. In a suite of empirical studies comparing RAG and fine-tuning, fine-tuning has been shown to effectively enhance model performance across entities of varying popularity, including the least popular ones \cite{ref77}, while RAG can surpass other methods when retrieval is effective \cite{ref77}. Both RAG and fine-tuning efficacy are amplified by advancements in retrieval and data augmentation techniques \cite{ref77}. To improve retrieval when facing specialized domains, model augmented fine-tuning (Mafin) offers a strategy to enhance black-box embeddings by adding trainable components \cite{ref54}. Mafin improves the retrieval pipeline's performance on domain-specific data---insensitive to whether the embedding models are API-accessed or completely frozen---and thereby broadens the applicability of RAG by allowing the adaptation of proprietary embedding models \cite{ref54}. This technique is broader than just improving domain recall; it creates a pathway to robust retrievers for niche knowledge areas, without requiring internal access or exposing weights \cite{ref54}.

Personalization efficiency is however another crucial dimension, especially when dealing with very large user bases. Unlike the model-level adaptation techniques discussed so far, personalization must be specifically optimized for storage, lookup, and response latency. Rather than relying on expensive fine-tuning or entire-corpus retrieval, augmentation can be optimized to store and recall user-specific attributes in a compact and adaptive form. The Persona-DB model serves as an example of this principle, focusing on how user data is represented rather than trying to enrich the retrieval stage directly \cite{ref220}. Persona-DB constructs a hierarchical representation of user data to enhance generalization across contexts and employs a collaborative refinement process to bridge knowledge gaps among users \cite{ref220}. This representation allows for high accuracy even when the number of characters is substantially reduced, a significant advantage in the context of extensive user histories. This suggests that focusing on improving data representation---the crust of the personalization bottleneck---offers a more direct pathway to enhance efficiency and minimize inference latency in personalized RAG, thereby highlighting that database and memory construction, rather than model architecture, can become the core determinant of system performance.

In conclusion, handling entity-centric queries, promoting effective personalization, and retrieving long-tail knowledge require a deliberate departure from simple ``embedding and retrieve'' pipelines. These challenges are addressed through the integration of structured knowledge representations for entities \cite{ref218}, the unification of processes around complex tasks \cite{ref219}, and advanced retrieval and fine-tuning strategies to adapt to specialized, low-frequency domains \cite{ref77,ref54}. The effectiveness of these techniques relies on representing knowledge with enough structured detail, dynamically evaluating what to retrieve, and adapting existing components without full access. As such, future retrieval-augmented systems will increasingly lean toward designs that preserve entity structure, adapt to user preferences implicitly, and are engineered to be efficient and robust across various and scarce real-world data---an evolution that also aligns with the enterprise and high-stakes requirements discussed in the next subsection.

\subsection{Enterprise, Industrial, and Practical Deployment of RAG}

Building directly on the entity-centric, long-tail, and personalization challenges discussed in the preceding subsection, the transition of Retrieval-Augmented Generation (RAG) from academic research to enterprise and industrial production environments marks a critical maturation phase for the technology. While the preceding discussion focused on how RAG systems can be tailored to individual users, hierarchical entities, and uncommon knowledge, this subsection addresses the equally demanding constraints that emerge when these capabilities must be delivered reliably, securely, and cost-effectively in real-world operational settings. The practical deployment of RAG systems introduces a distinct set of challenges primarily centered around privacy, cost, security, and operational resilience. In real-world settings, the primary value proposition of RAG is its ability to facilitate language models with proprietary, private, and often highly sensitive data that cannot be used to train or fine-tune public-facing foundation models \cite{ref221}. This core use case immediately places privacy preservation and data security at the forefront of any enterprise deployment strategy.

The greatest practical challenge for entities is the operation of RAG in privacy-sensitive and cost-constrained environments. The reliance on external knowledge bases introduces a new attack surface, as the stored, non-parametric knowledge becomes a potential target for extraction. Despite extensive research on the privacy risks of LLMs, the RAG technique could potentially reshape the inherent behaviors of LLM generation, posing new privacy issues that are currently under-explored and demonstrated by the vulnerability of RAG systems in leaking the private retrieval database \cite{ref221,ref222}. This is a critical concern, as the confidentiality of the augmented data is often the sole reason for deploying RAG in the first place. Consequently, enterprise architectures must proactively include significant measures to secure the data pipeline from storage to inference, often necessitating the use of hardware-backed Trusted Execution Environments (TEEs) or other enclave-based technologies to isolate retrieval and computation processes. The engineering cost to adapt existing third-party dependencies for such secure hardware is a substantial obstacle, as developers have to invest significant operational costs in porting their stack to adhere to restrictive programming models---such as those required by Intel SGX---to utilize effective memory encryption and hardware-isolated computation \cite{ref223}. This, in turn, impacts overall system design and often leads to hybrid deployment strategies where less-sensitive queries reside on standard infrastructure while confidential operations are handled within secure enclaves.

Practical enterprise deployments often build a RAG system around a ``bring your own vector database'' model, requiring a hybrid integration pipeline that handles evolving knowledge bases. Yet, this challenge is not a one-time routine ingestion, but the continuous building and maintenance of the source of data; ironically, these pipelines often end up mirroring the same challenges as managing a software supply chain. Maintaining relevance and freshness for a RAG system requires the data pipeline with the same rigor as a software supply chain---with regular ``upstream'' updates, auditing for quality and security, and versioning to reflect changes in source documents \cite{ref223}. Whenever a vendor updates its internal documents, product catalogs, or standard operating procedures, the changes must be merged, reflected in the retrieval index, and processed for relevance. Unlike software updates that are activated by a package manager, these data mutations must be analyzed for their impact on search quality and integrated into the system without introducing downtime or service degradation.

Operational constraints in production settings have implications for latency and cost. If a RAG-based system suffers from prolonged high latency or excessive computational cost, it will not provide users with a tangible advantage over a simpler, classical information retrieval system. Enterprise RAG pipelines often involve careful orchestration of multiple components---a retriever, a reranker, and a generation model---and the orchestration of these components introduces multiples of latency. Moreover, the need to process incoming queries in real-time while managing different retrieval indexes, ensuring that the system remains responsive and within budget, is a constant operational burden. This challenge is further amplified when RAG is used for complex tasks like interactive robotics control or secure data exchange, which demand extremely low latency and high up-time \cite{ref224}. Given that enterprise RAG systems typically operate in tandem with other enterprise software, fine-tuning every component of the pipeline for speed, resource usage, and efficiency is a prerequisite for successful deployment.

Prioritizing the reduction of hallucinations and incorrect answers is the optimization target of RAG in production deployments; however, ensuring correctness in the face of noisy documents requires engineering the entire chain of interactions, much like the robustness considerations introduced in the previous discussion on entity-aware augmentation. A RAG pipeline evaluated only in a gold-standard academic setting is severely incomplete: to be robust in production, the system must withstand a wide variety of adversarial and low-level perturbations, such as structural or typographical errors in the documents being ingested. The presence of even a small number of typos or syntactically malformed text within retrieved documents can severely degrade the system's overall performance \cite{ref97}. The stability of the retrieval and generation models when handling irrelevant, incomplete, or conflicting information is a crucial factor for a successful launch. If the retrieval component fails to locate the exact pieces of content that answer a user's query, the RAG strategy becomes ineffective. The system must be designed to ensure that its performance remains stable and can gracefully handle knowledge conflicts, which are inevitable as real-world data evolves over time \cite{ref9}.

Finally, when RAG systems are integrated into regulated industrial environments, both security and availability become paramount, echoing the enterprise-level data access and structural constraints highlighted by entity-centric RAG strategies. As in any real-time control loop where data is continuously arriving, feeding the output of one LLM into another may create hallucinations or unstable behaviors, which in a business setting is not simply an inconvenience but a critical liability. An industrial setting requires an engineering approach that focuses not only on retrieval quality but also on the communication and data security across different stakeholders \cite{ref225}. This has significant implications for the trust decisions an organization is willing to place in cloud services and third-party vendors. For these sensitive environments, the company's infrastructure must be architected as a distributed continuum: from local plant-floor devices to large-scale cloud-based services, with secure channels and the ability to uphold explicit service-level agreements \cite{ref225}. As such, successful deployment of RAG is not merely a matter of choosing the correct algorithm but a systemic engineering effort that balances knowledge base management, system integration, and a deep awareness of the entire security landscape, directly extending the requirement for preserving entity structure and adaptive fine-tuning discussed in the previous subsection. In this way, the future advances in RAG will rely heavily on addressing these security, performance, and operational hurdles to fully unlock its potential in industry and enterprise settings---a transition that also sets the stage for the robust, scalable, and knowledge-aware systems this survey envisions as the next generation of language technologies.

\section{Open Challenges, Unsafe Practices, and Future Research Directions}

\subsection{The Efficiency Bottleneck: Computational and Latency Challenges in Retrieval and Generation}

The integration of retrieval mechanisms into large language models (LLMs) introduces a fundamental tension: while retrieval dramatically enhances factual accuracy and knowledge coverage, it also imposes severe computational and memory overheads that threaten the practical deployability of RAG systems. This efficiency bottleneck manifests at multiple levels, from the long-sequence generation necessitated by knowledge injection to the added latency of retrieving and processing external documents. Moreover, as the preceding discussion on challenges in RAG highlighted, even when systems generate accurate outputs, the interpretability of those outputs requires traceable evidence and transparent reasoning---yet such functionality further amplifies the computational burden. For instance, self-critique mechanisms and credibility-aware generation introduce additional inference steps and require more extensive context windows, directly exacerbating the latency and memory pressures discussed here. Efficient deployment is therefore not merely an engineering concern but a prerequisite for enabling the kind of robust, auditable, and user-trustworthy behavior described in the subsequent section. This subsection synthesizes the primary sources of these inefficiencies and surveys emerging techniques---including system-level caching, speculative decoding, and energy-aware system design---that are necessary to mitigate them.

At the core of the efficiency problem is the memory footprint of the key-value (KV) cache, which scales linearly with the sequence length and batch size. During long-sequence generation, a large amount of transient state information, stored as the KV cache, must be held in GPU memory in addition to model parameters \cite{ref226}. This growth is a major obstacle, as the entire KV cache is loaded for every generated token in the auto-regressive generation process, leading to low utilization of computational cores and high inference latency \cite{ref227}. In the RAG context, this issue is exacerbated because the retrieved knowledge is injected into the prompt, creating a long sequence that drives both computation and memory costs to prohibitive levels \cite{ref58}. The basic design choice of moving data to computation in modern computing systems goes directly against these trends, causing performance, scalability, and energy bottlenecks, since data access is a key challenge for data-intensive applications where memory bandwidth and energy do not scale well \cite{ref228,ref229}. Therefore, the cost of data movement---off-chip to on-chip---is very expensive in terms of bandwidth, energy, and latency, more so than computation traditionally, making it a central challenge in RAG system design \cite{ref230}.

To address the KV cache bottleneck independent of retrieval, one of the leading strategies is token eviction. The observation that many tokens contribute little to the attention score in a sequence has driven research into compression methods. The Heavy Hitter Oracle (H\(_2\)O) leverages the insight that a small portion of tokens---the ``Heavy Hitters''---accounts for most of the value in attention calculations \cite{ref226}. By formulating the KV cache eviction as a dynamic submodular problem and proving theoretical guarantees under mild assumptions, H\(_2\)O dynamically retains a balance of recent and heavy-hitter tokens, reducing cache memory usage by up to 70\% without performance degradation \cite{ref226}. Further, the observation that the similarity between adjacent tokens' query vectors is remarkably high in LLMs such as LLaMA, and the current query's attention calculation can minimize reliance on a stream of preceding queries, powers simpler eviction policies like CORM, which significantly reduces inference memory \cite{ref231}. While such approaches reduce memory, they mainly address the generative component and not the full RAG pipeline latency.

A more comprehensive approach to reducing the latency of RAG specifically is system-level caching of the intermediate states of retrieved knowledge. The RAGCache pipeline effectively organizes the intermediate states of retrieved knowledge into a knowledge tree and caches them across the GPU and host memory hierarchy \cite{ref58}. By proposing a replacement policy that is aware of LLM inference characteristics and RAG retrieval patterns, and dynamically overlapping retrieval and inference steps, RAGCache minimizes end-to-end latency, achieving up to 4× reduction in time to first token (TTFT) and a 2.1× improvement in throughput compared to a state-of-the-art LLM inference system integrated with a vector database \cite{ref58}. This system-level co-design is key to tackling the long-sequence problem because it converts a costly recomputation of data into a cached state, and it operates synergistically with other latency reduction techniques.

Beyond caching, techniques based on speculative decoding provide another layer of acceleration. TriForcing extends the concept of speculative decoding to the long-sequence generation domain by introducing a hierarchical system that scales to long contexts \cite{ref227}. Instead of relying solely on a smaller draft model (and its associated overhead), TriForce leverages the original model weights and dynamic sparse KV cache via retrieval as an intermediate draft layer, which is itself speculatively drafted by a smaller model. This hierarchy allows TriForce to achieve impressive speedups (up to 2.31× for Llama2-7B-128K on an A100 GPU) while remaining lossless regarding the output distribution, outperforming competing offloading baselines such as DeepSpeed-Zero-Inference \cite{ref227}.

In broader data-centric terms, the principles of processing-in-memory (PIM) and near-data processing offer a high-leverage alternative to moving data to computation. PIM places computational mechanisms in or near where the data is stored, and thereby can reduce the cost of data movement \cite{ref228,ref230}. This approach is directly relevant for retrieval, such as in the similar similarity search, a key primitive in RAG. Application-driven near-data processing for similarity search demonstrates how high-capacity associative memory can execute distance calculations and top-k sorts in-memory with high bandwidth and energy efficiency compared to multicore CPUs \cite{ref232}. This suggests that future RAG pipelines could offload the retrieval stage entirely to memory such as PIM-capable accelerators, to avoid the system-wide communication bottleneck, as data movement remains a primary source of energy and latency, as shown in analysis of disaggregated memory systems \cite{ref233,ref229}.

System-level optimization frameworks that focus on the computation side aid to address latency when hardware bottlenecks persist. For instance, frameworks such as Cello use predictive early termination and censored regression to reduce the cost of collecting slow configurations in system optimization, improving latency and energy efficiency \cite{ref234}. Similarly, targeting edge devices, EdgeBERT implements entropy-based early exit prediction to perform dynamic voltage-frequency scaling (DVFS) at a sentence granularity, achieving 53× lower energy consumption compared to mobile GPU baselines without sacrificing accuracy \cite{ref235}.

The persistence of data movement across the system, from storage to memory and processor, has further driven research into computational storage. Systems like Skytether decompose queries across paradigms and aim to minimize movement of data while accumulating compute capabilities near storage hardware \cite{ref236}. For RAG, this dynamic partitioning of data and computation ensures that long-sequence generation and knowledge retrieval do not become I/O bottlenecked, while avoiding over-provisioning of traditional memory capacity.

In conclusion, addressing the computational and latency challenges in RAG requires a combined approach: algorithmically compressing the KV cache and enhancing long-sequence inference through speculative decoding, and architecturally rethinking memory scheduling and data placement to use caching and near-data processing. While early breakthroughs in system caching \cite{ref58} and hierarchical speculative decoding \cite{ref237} are significant, achieving performance parity with non-RAG LLMs in standard deployment will be contingent upon scaling these models across increasingly distributed, and memory-bound, computing platforms. As the next subsection will explore, the efficiency gains outlined here directly underpin the viability of implementing richer explainability mechanisms---such as reflective tokens and credibility checks---which would otherwise become prohibitively expensive in practice. Thus, optimizing for speed and memory is not solely about performance metrics; it is a foundational step toward trustworthy and reliable RAG systems.

\subsection{The Quest for Explainability and Interpretability in RAG Systems}

While Retrieval-Augmented Generation (RAG) has emerged as a powerful paradigm to enhance the factual accuracy and knowledge capacity of large language models, it simultaneously introduces a profound interpretability paradox. On one hand, RAG ostensibly offers a pathway towards greater transparency by grounding model outputs in verifiable, external sources; on the other, the intricate interplay between parametric memory and non-parametric retrieval creates new ``black boxes'' that are often more difficult to explain than the components they augment. This paradox is further compounded by the efficiency challenges outlined in the previous subsection: the very optimizations---such as KV cache compression, speculative decoding, and system-level caching---that make RAG computationally tractable also obscure the internal reasoning trace, as they operate on transient or evicted states that are no longer fully inspectable. Furthermore, as the storage, retrieval, and ranking of documents become increasingly abstracted and accelerated---sometimes offloaded to near-data processing units---the need for explainability becomes even more pronounced and harder to satisfy. Conversely, as the following subsection will argue, the systemic fragility of RAG's knowledge sources necessitates even greater interpretability: if a retrieved source is stale, biased, or has been poisoned by feedback loops, the user must be able to see and challenge the evidence trail rather than blindly accept the output.

The primary interpretability challenge in RAG systems lies in disentangling the exact contribution of each knowledge source and understanding how they interact to produce a final output. Traditional language models process information through a deterministic lens of parametric weights, whereas RAG introduces a second, heterogeneous stream of information---retrieved documents---whose selection, ranking, and integration are governed by often-opaque similarity metrics and contextual relevance assessments. This dual-stream nature obscures the reasoning trace, making it difficult for users and developers to answer the core question: \emph{why} did the system generate this particular response based on \emph{which} specific evidence?

Transparency in knowledge integration has been a multi-faceted issue, encompassing each phase of the RAG pipeline. Naive RAG systems often retrieve documents through lexical matches or dense vector similarities without providing a clear rationale for the selection or ranking of each passage \cite{ref1}. Without a clear trace linking generated claims through composition to their specific source passages, users cannot judge whether the answer is a faithful reasoning over evidence or a conflation pretending to be grounded.

The need for explainability extends beyond mere internal auditability; user trust and credence generation fundamentally depend on it. When a system presents a response augmented by a citation, the implied contract is that the citation contains the evidence for the claim. However, this trust is fragile. As research indicates, conflicts between an LLM's internal parameterized knowledge and retrieved evidence are common, and even when correct context is provided, a model might still fail to update its incorrect prior \cite{ref238}. This highlights a failure in auditable decision-making that is central to the ``trust-busting'' paradox. Without proper interpretability, users cannot identify \emph{when} such a conflict occurs or \emph{why} the system defaulted to parametric memory over the ground truth, making it impossible to correct faulty behavior in practice.

To address this, the field has begun to move from ``traceability-centric'' post-hoc explanations toward self-critique mechanisms designed to embed awareness of its own evidence usage. The framework \cite{ref109} embodies a significant departure from static retrieval-augmented pipelines. Instead of blindly retrieving a fixed number of passages, Self-RAG deploys reflection tokens, specialized output tokens that allow the model to generate self-evaluations on which the system can make decisions and adaptations regarding retrieval needs \cite{ref109}. This moves interpretability from an external audit exercise into an intrinsic generation component, allowing the system to confirm---through a self-assessment of practical benefit to the equations themselves. This shift is a profound step toward auditing the model's evidence usage and making its reasoning more transparent.

Experiments in specialized areas such as biomedicine have shown similar promise, introducing domain-specific models like Self-BioRAG with ``customized reflective tokens'' that indicate whether a claim is supported by retrieved documents or by domain-specific knowledge, creating a clear and auditable trail of epistemic justification \cite{ref176}. This is particularly crucial in high-stakes fields where incorrect conclusions can lead to user harm or systemic errors, and it explains not only \emph{what} the output is but \emph{how} it reached a conclusion.

Alongside self-critique, \textbf{credibility-aware generation} has emerged as a complementary paradigm. Classical retrieval augmentation risks treating all retrieved documents as equally trustworthy. \cite{ref239} directly addresses the negative impact of low-quality passages. They propose training models on objective (i.e., risk-aware) behaviors that enable them to strike a balance of trustworthiness in the generated output, empowering the user to see the pathway of the model's judgment \cite{ref239}. This aspect is especially urgent in the context of the enhanced systemic fragility, as RAG systems may inadvertently retrieve hallucinated or synthetic content from flawed sources and propagate it as authoritative.

Beyond the correctness of generation itself, explainability requires the ability to assess whether a retrieved passage is actually \emph{useful} for the generation. This is a skill that a well-instructed LLM can exhibit by distinguishing between the relevance and utility of a passage \cite{ref53}. A further complication arises when retrieval---a high-level judgment that is not directly aligned with relevance---causes the generator to over-rely on a low-ranked passage or even a single retrieval error. Thus, assessing the \emph{utility} alongside \emph{relevance} is essential for the system's overall trustworthiness.

The technical interpretability challenge is tightly coupled with practical debugging. Developers need the ability to locate a defect behind a RAG failure by understanding which component misbehaved in an untraceable way \cite{ref51}. Providing a system-level view of the retrieval trajectory, ranking scores, and the finally gathered set is foundational for error analysis.

To provide that visibility, we need more than a final output---we need provenance (the retrieved set), attribution (the relevance vector), and introspection (the reflective tokens). Without these, the retrieval system's influence becomes a hidden oracle inside the model. While transparent components are useful, they are not sufficient alone: the challenges of stale or biased sources highlighted in the upcoming section underscore that transparency must be coupled with \emph{robustness} to degraded, malicious, or simply uncurated external data. Only then can the user or developer trace how evidence flows, spot corrupt, and intervene.

Existing approaches are beginning to address multi-hop and decomposable evidence. In fact, the reasoning problem is especially acute when retrieval must combine multiple siloed pieces of evidence across a multi-step chain, as is typical in knowledge-dense domains like law or medicine. For instance, a framework such as RAGAR (Retrieval-Augmented Reasoning for Misinformation Detection) decomposes claims and collects sub-evidence to support a final verdict but leaves the aggregation mechanism opaque \cite{ref240}. To that end, practitioners have proposed multi-view frameworks that operate on \emph{intent-aware query rewriting} to retrieve from multiple angles, but their composition logic remains a challenging interpretability target for future exploration \cite{ref182}. Providing insight into these decomposed reasoning paths is the next frontier.

Ultimately, the search for causal and counterfactual interpretability---the ability to ask ``what if this document was absent?'' or ``what if the retrieved evidence had been different?''---represents the frontier for explainable RAG. Awareness of these missing links not only increases transparency; it also facilitates debugging and reduces risk. It is not simply about ``peering inside'' the model, but rather enabling open, verifiable communication between the user and the external indexed world, while ensuring that the \emph{user} remains in the loop. By combining the self-reflection capability of reflective tokens with the reliability insights from credibility-aware generation and next-generation interface for inspection, RAG systems can truly offer both performance and integrity---trust that is well earned and explainable.

\subsection{The Fragility of Static and Dynamic Knowledge: Long-Tail, Freshness, and Sustaining Knowledge Bases}

The reliance of Retrieval-Augmented Generation (RAG) systems on external knowledge bases introduces a fundamental fragility that is often overlooked in the pursuit of improved accuracy and reasoning. While the previous discussion has highlighted how the tug-of-war between parametric memory and retrieved evidence shapes the trustworthiness of individual outputs, a deeper, systemic vulnerability emerges from the inherent dynamism of the real world and the corresponding static, incomplete nature of the knowledge sources these systems depend upon. RAG mitigates some issues of parametric knowledge staleness in LLMs, yet it creates a new dependency on a brittle infrastructure that must be meticulously curated and updated. This subsection explores the open challenges of static and dynamic knowledge within RAG, focusing on long-tail coverage, the constant need for freshness, and the precarious feedback loops that threaten sustainable knowledge creation---all of which are further vitiated by the credibility and overconfidence issues discussed in the preceding section.

The first critical dilemma is the static nature of knowledge bases in contrast to their dynamic intent. Knowledge Bases (KBs) are not monolithic, static repositories; they are deeply unstable and highly dynamic, constantly subject to change as the real world evolves \cite{ref241}. This instability poses a central problem for RAG systems, which often rely on snapshots of these KBs for downstream tasks such as open-domain QA. As the real-world changes---a company's board member changes, a political leader loses office, or a scientific discovery overturns a previous theory---the KB must be updated to remain relevant. However, the task of updating these repositories is non-trivial. Works have long investigated the nature of such evolution and contraction for formal knowledge bases, highlighting that incorporating new information while discarding obsolete knowledge is a complex logical task \cite{ref242}. RAG frameworks that do not account for these updates risk retrieving outdated information, thereby propagating the same factual errors they aim to mitigate. During pretraining, a major issue is the decay of world knowledge in LLMs and the difficulty of ``knowledge editing'' to align them with up-to-date facts \cite{ref243}. While RAG offers a shortcut by promising retrieval instead of constant retraining, the underlying repository must itself be maintained with high fidelity.

Specifically, evaluating and filtering for stability is as important as adding new facts. Research on KB stability has delved into how real-world evolution differs from simple data corrections or delayed insertions, showing that prediction models can anticipate which facts or properties are likely to change, achieving considerable F1 scores on change prediction tasks \cite{ref241}. However, this predictive ability is not absolute, and entities previously deemed ``unstable'' may change later, highlighting the unpredictable nature of the domain. For RAG systems, this creates a critical tension: retrieving a current fact versus an outdated one. The difficulty is that the correct answer requires temporal understanding (e.g., ``Einstein-children'' is stable, but ``Ronaldo-occupation as of 2022'' is not) \cite{ref241}. RAG systems often struggle to perform the required temporal reasoning to address such queries, making them susceptible to the sorts of epistemic misalignment that users may not easily detect, amplify the trustworthiness concerns of the previous subsection.

Beyond temporal instability, RAG systems struggle with the ``long-tail'' of knowledge. The most frequently discussed or archived facts are easy to retrieve, but the long-tail consists of rare, obscure, or novel entities that are sparsely represented or entirely missing from reference KBs \cite{ref244}. This fragility is paramount when serving user queries about niche scientific discoveries, recent startups, or local politicians; while a parametric model may hold ``vague'' memories, the retrieval module is often unable to fetch specific passages. Consequently, the generation module must fall back on its internal knowledge, leading to more hallucinations or unsupported claims. This issue is aggravated in fast-evolving, high-stakes domains such as medicine or law, where the volume of novel information can outpace the curation speed of static sources, and where the reliability of provider-rated ``utility'' and credibility is limited, as noted earlier.

The necessary but expensive nature of ``freshness'' further introduces an underexplored danger: the risk of harmful feedback loops affecting the quality of the very sources the system uses. A central and emerging concern is the influx of AI-generated content onto the internet. Future RAG systems may retrieve data from a web increasingly synthesized by their predecessors, conflating the source authority and exacerbating the credibility judgment problems described in the Previous Subsection. While some approaches mitigate this through ``repeatedly referencing'' the same source, there's limited attention to the deterioration of link and source stability. Studies show that a vast percentage of links in open collaborative environments decay and disappear over time; remote external domains and software homepages are especially vulnerable \cite{ref245}. As RAG systems generate and disseminate new content, the long-term integrity of the web content ecosystem is threatened, creating a feedback loop where synthetic content generated today is used as a reference in future retrieval. If a flawed generative response is used as a reference, future retrievals will be augmented with hallucinated content, thereby cementing misinformation into ``ground truth'' and amplifying the overconfidence and misinformation propagation highlighted earlier. These generative feedback loops do not simply add noise; they actively poison high-quality sources, turning stable sources into deteriorating ones.

The sustainability of knowledge bases can be examined through social and community dynamics, which offer a parallel to the need for human \emph{curation}. After all, the sustainability of any content system depends heavily on the trust of its contributors, a theme that resonates with the earlier emphasis on credibility-aware generation \cite{ref239}. Findings from community dynamics show that active communities develop stable, interconnected cores, but degrade when trust between core members is broken \cite{ref246}. For RAG, this underscores that web knowledge is maintained by human communities, and their trust signals are essential for robust reliability and long-term quality.

In sum, this systemic fragility---arising from static KBs, the rare long tail, and destructive feedback loops---must be addressed head-on to ensure RAG's promises of grounded generation are sustainable. Future research must develop retrieval models that handle long-tail scarcity robustly, build context-aware, self-correcting mechanisms that treat sources as dynamic, and design governance protocols that mitigate harmful feedback loops. Without such advances, RAG systems risk becoming brittle, self-consuming engines, where the very infrastructure they rely on decays under the weight of \cite{ref247}{]}.

\subsection{The Trustworthiness Dilemma: Overconfidence, Hallucinations, and Misinformation Propagation}

The integration of retrieval mechanisms into large language models (LLMs) has been widely championed as a primary solution for mitigating hallucinations and grounding outputs in verifiable facts. However, a deeper examination of RAG systems reveals a profound trustworthiness dilemma: while retrieval can enhance factual accuracy, it can also inadvertently reinforce model overconfidence, propagate misinformation, and create new failure modes that are more subtle and therefore more dangerous than the hallucinations they aim to correct. The central tension lies in the interaction between a model's parametric knowledge (its internal prior) and the non-parametric evidence presented in retrieved documents, and the system's uncritical acceptance of retrieved context as ground truth. This dynamic of internal parametric memory versus external evidence---and the resulting overconfidence and misinformation risks---forms the foundation upon which the systemic fragility of static and dynamic knowledge bases, as well as the eventual security vulnerabilities of the RAG pipeline, must be understood.

At the heart of the trustworthiness problem is the ``tug-of-war'' between a model's internal prior and retrieved evidence. Research has systematically analyzed scenarios where these two sources conflict, revealing that LLMs are not simply passive recipients of retrieved data \cite{ref238}. When correct retrieved information is provided, models can fix many of their initial mistakes, achieving high accuracy. Yet, when retrieved reference documents are perturbed with incorrect values, the model's likelihood of reciting false information is critically dependent on the strength of its prior. Specifically, when the model's internal knowledge is weak, it is highly susceptible to being led astray by wrong documents; conversely, when the prior is strong, the model may ignore the contextual correction. This dynamic is not uniform: the greater the deviation between the modified information and the model's prior, the less likely the model is to prefer the retrieved content. This tug-of-war establishes that RAG does not guarantee a correction of the model's response but rather creates a conflict---which leads to the epistemic misalignment and brittleness that become even more pronounced when knowledge sources are static or temporally outdated, as discussed in the following section.

This phenomenon is further confounded by findings from studies on realistic knowledge conflicts. Previous research that presented models with synthetic, misleading documents to contradict correct parametric answers often overstated the failure rate of knowledge updates \cite{ref9}. When knowledge conflicts are constructed in a realistic setting with genuinely contradictory documents, knowledge updates fail less often than previously reported. However, the studies reveal a ``parametric bias'': when the incorrect parametric answer is present in the retrieved context, the model is much more likely to fail to update its knowledge. Rather than independently evaluating the evidence in the retrieved documents, the model's internal (possibly outdated) knowledge can dominate its reading comprehension and behavior. This insight is crucial because it suggests that simply appending retrieved documents to the prompt is insufficient for reliable factual grounding in ambiguous or contradictory cases---a limitation that becomes highly consequential when the contradiction arises from the rapid evolution of real-world events, thereby directly connecting to knowledge-updating challenges discussed in the subsequent subsection.

This dynamic leads to significant overconfidence issues, exacerbated by the imperfections of the retrieval process itself. Conformal risk analysis frameworks have explored providing theoretical guarantees for generation risks, but their findings also show that RAG's benefits are conditioned on the quality of the retrieval model and the transformer \cite{ref135}. There is no unconditional guarantee that a RAG system has a lower generation risk than a vanilla LLM. Therefore, when retrieval fails to identify relevant documents or retrieves contradictory information, the RAG framework does not fail gracefully. Instead, it can project uncertainty in a misleading way. Theoretical analysis demonstrates that RAG can lower its risk if retrieval is reasonable, but the assumptions behind such guarantees often break down under the noisy, messy conditions of real-world data---precisely the conditions where stale and long-tail knowledge sources begin to exert their pull, leading to non-transparent overconfidence.

This overconfidence manifests particularly strongly in the presence of flawed, top-ranked results. The standard RAG paradigm treats every retrieved chunk as equally valuable, but the concept of ``credibility-aware generation'' highlights that not all contexts are equal \cite{ref239}. Existing frameworks ignore the credibility of sources, and inevitably suffer from the flawed information introduced during retrieval. When the model is trained only on retrieved data, it lacks the ability to weigh the credibility of different snippets, leading to a tail-wagging-dog effect where a low-quality retrieved passage can undermine an otherwise strong parametric response. The success of the generation hinges on the model's ability to judge the utility and factual quality of the given evidence, a capability that is often absent---a critical flaw when relying on unstable, community-maintained knowledge sources, as described in the previous subsection's analysis of curation sustainability.

Despite these inherent fragilities, the system is often designed to enforce confidence, without the necessary safeguards. Approaches to detect hallucination, such as internal state analysis and uncertainty estimation, have been proposed \cite{ref134}. Techniques such as the EigenScore metric exploit the eigenvalues of the covariance matrix of the response to measure semantic consistency in dense embedding space, alongside test-time feature clipping to truncate extreme activations. While these methods can help identify overconfident hallucinations, integrating such uncertainty detection into a deployed RAG pipeline remains a challenge. RAG systems produce one output without communicating a degree of confidence or the confidence level of the retrieved grounding, misrepresenting the underlying reliability of the output. Moreover, a system that cannot communicate its own uncertainty is incapable of flagging when the knowledge base is stale, sparse, or poisoned---blinding the system and user to the type of the attacks detailed in the following section.

The situation becomes more alarming when we evaluate the robustness and the propagation of misinformation. Studies assessing the robustness of RAG systems against errors in external knowledge across multiple typologically diverse languages expose substantial variability among models: many exhibit over 88\% hallucination rate on non-relevant passages, while others demonstrate up to 74.9\% error rate on relevant passages, underscoring that most models struggle to balance the necessary capacities for robust grounding \cite{ref125}. Furthermore, benchmarks that verify whether generated text is actually supported by retrieved documents---such as FACTAL entailment approaches---indicate that conventional textual entailment is insufficient for spotting subtle factual errors in LLM-generated content \cite{ref146}. This indicates that RAG's foundation---the verification of factual entailment between the output and the retrieved docs---is itself fragile. Moreover, models' varying sensitivities to the degree of their ``pull'' from internal and external sources. The tendency to refuse or ignore conflicting information depending on their style (dependent vs.~intuitive) can make RAG notoriously rigid in misinformation scenarios \cite{ref248}. Unverified or erroneous content can become a ``ground truth'' in future retrieval cycles, a mechanism that poisons static and dynamic knowledge bases and serves as a crucial link to the feedback loops addressed in the next section.

Prompt-based conflicting patterns are not just a theoretical concern; they are observed during empirical evaluations. For example, when the top-1 retrieved document contains a biased red herring and a top-2 contains the correct answer, the model may fail to reconcile these sources, further propagating the error. While some efforts have been made to evaluate an LLM's ability to judge the utility of retrieved passages---as opposed to their relevance---these attempts reveal persistent limitations requiring a more systematic approach to utility evaluation \cite{ref53}. However, without explicitly modeling source credibility, utility, or temporal validity in retrieval, RAG systems become unable to handle the length-tail entities and rapidly updated knowledge sources---key shortcomings of over-reliance on static KBs---and consequently remain vulnerable to error propagation from suboptimal retrieval results.

The core of the dilemma is this: effective RAG requires not only access to retrieved evidence but also the model's ability to critically assess and integrate this evidence with its parametric knowledge. Yet in practice, the model often lacks a mechanism to verify and detect when its output is based on unverified or irrelevant premises. Frameworks like ``Retrieve Only When Needed'' \cite{ref150} emphasize that hallucination mitigation is even more difficult than it appears: models tend to generate external hallucinations when they over-rely on irrelevant evidence---and overconfidence is amplified when the model does not have the introspection to selectively retrieve information at all. The ``tug-of-war'' literature and subsequent probes demonstrate clearly that the issue is not binary. Rather, RAG systems face a continuum of trust: from conflicting memories to static failures. The solution arguably lies not in a single improvement, but in a broader engineering of the process and architecture, incorporating capabilities such as hallucination-aware decoding, upper-confidence calibrated generation, and source-aware attention mechanisms that actively verify sources \cite{ref47}. Only when there is a calibrated awareness of the limits of both retrieval and parametric knowledge can we avert the persistent misinformation propagation, which can otherwise be reinforced and propagated by the very system designed to reduce it.

In summary, the trustworthiness dilemma is as follows. An ideal RAG system must be confident, grounded, and robust. In current practice, however, accurate retrieval can help a system improve, but retrieval failures can quietly transform a plausible output into a confidently stated false one. When such misinformation is retrieved by future queries, it risks creating self-reinforcing feedback loops---intricately tied to the systemic fragility of knowledge bases described in the next subsection. And when an attacker can manipulate this process by injecting corrupted references, as the following section details, the trustworthiness problem transitions from a question of accuracy to a security threat. As a critical need, research in calibrated generation and the role of internal states \cite{ref47} must move from considering model errors to towards refined methods that preserve the integrity of generated content in the face of adversarial retrieval disturbances, while explicitly acknowledging that the danger of an uncritical system is not only that it fails, but that it fails confidently, in ways that are harder to detect and thus more dangerous.

\subsection{Security and Privacy Threats: Prompt Perturbation, Poisoning, and Data Leakage}

The integration of external knowledge bases through Retrieval-Augmented Generation (RAG) introduces not only enhancements in factual grounding but also a significantly expanded attack surface that exposes LLMs to a new class of security and privacy threats. Unlike purely parametric models, RAG systems intertwine the generative capabilities of LLMs with dynamic, often unvetted, external datastores, creating vulnerabilities at the intersection of prompt engineering, retrieval mechanisms, and knowledge storage. These threats can be broadly categorized into prompt perturbation (including leakage), data poisoning, and explicit extraction attacks, each requiring a nuanced understanding of the system's architecture to adequately defend against. Critically, these security vulnerabilities do not exist in isolation---they compound the trustworthiness challenges discussed earlier. A system that cannot reliably distinguish credible from misleading retrieved evidence is precisely the system most susceptible to adversarial manipulation of that evidence. Moreover, the lack of calibrated confidence and uncertainty communication noted in the previous subsection means that successful attacks often produce outputs that appear just as confident and well-grounded as benign ones, making silent compromise even more insidious. Therefore, addressing these security vectors is not a separate endeavor but a deepening of the same critical assessment that RAG systems must conduct on their own retrieved knowledge.

A particularly insidious vulnerability arises from \textbf{prompt leakage}, a privacy threat where malicious users manipulate the model into revealing the underlying system prompts or contextual knowledge, which often contains proprietary instructions or sensitive reference data. Research has shown that this vulnerability is significantly amplified in multi-turn interactions, where attackers can leverage the model's inherent tendency to agree with the user, progressively coaxing it into leaking its directive. This threat model has demonstrated a startling efficacy, elevating the average attack success rate (ASR) to 86.2\% across various domains and LLMs, with GPT-4 and claude-1.3 exhibiting a 99\% leakage rate \cite{ref159}. The danger extends beyond simple instruction disclosure to the leakage of contextual knowledge used in the RAG pipeline, and notably, defense is non-trivial; even a proposed combination of multi-tier black-box defenses, including query rewriting, cannot completely mitigate the risk, retaining an ASR of over 5\% \cite{ref159}.

Beyond the system prompt, the datastore itself is a prime target for extraction attacks. RAG systems are often built with proprietary, private data to enable customized adaptations, which raises a critical question: can an adversary retrieve entire records verbatim? Research shows that an attacker can exploit the powerful instruction-following capabilities of modern LLMs to perform severe data extraction from the retrieval datastore \cite{ref222}. The success rate of this attack scales with model size, and it has achieved a 100\% success rate on customized production GPT instances, demonstrating that extracting sensitive text from the integrated knowledge base is not only possible but practical to execute. This is consistent with the previously discussed problem of overconfidence: a RAG pipeline that treats all retrieved chunks as equally valid is particularly vulnerable to extraction attacks because the model lacks the mechanism to distinguish genuinely relevant sensitive data from a maliciously crafted request.

The security and integrity of the retrieval database are further compromised by specific poisoning attacks designed to manipulate its contents. It is not only confidentiality that is at risk; the integrity of the knowledge database is a foremost security concern, as RAG systems are vulnerable to a particularly stealthy form of manipulation known as knowledge poisoning. In this attack strategy, the retrieval database is corrupted with minimal number of carefully crafted malicious texts, and studies show that injecting merely 5 poisoned texts into a database containing millions of entries can achieve a 90\% attack success rate \cite{ref249}, highlighting that even nearly pristine datastores are not safe. This attack treats the RAG process as a whole---constructing poisoned texts that maximize the likelihood of the target output becoming the model's answer via an optimization step considering both the retriever's embedding space and the generation. It is implementable in both black-box and white-box settings, underscoring that knowledge poisoning can be executed with minimal resources. The fact that a small injection can so dramatically alter the behavior of a RAG system also serves to emphasize the hidden fragilities of the pipeline: the system has little ability to detect when it has been provided poor or malicious evidence, and the lack of verification mechanisms makes the attack both easy and dangerous.

While the previous attack seeks to corrupt the database, the integration and life-cycle of the RAG brings additional attack vectors. Data poisoning can also manipulate the model's behavior during update phases, where attackers replace features during instruction tuning to insert backdoor triggers \cite{ref250}. This shows that an adversary can degrade a target model's performance with few resources, poisoning as little as 1\% of the training set. Relatedly, model extraction attacks allow adversaries to steal the architecture and parameters of the component models itself, enabling the model to be replicated or exploited \cite{ref251}. These vulnerabilities in the broader AI lifecycle must also be considered, since a large infrastructure supporting RAG is made up of heavily trained components. It is, therefore, necessary to focus the analysis on real-world attack goals and potential actual harm, rather than on the imperceptibility of the perturbations on their own \cite{ref252}.

Adversarial inputs can also be systematically scaled across topics \cite{ref253}, demonstrating how sophisticated attackers can manipulate domain-specific knowledge and reveal the brittleness of RAG grounding beyond general-purpose data. This scaling effect is particularly relevant to the dynamic reasoning discussion in the following subsection---the need to assess the utility of the retrieved context is elevated to a security imperative, not only a performance one, because a poisoned context that is trusted will directly propagate the attack's intended false content. Interactive and generation systems increasingly need robustness to adversarial neighborhoods---setting up data augmentation for defense. For the architecture, it is important to ensure that interactions in a distributed setting are safe and maintain, mitigating against adversarial neighborhood attacks \cite{ref254}.

Additionally, and between security and the previously discussed model updating, the RAG does not provide a shield from training data leaks. Novel attack showed that RAG's mechanism of augmenting the LLMs with retrieval is a double-edged sword: it introduces new privacy leaks while potentially mitigating training data leakage from the original model \cite{ref221}. This demonstrates a critical paradox: while RAG may protect the model's parametric memory, the non-parametric retrieval component emerges as a highly susceptible differential avenue directly siphoning sensitive knowledge. Together, these vectors require that security evaluations consider the entire RAG pipeline---including the `hard' facts and source retrieval medium---as a system: the pertinence hind in the retrieved data is, as previously highlighted, not a problem that can be fine-grained and solved simply by optimizing the retriever alone.

When asked about defense, it is essential to recognize that the current landscape is additionally unsolid. Defenses are currently battling a multi-stage problem that spans across the data store, the model access interface and the prompt. High-level mitigation strategies---such as input filtering or query rewriting---have shown some defense in black-box settings, but architecture-level or multi-tier being susceptible \cite{ref159}. Because RAG's success depends on the interaction of the different components, from the security field
redundant---we need to think of defense in terms of ensuring secure interactions with constant monitoring and anomaly detection for usage, embedded indexing and retrieval quietly, and hiding instructions.

In answering what a complete RAG defense blueprint looks like, we need to consider harm on the retrieval file, spending augmented data with isolated walls and presumable with measure. On the defense, as adversarial methods interplay with the defensive updating loop, the security and the defensive method are two closed-loop outputs. We must design pipelines robust to adversarial perturbation across \emph{and} the retrieval ; detection unit, like deployment-time monitoring for prompt leakage from a black box might be needed. Following trust from previous section, metacognitive reasoning seems an obvious avenue for a security-savvy RAG, where the system is trained to recognize and flag internal state inconsistencies between prompt, retrieved data, and prior knowledge---arising from malicious interactions or tampered databases. In sum, defending RAG is not an independent extra layer of `guardrails' but design a general, credible integration for know-what-they-know, where security and trust remain \emph{analytically unified} in the generative response itself.

\subsection{The Goal of Dynamic and Adaptive Reasoning: From Static Retrieval to Self-Aware Generation}

The evolution of retrieval-augmented generation (RAG) has progressively shifted from static, one-size-fits-all pipelines toward dynamic, context-aware reasoning systems. This transition is not merely an engineering preference but a direct response to the security and integrity vulnerabilities identified in the preceding subsection. If a system indiscriminately retrieves and incorporates a fixed number of passages---regardless of whether retrieval is necessary or whether the datastore has been compromised---it not only diminishes the model's versatility and produces unhelpful responses \cite{ref109}, but it also amplifies the attack surface by giving equal weight to poisoned or maliciously injected content. The central open problem in this domain is therefore no longer \emph{how} to retrieve accurately, but \emph{when} to retrieve, \emph{how much} to retrieve, and \emph{how} to \emph{integrate} retrieved knowledge in a manner that respects both the task's demands and the model's ability to detect unreliable evidence. This subsection explores these open problems, highlighting pivotal frameworks such as self-reflection, active retrieval, and the hidden gaps that influence robust dynamic behavior---capabilities that are as essential for defending against poisoned knowledge bases as they are for improving benchmark performance.

\textbf{The Question of When to Retrieve: Moving Beyond Fixed Pipelines.} Traditional RAG systems treat retrieval as an unconditional step that occurs before generation. This static approach often introduces noise: irrelevant passages can impair the model's performance and lead to factual degradation \cite{ref5,ref109}. Critically, in a security context, this unconditional trust in the retrieval pipeline means that the model will ingest whatever the datastore returns---including the crafted malicious texts described in the prior subsection, which achieve a 90\% attack success rate with just five poisoned entries. The central technical challenge is thus determining a query's genuine dependence on external knowledge, both for efficiency and for security. Some questions can be answered correctly from the model's parametric memory alone, and introducing retrieval in these cases can confuse the generation process, introduce conflicting evidence, or activate a hidden trigger. Research into self-knowledge guided retrieval has revealed that LLMs possess latent ``self-knowledge''---an ability to recognize what they don't know---and that they can reference previously encountered similar questions to adaptively call for external resources \cite{ref255}. Leveraging this capability, systems can avoid unnecessary retrieval, reducing the attack surface by limiting the number of times the model touches a potentially poisoned knowledge store.

\textbf{Self-Reflection and Introspective Reasoning Mechanisms.} A crucial step toward robust dynamic behavior is the integration of self-reflection capabilities. Self-RAG is a pioneering framework that reframes RAG as a reinforcement learning problem where a model is trained to retrieve passages on-demand, generate responses, and critically evaluate its own generations through special ``reflection tokens'' \cite{ref109}. This approach allows the model to adaptively decide when retrieval is necessary and to critique its own confidence against retrieved passages, moving beyond simple generation to a form of self-critique. Yet, the security analysis in the previous subsection reveals a critical boundary: self-reflection can reject noisy contexts, but it cannot reject a skillfully poisoned text that is unknowingly well-aligned with the attacker's goal. Therefore, while self-reflection is necessary, it is insufficient for handling inputs that are structurally indistinguishable from benign text. When internal knowledge is insufficient or incomplete, frameworks such as RA-ISF introduce an iterative process of self-feedback that decomposes tasks into three submodules---retrieval, feedback, and generation---creating a loop that allows for continuous refinement of the answer \cite{ref5}. This iterative loop provides renewed opportunities to detect and discard an anomalous retrieval that contradicts the coherence of the accumulating answer.

\textbf{Active, Iterative and Error-Driven Retrieval.} Dynamic reasoning demands more than a single retrieval pass when handling multi-hop or knowledge-intensive tasks. The ``retrieve-once-then-read'' paradigm struggles with queries that require intermediate reasoning steps, and it also unnecessarily deepens the attack surface by committing early to a single leg of potentially flawed evidence. To address this, active retrieval mechanisms that revisit the knowledge source based on the progress of generation have been developed. Frameworks like ActiveRAG shift the role of the LLM from a passive receptor to an active learner, employing a ``Knowledge Construction'' mechanism to link newly retrieved knowledge with previously acquired knowledge, followed by a ``Cognitive Nexus'' mechanism that integrates chain-of-thought patterns with constructed knowledge, calibrating the model's intrinsic cognition \cite{ref177}. This iterative re-querying allows the model to attempt the evidence path and verify the intermediate results before integrating them into the final composite answer.

\textbf{Metacognitive Regulation.} Beyond active learning, the integration of metacognition---the ability to monitor, evaluate, and plan---offers a new frontier in dynamic reasoning. MetaRAG presents an approach that combines RAG with metacognitive regulation, allowing models to introspect on their own cognitive processes through a three-step pipeline (monitoring, evaluation, and planning) \cite{ref44}. This is especially salient following the discussion of security defenses: a metacognitive process that can identify inadequacies in its initial cognitive response is precisely the mechanism needed to flag the internal state inconsistencies that arise from malicious interactions or tampered databases. This shift from a purely reactive system to a proactive one---which identified inadequacies \emph{before} they locate final output---is the unifying first step toward the security-aware architecture recommended in prior passages.

\textbf{Responding to Retrieval Failures: When the Context is Wrong.} A truly adaptive system must also maintain its integrity when retrieval fails or when the retrieved content contradicts the model's prior knowledge. The phenomenon of knowledge ``conflicts'' lies at the core of the dynamic problem in RAG. Research reveals that knowledge updates often fail less frequently than expected, yet a clear ``parametric bias'' exists: when an incorrect parametric answer appears in context, updates become more likely to fail \cite{ref9}. This reveals a fragile model behavior: they tend to rely on their own parametric knowledge over the provided textual premises. Empirically, the ``tug-of-war'' between prior and retrieved context demonstrates that while correct retrieval fixes most errors, when reference documents are perturbed with increasingly incorrect values, the LLM will recite the incorrect \emph{modified} information more readily when its internal prior is weak, but resists when the prior is strong \cite{ref238}. This finding serves as an important contact point with security concerns: the detection of poisoned texts fundamentally relies on the model's ability to surprise a confident, contradictory prior. A self-aware generation system must know the strength of its own prior; when the prior is strong, the system should be more resistant to misleading noise, whereas when the prior is weak, it should defer more heavily but also place the retrieved texts under stricter provenance scrutiny---preventing a single compromised store from achieving its intended deception.

\textbf{Ignoring and Being Aware.} A complete feedback system is not yet complete without robustly ignoring harmful context. The stream of work on understanding when retrieval is not helpful positions the model's preferences into an ``intuitive'' or ``dependent'' style, where dependent models defer to context while intuitive models rely exclusively on internal memory \cite{ref248}. Role-play interventions can change the style, but with varying degrees of success and different upper-bound between models. This reveals that overconfident or wrongly grounded behavior in retrieval is not purely a function of good embeddings: it reflects the model's personality, which can be manipulated dynamically (e.g., adopt a ``dependent'' style for highly specific queries, but with an ``intuitive'' fallback clause when the retrieval content smells suspicious and contradicts strong prior knowledge). These insights point highlight the necessity of dynamic behavior calibrated according to the quality of the retrieval context as well as the model's confidence in its internal memory---an integration that strengthens the integrity gaps identified in the security discussion.

\textbf{Hierarchical, Graph-Structured, and Multi-Hop Reasoning.} From dynamic reasoning beyond linear integration, real-world queries require multi-step planning and knowledge integration across multiple documents. Multi-vector and multi-hop retrieval methods break complex queries into manageable sub-queries. Building on this, frameworks use multiple thought nodes in a hierarchical graph to answer complex questions (HGOT), applying a divide-and-conquer strategy and refining self-consistency majority voting weighted by citation quality of thought paths \cite{ref48}. This demonstrates that only reveal answers are equally well-supported by their retrieval passages; the system weights individual answer candidates dynamically by the quality of the citations they rely on. Extending this reasoning to the plotting in the next subsection: when these graph edges are sourced from the model's own untampered extraction, they become valuable; but if the knowledge graph has a poisoned source (for example, the attack discussed earlier), the entire edge weight is fraudulent. This requires the addition of the dynamic reasoning capabilities we propose here---specifically, embedding a verification step at every junction from the source graph to meet the attribution standards that become central in the following subsection.

\textbf{The Missing Unified Framework and Future Directions.} Despite these advances, a fully self-aware RAG framework that seamlessly integrates query complexity assessment, self-knowledge calibrated confidence, adaptive retrieval, and iterative refinement accounting for uncertainty about the model's own internal knowledge remains an open challenge. Future systems must incorporate uncertainty quantification not just at the final answer stage but at every intermediate stage, especially the retrieval connection. We recommend a self-aware loop: a system that knows when it has achieved a correct answer and knows to stop; knows when the retrieval path has hit a conflict and re-routes; and ultimately, as the security discussion above implies, knows when a retrieval result is internally inconsistent with its foundational knowledge. More work is needed to weave these components together, from independent mechanisms such as conformal prediction-based retrieval confidence \cite{ref47} to quantify retrieval confidence with guarantee against unanswerable chunks, into the dynamic loop. The future of RAG is not merely as a retriever, but as a self-aware generation agent that orchestrates retrieval and self-monitoring of its own inferential confidence---simultaneously task-oriented and security-aware, protected against the adversarial dynamics that undermine the retrieval process, and designed to move from the promise of dynamic retrieval, to the stricter integrity and attribution requirements of modern knowledge grounding.

\subsection{The Risky Integration of Attribution and New Frontiers of Concern}

As Retrieval-Augmented Generation (RAG) systems move from the controlled, accuracy-driven benchmarks discussed in the previous subsection toward the high-stakes deployment realities examined in the following subsection, a new class of risks emerges at the point where retrieved knowledge is actually integrated. Where earlier concerns centered on dynamic retrieval decisions---when to retrieve, how much to trust, and how to self-reflect---the most pressing frontier has shifted to attribution: the capacity to trace an output back to precisely which retrieved source (or combination of sources) produced it. The naive assumption that retrieved passages are definitively correct and inherently attributable to a reliable source collapses in real-world environments, where attestation is not a post-hoc labeling task but is interwoven with the reasoning process itself. When attribution fails, the system does not merely lose explainability; it loses the very conditions that make human oversight possible, and it can propagate errors that overwhelm the safeguards described in the following subsection.

\textbf{From Dynamic Retrieval to an Attribution Gap.} The previous subsection showed that dynamic reasoning requires the system to decide \emph{when} and \emph{what} to retrieve. In multi-hop reasoning, this challenge transforms into multiple retrieval events whose results must be synthesized. The model incrementally composes a line of argument across documents, in each step creating an ``attribution gap'': a final output that reflects a derivative of knowledge originating from several sources, not fully containing each source. The risk is compounding errors---if one intermediate document is misrepresented or contains a subtle contradiction, the final answer is confidently generated but fundamentally flawed, and there is no traceability to the root cause \cite{ref256}. Perhaps most problematic: the self-aware retrieval and iteration mechanisms described earlier have no way to know that these compounding errors have occurred, because valuable ``confidence'' signals are applied to the answer, not to the attribution chain. The system knows it is confident; the attribution system does not know \emph{why}.

\textbf{Multimodal and Cross-Channel Attribution.} This attribution gap becomes critical in the multimodal contexts that advanced RAG is increasingly entering. A multi-modal RAG system does not retrieve only text; it integrates visual information, LIDAR/swarm data, sensor outputs, and textual corpora. Consider a model that retrieves a legal rule (``yield to pedestrians at all times'') and fuses it with real-time vision data capturing a jaywalker. The resulting driving decision is a synthesis of these. If that decision is unsafe, the question ``was this caused by a wrong retrieval or a wrong vision analysis?'' is not answerable. In frameworks such as LMDrive and RAG-Driver (which authors demonstrate positive potential for language and sensor fusion \cite{ref257,ref206}), attribution remains a rhetorical rather than a technical question. Because the failure cannot be pinned to a particular channel, the feedback loop that we must use to react to that failure is broken. The system cannot even perform a useful self-reflection, because it does not know what to reflect on. The only realistic path to rehabilitation of these systems is a ``token-based tracing'' layer that records exactly which input (sensory or textual) was in context at every stage of a generated chain.

\textbf{Global Knowledge-Graph Risk: an Attribution Chain that is Lost.} A promising response to multi-hop retrieval of heterogeneous relations is to represent global knowledge through knowledge graphs, which create multi-relational links between sources. However, integrating AGraph into RAG introduces an attribution risk of its own. When the model traverses a knowledge graph, it is not using the original text; it is using a series of relational edges that have been extracted from that text. The provenance of the extraction---the original publication, the paragraph, even the human author who extracted it---is lost. If the source document contained a flawed triple relation, the downstream model receives a statement identical to truth, with no visible path to the flawed source. This is the ``missing attribution chain'' of the global knowledge. While frameworks like RAG-Driver attempt to address a piece of this gap, ``\cite{ref206}'', the field lacks metadata conventions to manage this graph-level provenance. Integrating a graph without its meta-layer is functionally equivalent to having a retrieval component that hides the source of its own source---arguably worse than no graph at all.

\textbf{Scientific Reasoning and Autonomous Driving as Open Safety Frontiers.} The concealed attribution failure is not a mere ``research gap''; it is a tripping hazard for highly visible application domains. With scientific discovery as mention dev, LLM-agents can now draw multiple retrieved resources to propose experiments. But if the supporting knowledge graph contains a corrupted or partially hallucinated edge, the agent can propose a hazardous protocol while still quoting a plausible paper. \cite{ref258} \cite{ref258} demonstrates that the absence of a causal attribution trace transforms a research assistant from a tool into a mode of dangerous autonomy that is inherently in early, run-away optimizing a loop of harmful actions unless we embed ``attribution-at-checkpoints'' and the authority to stop. The inability to attribute the mistake to the correct edge also completely compromises the ``self-reflection'' we described in the previous subsection---the system cannot ``know when it is wrong'' if it cannot track what evidence it used and why.

The same holds for autonomous driving agents. The previous subsection on the development of knowledge-driven autonomous driving points out that the system must rely on ``world knowledge''-expected vehicle behavior for location types or time of day. This system could retrieve a behavioral assumption, then integrate it with a particular vehicle's on-board sensor reading (such as a leading car and its visual cue). If the final trajectory is rejected by a safety monitor, a human reviewer cannot clearly tell whether the unsafe maneuver was caused by the retrieved ``world knowledge'' (actively) or by the incorrect vision det** that made them work. However enormous advancements in Language-MPU \cite{ref259} show the correctness of the integration, the field is equally only as good as its capacity to attribute \textbf{decision-making actions} to either internal or external anchors, and to that end we have not yet created standards for proving that a retrieval change in context is isolated from a physical environmental change. Without the ``link'' to the retrieved semantic kernel, we cannot distinguish - a model that ``misunderstood the language data'' from the traffic scenario is the reason.

\textbf{Feedback Loops and Long-Tail Attribution Decomposition .} The attribution problem becomes recursive when systems begin to close feedback loops. A RAG system generates content based on a retrieved source---then that output becomes part of a second set. A future generation might retrieve the first generation's output as an authoritative representation of the first source. If the first generation contained misinterpretation of a long-tail event, that error becomes a new artifact, ready to guide future retrieval loops that construct new chains of logic. This problem is more than the ``echo chamber'' of engineering; it is a \textbf{de facto} design mechanism for malice or malfunction. Instead of simply marking attention as a stop token, a highly dynamic RAG must dynamically annotate the source graph of the original iteration and retain those annotations during downstream retrieval. The question, ``How do we attribute the attribution itself'' has no answer unless the system stores provenance alongside the response generation. In the case of high-stakes systems, we propose a ``dynamic, dry rot-filtering'' that \emph{tracks} the lifecycle of the content: if a generated artifact is re-retrieved and re-integrated into a /margin/store, the RAG framework must be configured to verify provenance. orphans.

\textbf{The Road Ahead for Safe Integration}: Future RAG design is not solely about retrieval. It is about \textbf{defending the attribution path} from errors and mistakes. For Multi-Hop approaches, new network architectures that implement token-level attention combined with passage-level evidence must be used to \emph{trace} the path from a list of sources to a proposition; this would allow Eloquently the ``chain of reasoning'' to be recreated by a ground-truth reader. Second, high-quality global knowledge graph without metadata layers must be considered as out-of-context (having an empty provenance is not a bug, but a security vulnerability); the metadata tracking layer is a prerequisite. Finally, our review of the high-stakes domain indicates we need ``rules of reflection'' that inevitably couple attribution logic with the self-reflection and the retrieval act: a RAG should be able to stop asking ``what did the model believe'' and ask ``what did the model \emph{retrieve}? And why did it retrieve it?'' These systematic changes are not experiments. Without transforming the attribution function from a nice-to-have to a core function, architectural reasoning in AI will remain blind---not just legally and philosophically---in the environment we reinforcing feedback loops with unsourced trouble, leading to the opposite of the governance-friendly safety systems we describe below.

\subsection{The Unsafe Practices: Human-in-the-Loop, Benchmarking, and Legal Accountability in Real-World Deployments}

The deployment of Retrieval-Augmented Generation (RAG) systems in real-world, high-stakes environments introduces a constellation of unsafe practices and governance challenges that extend beyond the purely technical limitations discussed in the preceding subsection. While RAG is frequently lauded for mitigating hallucination and grounding responses in verifiable sources, its integration into operational pipelines brings to the forefront issues related to misinformation, legal accountability, and the effectiveness---or ineffectiveness---of mandated human oversight. Critically, these risks do not merely coexist with the attribution failures described above; they are often directly exacerbated by them. When a system cannot trace the provenance of its own output, the human and institutional checks designed to catch errors are cut off from the very information they require to function. This subsection decomposes these socio-technical risks, critically examining the paradoxes of human-in-the-loop mechanisms, the fragility of evaluation benchmarks, and the profound accountability gaps that arise when RAG systems interact with dynamic information ecosystems.

A primary concern is the amplification of misinformation through RAG pipelines. By retrieving and synthesizing content from the web or internal knowledge bases, RAG systems can inadvertently generate fluent, authoritative-sounding responses that propagate false or misleading claims, particularly when the source documents themselves are unreliable. The challenge is compounded by the fact that misinformation is rarely a simple binary of true or false; it often contains partially correct but misleading elements that undermine public trust \cite{ref260}. This ``factual but misleading'' nature is uniquely dangerous in a RAG system, as the presence of a cited source can confer an undue sense of legitimacy. This risk is directly aggravated by the ``missing attribution chain'' of knowledge graphs discussed earlier: when a retrieval component hides the provenance of its own source, downstream systems cannot detect whether the retrieved statement was a faithful rendering of a primary source or a corrupted edge extracted by an upstream process. Without upstream filtering and provenance-aware credibility assessment, RAG risks becoming an automated engine for recycling and re-asserting harmful narratives, verbatim, with a simulation of authority.

The human-in-the-loop paradigm is frequently proposed as a universal safeguard, yet it is fraught with operational fallacies and misaligned incentives---and in some cases, it constitutes its own form of unsafe practice. Audits of algorithmic governance policies reveal a problematic assumption: that imperfect, time-constrained humans are actually capable of performing the oversight functions expected of them at algorithmic scale and speeds \cite{ref261}. In the context of RAG, effective oversight depends on a human reviewer having the ``causal power,'' ``epistemic access,'' and cognitive agency to audit the system's conclusions \cite{ref262}. However, exactly when the attribution architecture described earlier is broken, the overseer stands no chance of exercising that access---if the system cannot trace its own confident output to a source or a reasoning chain, a human reviewer cannot either, no matter how much time or will is available. Indeed, the massive volume of RAG outputs and the cognitive load required to fact-check them in high-velocity environments---such as a public dashboard generating COVID-19 treatment claims---often overwhelms human capacity \cite{ref263}. This creates a dangerously seductive dynamic: human oversight becomes a symbol of accountability---a sign that error is being caught---but the human reviewer is effectively incapable of performing the verification. In high-stakes public sector contexts, structural difficulty is amplified: decisions are embedded in nested teams of experts, but AI systems are integrated without careful attention, disrupting the ``team-in-the-loop'' dynamics that sustain professional reasoning and evaluation, ultimately eroding the quality of the decisions \cite{ref264}.

Benchmarking and evaluation represent another set of ``safe practices'' that fail in the field. Benchmark data itself is often constructed with a filter that smooths over the ragged edges of reality; as a result, a system that excels in a controlled benchmark can still be catastrophically wrong on the long tail of naturally occurring inputs, precisely because the benchmark excludes the real-world distributions that a deployment will face. Worse, multiple models that each perform well on the same benchmark can produce conflicting predictions on the same inputs \cite{ref265}---the prelude to the ``black box with a new surface'' of RAG. The design of evaluation sets can also embed normative and demographic biases \cite{ref266}. Consequently, a system that earns high marks in a benchmark may fail in human-centered contexts and harm those are not considered during its design.

The governance and feedback dynamics of RAG deployment add another layer of concern. RAG outputs that perform well in benchmarks continue to interact with the environment in a recursive loop. When partially hallucinated content is scraped by search engines or inserted into synthetic datasets, it enters the training corpora of future generations. This spawns a closed cycle: a model's own hallucinations---or the retrieved edge of an unattributed knowledge graph---become the ``verified facts'' by subsequent retrievers, creating an ``echo that never stops.'' The attribution breakage described in the previous subsection is the ultimate enabler of this runaway loop: if a chain of provenance is never retained, there is no way to quarantine or ``dry-rot filter'' the erroneous artifact, exactly as we highlighted in the context of recursive attribution. Cases of algorithmic drift demonstrate the added danger of such self-reinforcing, recursive biases; with the model influencing the outer environment, even a human ``in the digital loop'' will not supply the required signal to stop it.

Underneath all these operational risks lies the profound legal and methodological vacuum. RAG's language models are not formally deterministic and cannot be subjected to the kinds of formal counterfactual verification that for decades anchored algorithmic accountability \cite{ref267}. They operate like stochastic parrots rather than logical computers. The inability to interrogate an agent on its behavior in a specific sequence of events---like a car crash---or a decision flowchart showcases a systemic lace of legal accountability in the deployment lifecycle. At the same time, the structural inequities of the real world cannot be abstracted away or fixed: accountability is socially constructed, and in contexts where users are financially stressed, they may ``blame themselves'' for the harms they experience from the interface, foreclosing even the possibility of feedback \cite{ref268}. When feedback loops are broken in this way, the system flips to the dynamic ``run-away optimizing a loop of harmful actions'' foreshadowed in the scientific-use risks of the previous subsection.

The path forward must therefore abandon the illusion that human review is a sufficient safeguard by itself. While legal frameworks such as the EU AI Act mandate oversight, the reality is that without humanly accessible causal attribution---described in the previous subsection: token-level traces and metadata layers with authority to stop---human oversight becomes a legal ritual rather than an operational control. On the contrary, policy must be designed with the assumption that not the technology alone, but the users---fact-checkers, decision-makers, community reviewers---are functionally capable in a small-scale setting of exercising real judgment. The future of RAG does not depend solely on stronger retrievers or more powerful a generator, but on rebuilding the \emph{social} context in which those algorithms are deployed: redefining the relationship between system, knowledge, and environment, and designing for secure systems where attribution is integrated to enable agency for whoever can---and must---say ``no.'' Without such a foundation, RAG's speed and scale become an unresolved fire hazard for the very variables and communities it claims to improve---setting up expect failures not at the retrieval component, but in our ability to see, intervene, and correct the procedural unfairness.


\begin{thebibliography}{268}
\addcontentsline{toc}{section}{References}
\bibitem{ref1} Retrieval-Augmented Generation for Large Language Models A Survey. arXiv:2312.10997v5, 2023. \url{https://arxiv.org/abs/2312.10997v5}
\bibitem{ref2} A Survey on Retrieval-Augmented Text Generation for Large Language Models. arXiv:2404.10981v1, 2024. \url{https://arxiv.org/abs/2404.10981v1}
\bibitem{ref3} Improving Retrieval for RAG based Question Answering Models on Financial Documents. arXiv:2404.07221v1, 2024. \url{https://arxiv.org/abs/2404.07221v1}
\bibitem{ref4} ARAGOG Advanced RAG Output Grading. arXiv:2404.01037v1, 2024. \url{https://arxiv.org/abs/2404.01037v1}
\bibitem{ref5} RA-ISF Learning to Answer and Understand from Retrieval Augmentation via Iterative Self-Feedback. arXiv:2403.06840v1, 2024. \url{https://arxiv.org/abs/2403.06840v1}
\bibitem{ref6} Fine-tune the Entire RAG Architecture (including DPR retriever) for Question-Answering. arXiv:2106.11517v1, 2021. \url{https://arxiv.org/abs/2106.11517v1}
\bibitem{ref7} RAGGED Towards Informed Design of Retrieval Augmented Generation Systems. arXiv:2403.09040v1, 2024. \url{https://arxiv.org/abs/2403.09040v1}
\bibitem{ref8} MultiHop-RAG Benchmarking Retrieval-Augmented Generation for Multi-Hop Queries. arXiv:2401.15391v1, 2024. \url{https://arxiv.org/abs/2401.15391v1}
\bibitem{ref9} Studying Large Language Model Behaviors Under Realistic Knowledge Conflicts. arXiv:2404.16032v1, 2024. \url{https://arxiv.org/abs/2404.16032v1}
\bibitem{ref10} The Power of Noise Redefining Retrieval for RAG Systems. arXiv:2401.14887v3, 2024. \url{https://arxiv.org/abs/2401.14887v3}
\bibitem{ref11} On Single and Multiple Representations in Dense Passage Retrieval. arXiv:2108.06279v2, 2021. \url{https://arxiv.org/abs/2108.06279v2}
\bibitem{ref12} Dense Passage Retrieval for Open-Domain Question Answering. arXiv:2004.04906v3, 2020. \url{https://arxiv.org/abs/2004.04906v3}
\bibitem{ref13} GNN-encoder Learning a Dual-encoder Architecture via Graph Neural Networks for Dense Passage Retrieval. arXiv:2204.08241v2, 2022. \url{https://arxiv.org/abs/2204.08241v2}
\bibitem{ref14} Topic-DPR Topic-based Prompts for Dense Passage Retrieval. arXiv:2310.06626v1, 2023. \url{https://arxiv.org/abs/2310.06626v1}
\bibitem{ref15} Retrieval-Augmented Generation Is Dense Passage Retrieval Retrieving. arXiv:2402.11035v2, 2024. \url{https://arxiv.org/abs/2402.11035v2}
\bibitem{ref16} A Replication Study of Dense Passage Retriever. arXiv:2104.05740v1, 2021. \url{https://arxiv.org/abs/2104.05740v1}
\bibitem{ref17} CITADEL Conditional Token Interaction via Dynamic Lexical Routing for Efficient and Effective Multi-Vector Retrieval. arXiv:2211.10411v1, 2022. \url{https://arxiv.org/abs/2211.10411v1}
\bibitem{ref18} M3 A Multi-Task Mixed-Objective Learning Framework for Open-Domain Multi-Hop Dense Sentence Retrieval. arXiv:2403.14074v1, 2024. \url{https://arxiv.org/abs/2403.14074v1}
\bibitem{ref19} Hybrid Retrieval and Multi-stage Text Ranking Solution at TREC 2022 Deep Learning Track. arXiv:2308.12039v1, 2023. \url{https://arxiv.org/abs/2308.12039v1}
\bibitem{ref20} Phrase Retrieval Learns Passage Retrieval, Too. arXiv:2109.08133v1, 2021. \url{https://arxiv.org/abs/2109.08133v1}
\bibitem{ref21} Embedding-based Zero-shot Retrieval through Query Generation. arXiv:2009.10270v1, 2020. \url{https://arxiv.org/abs/2009.10270v1}
\bibitem{ref22} Coupled intrinsic and extrinsic human language resource-based query expansion. arXiv:2004.11083v1, 2020. \url{https://arxiv.org/abs/2004.11083v1}
\bibitem{ref23} Query Expansion Strategy based on Pseudo Relevance Feedback and Term Weight Scheme for Monolingual Retrieval. arXiv:1502.05168v1, 2015. \url{https://arxiv.org/abs/1502.05168v1}
\bibitem{ref24} Query Expansion Using Contextual Clue Sampling with Language Models. arXiv:2210.07093v1, 2022. \url{https://arxiv.org/abs/2210.07093v1}
\bibitem{ref25} Exploring Query Categorisation for Query Expansion A Study. arXiv:1509.05567v1, 2015. \url{https://arxiv.org/abs/1509.05567v1}
\bibitem{ref26} A new approach for query expansion using Wikipedia and WordNet. arXiv:1901.10197v2, 2019. \url{https://arxiv.org/abs/1901.10197v2}
\bibitem{ref27} A Linguistically Driven Framework for Query Expansion via Grammatical Constituent Highlighting and Role-Based Concept Weighting. arXiv:2004.13481v1, 2020. \url{https://arxiv.org/abs/2004.13481v1}
\bibitem{ref28} Natural language technology and query expansion issues, state-of-the-art and perspectives. arXiv:2004.11093v1, 2020. \url{https://arxiv.org/abs/2004.11093v1}
\bibitem{ref29} Query Expansion by Prompting Large Language Models. arXiv:2305.03653v1, 2023. \url{https://arxiv.org/abs/2305.03653v1}
\bibitem{ref30} MILL Mutual Verification with Large Language Models for Zero-Shot Query Expansion. arXiv:2310.19056v3, 2023. \url{https://arxiv.org/abs/2310.19056v3}
\bibitem{ref31} GenQREnsemble Zero-Shot LLM Ensemble Prompting for Generative Query Reformulation. arXiv:2404.03746v1, 2024. \url{https://arxiv.org/abs/2404.03746v1}
\bibitem{ref32} ConQueR Contextualized Query Reduction using Search Logs. arXiv:2305.12662v1, 2023. \url{https://arxiv.org/abs/2305.12662v1}
\bibitem{ref33} Query Rewriting via Large Language Models. arXiv:2403.09060v1, 2024. \url{https://arxiv.org/abs/2403.09060v1}
\bibitem{ref34} Massive Query Expansion by Exploiting Graph Knowledge Bases. arXiv:1310.5698v1, 2013. \url{https://arxiv.org/abs/1310.5698v1}
\bibitem{ref35} Query Augmentation by Decoding Semantics from Brain Signals. arXiv:2402.15708v2, 2024. \url{https://arxiv.org/abs/2402.15708v2}
\bibitem{ref36} Deep Neural Networks for Query Expansion using Word Embeddings. arXiv:1811.03514v1, 2018. \url{https://arxiv.org/abs/1811.03514v1}
\bibitem{ref37} Multi-Task Retrieval-Augmented Text Generation with Relevance Sampling. arXiv:2207.03030v1, 2022. \url{https://arxiv.org/abs/2207.03030v1}
\bibitem{ref38} Improving Natural Language Understanding with Computation-Efficient Retrieval Representation Fusion. arXiv:2401.02993v1, 2024. \url{https://arxiv.org/abs/2401.02993v1}
\bibitem{ref39} Re2G Retrieve, Rerank, Generate. arXiv:2207.06300v1, 2022. \url{https://arxiv.org/abs/2207.06300v1}
\bibitem{ref40} Fine- and Coarse-Granularity Hybrid Self-Attention for Efficient BERT. arXiv:2203.09055v1, 2022. \url{https://arxiv.org/abs/2203.09055v1}
\bibitem{ref41} Continuous Decomposition of Granularity for Neural Paraphrase Generation. arXiv:2209.01765v2, 2022. \url{https://arxiv.org/abs/2209.01765v2}
\bibitem{ref42} You Need to Read Again Multi-granularity Perception Network for Moment Retrieval in Videos. arXiv:2205.12886v2, 2022. \url{https://arxiv.org/abs/2205.12886v2}
\bibitem{ref43} Bifurcated Attention for Single-Context Large-Batch Sampling. arXiv:2403.08845v1, 2024. \url{https://arxiv.org/abs/2403.08845v1}
\bibitem{ref44} Metacognitive Retrieval-Augmented Large Language Models. arXiv:2402.11626v1, 2024. \url{https://arxiv.org/abs/2402.11626v1}
\bibitem{ref45} Task-Aware Information Routing from Common Representation Space in Lifelong Learning. arXiv:2302.11346v1, 2023. \url{https://arxiv.org/abs/2302.11346v1}
\bibitem{ref46} UniRQR A Unified Model for Retrieval Decision, Query, and Response Generation in Internet-Based Knowledge Dialogue Systems. arXiv:2401.06811v1, 2024. \url{https://arxiv.org/abs/2401.06811v1}
\bibitem{ref47} CONFLARE CONFormal LArge language model REtrieval. arXiv:2404.04287v1, 2024. \url{https://arxiv.org/abs/2404.04287v1}
\bibitem{ref48} HGOT Hierarchical Graph of Thoughts for Retrieval-Augmented In-Context Learning in Factuality Evaluation. arXiv:2402.09390v1, 2024. \url{https://arxiv.org/abs/2402.09390v1}
\bibitem{ref49} Collaboration and Transition Distilling Item Transitions into Multi-Query Self-Attention for Sequential Recommendation. arXiv:2311.01056v2, 2023. \url{https://arxiv.org/abs/2311.01056v2}
\bibitem{ref50} Conditional Self-Attention for Query-based Summarization. arXiv:2002.07338v1, 2020. \url{https://arxiv.org/abs/2002.07338v1}
\bibitem{ref51} CRUD-RAG A Comprehensive Chinese Benchmark for Retrieval-Augmented Generation of Large Language Models. arXiv:2401.17043v2, 2024. \url{https://arxiv.org/abs/2401.17043v2}
\bibitem{ref52} Unsupervised Information Refinement Training of Large Language Models for Retrieval-Augmented Generation. arXiv:2402.18150v1, 2024. \url{https://arxiv.org/abs/2402.18150v1}
\bibitem{ref53} Are Large Language Models Good at Utility Judgments. arXiv:2403.19216v1, 2024. \url{https://arxiv.org/abs/2403.19216v1}
\bibitem{ref54} Mafin Enhancing Black-Box Embeddings with Model Augmented Fine-Tuning. arXiv:2402.12177v4, 2024. \url{https://arxiv.org/abs/2402.12177v4}
\bibitem{ref55} RAM Towards an Ever-Improving Memory System by Learning from Communications. arXiv:2404.12045v1, 2024. \url{https://arxiv.org/abs/2404.12045v1}
\bibitem{ref56} Dynamic Contexts for Generating Suggestion Questions in RAG Based Conversational Systems. arXiv:2403.11413v1, 2024. \url{https://arxiv.org/abs/2403.11413v1}
\bibitem{ref57} PipeRAG Fast Retrieval-Augmented Generation via Algorithm-System Co-design. arXiv:2403.05676v1, 2024. \url{https://arxiv.org/abs/2403.05676v1}
\bibitem{ref58} RAGCache Efficient Knowledge Caching for Retrieval-Augmented Generation. arXiv:2404.12457v2, 2024. \url{https://arxiv.org/abs/2404.12457v2}
\bibitem{ref59} Seven Failure Points When Engineering a Retrieval Augmented Generation System. arXiv:2401.05856v1, 2024. \url{https://arxiv.org/abs/2401.05856v1}
\bibitem{ref60} JORA JAX Tensor-Parallel LoRA Library for Retrieval Augmented Fine-Tuning. arXiv:2403.11366v2, 2024. \url{https://arxiv.org/abs/2403.11366v2}
\bibitem{ref61} Accelerator-level Parallelism. arXiv:1907.02064v5, 2019. \url{https://arxiv.org/abs/1907.02064v5}
\bibitem{ref62} Cache Performance Study of Portfolio-Based Parallel CDCL SAT Solvers. arXiv:1309.3187v1, 2013. \url{https://arxiv.org/abs/1309.3187v1}
\bibitem{ref63} SMURF Efficient and Scalable Metadata Access for Distributed Applications. arXiv:2105.14157v1, 2021. \url{https://arxiv.org/abs/2105.14157v1}
\bibitem{ref64} Coherence Traffic in Manycore Processors with Opaque Distributed Directories. arXiv:2011.05422v1, 2020. \url{https://arxiv.org/abs/2011.05422v1}
\bibitem{ref65} Local Self-Attention over Long Text for Efficient Document Retrieval. arXiv:2005.04908v1, 2020. \url{https://arxiv.org/abs/2005.04908v1}
\bibitem{ref66} Pseudo-Relevance Feedback for Multiple Representation Dense Retrieval. arXiv:2106.11251v2, 2021. \url{https://arxiv.org/abs/2106.11251v2}
\bibitem{ref67} Citations as Queries Source Attribution Using Language Models as Rerankers. arXiv:2306.17322v1, 2023. \url{https://arxiv.org/abs/2306.17322v1}
\bibitem{ref68} CODER An efficient framework for improving retrieval through COntextual Document Embedding Reranking. arXiv:2112.08766v3, 2021. \url{https://arxiv.org/abs/2112.08766v3}
\bibitem{ref69} No Parameter Left Behind How Distillation and Model Size Affect Zero-Shot Retrieval. arXiv:2206.02873v5, 2022. \url{https://arxiv.org/abs/2206.02873v5}
\bibitem{ref70} Incorporating Relevance Feedback for Information-Seeking Retrieval using Few-Shot Document Re-Ranking. arXiv:2210.10695v1, 2022. \url{https://arxiv.org/abs/2210.10695v1}
\bibitem{ref71} Re3val Reinforced and Reranked Generative Retrieval. arXiv:2401.16979v3, 2024. \url{https://arxiv.org/abs/2401.16979v3}
\bibitem{ref72} Less is Less When Are Snippets Insufficient for Human vs Machine Relevance Estimation. arXiv:2201.08721v1, 2022. \url{https://arxiv.org/abs/2201.08721v1}
\bibitem{ref73} Retrieval-Augmented Generation for Knowledge-Intensive NLP Tasks. arXiv:2005.11401v4, 2020. \url{https://arxiv.org/abs/2005.11401v4}
\bibitem{ref74} Fine-Tuning or Retrieval Comparing Knowledge Injection in LLMs. arXiv:2312.05934v3, 2023. \url{https://arxiv.org/abs/2312.05934v3}
\bibitem{ref75} FIT-RAG Black-Box RAG with Factual Information and Token Reduction. arXiv:2403.14374v1, 2024. \url{https://arxiv.org/abs/2403.14374v1}
\bibitem{ref76} Establishing Performance Baselines in Fine-Tuning, Retrieval-Augmented Generation and Soft-Prompting for Non-Specialist LLM Users. arXiv:2311.05903v2, 2023. \url{https://arxiv.org/abs/2311.05903v2}
\bibitem{ref77} Fine Tuning vs. Retrieval Augmented Generation for Less Popular Knowledge. arXiv:2403.01432v2, 2024. \url{https://arxiv.org/abs/2403.01432v2}
\bibitem{ref78} Improving the Domain Adaptation of Retrieval Augmented Generation (RAG) Models for Open Domain Question Answering. arXiv:2210.02627v1, 2022. \url{https://arxiv.org/abs/2210.02627v1}
\bibitem{ref79} Reinforcement Retrieval Leveraging Fine-grained Feedback for Fact Checking News Claims with Black-Box LLM. arXiv:2404.17283v1, 2024. \url{https://arxiv.org/abs/2404.17283v1}
\bibitem{ref80} Bridging the Preference Gap between Retrievers and LLMs. arXiv:2401.06954v2, 2024. \url{https://arxiv.org/abs/2401.06954v2}
\bibitem{ref81} Fine-tuning Language Models with Generative Adversarial Reward Modelling. arXiv:2305.06176v3, 2023. \url{https://arxiv.org/abs/2305.06176v3}
\bibitem{ref82} Unsupervised Large Language Model Alignment for Information Retrieval via Contrastive Feedback. arXiv:2309.17078v2, 2023. \url{https://arxiv.org/abs/2309.17078v2}
\bibitem{ref83} CLHA A Simple yet Effective Contrastive Learning Framework for Human Alignment. arXiv:2403.16649v2, 2024. \url{https://arxiv.org/abs/2403.16649v2}
\bibitem{ref84} Improving Reinforcement Learning from Human Feedback Using Contrastive Rewards. arXiv:2403.07708v2, 2024. \url{https://arxiv.org/abs/2403.07708v2}
\bibitem{ref85} Reinforcement Learning from Reflective Feedback (RLRF) Aligning and Improving LLMs via Fine-Grained Self-Reflection. arXiv:2403.14238v1, 2024. \url{https://arxiv.org/abs/2403.14238v1}
\bibitem{ref86} Parameter-Efficient Fine-Tuning for Large Models A Comprehensive Survey. arXiv:2403.14608v4, 2024. \url{https://arxiv.org/abs/2403.14608v4}
\bibitem{ref87} ComPEFT Compression for Communicating Parameter Efficient Updates via Sparsification and Quantization. arXiv:2311.13171v1, 2023. \url{https://arxiv.org/abs/2311.13171v1}
\bibitem{ref88} Composable constraints. arXiv:2112.06818v1, 2021. \url{https://arxiv.org/abs/2112.06818v1}
\bibitem{ref89} HyperLoRA for PDEs. arXiv:2308.09290v1, 2023. \url{https://arxiv.org/abs/2308.09290v1}
\bibitem{ref90} A Rank Stabilization Scaling Factor for Fine-Tuning with LoRA. arXiv:2312.03732v1, 2023. \url{https://arxiv.org/abs/2312.03732v1}
\bibitem{ref91} Algorand. arXiv:1607.01341v9, 2016. \url{https://arxiv.org/abs/1607.01341v9}
\bibitem{ref92} Mixtral of Experts. arXiv:2401.04088v1, 2024. \url{https://arxiv.org/abs/2401.04088v1}
\bibitem{ref93} Introducing Routing Functions to Vision-Language Parameter-Efficient Fine-Tuning with Low-Rank Bottlenecks. arXiv:2403.09377v1, 2024. \url{https://arxiv.org/abs/2403.09377v1}
\bibitem{ref94} Higher Layers Need More LoRA Experts. arXiv:2402.08562v1, 2024. \url{https://arxiv.org/abs/2402.08562v1}
\bibitem{ref95} A Note on LoRA. arXiv:2404.05086v1, 2024. \url{https://arxiv.org/abs/2404.05086v1}
\bibitem{ref96} Advancing Parameter Efficiency in Fine-tuning via Representation Editing. arXiv:2402.15179v2, 2024. \url{https://arxiv.org/abs/2402.15179v2}
\bibitem{ref97} Typos that Broke the RAG's Back Genetic Attack on RAG Pipeline by Simulating Documents in the Wild via Low-level Perturbations. arXiv:2404.13948v1, 2024. \url{https://arxiv.org/abs/2404.13948v1}
\bibitem{ref98} Catastrophic Overfitting A Potential Blessing in Disguise. arXiv:2402.18211v1, 2024. \url{https://arxiv.org/abs/2402.18211v1}
\bibitem{ref99} Ensemble Adversarial Training Attacks and Defenses. arXiv:1705.07204v5, 2017. \url{https://arxiv.org/abs/1705.07204v5}
\bibitem{ref100} Parameter-Saving Adversarial Training Reinforcing Multi-Perturbation Robustness via Hypernetworks. arXiv:2309.16207v1, 2023. \url{https://arxiv.org/abs/2309.16207v1}
\bibitem{ref101} Instance adaptive adversarial training Improved accuracy tradeoffs in neural nets. arXiv:1910.08051v1, 2019. \url{https://arxiv.org/abs/1910.08051v1}
\bibitem{ref102} Blind Adversarial Training Balance Accuracy and Robustness. arXiv:2004.05914v1, 2020. \url{https://arxiv.org/abs/2004.05914v1}
\bibitem{ref103} Learning to Learn from Mistakes Robust Optimization for Adversarial Noise. arXiv:2008.05247v1, 2020. \url{https://arxiv.org/abs/2008.05247v1}
\bibitem{ref104} Data Augmentation Can Improve Robustness. arXiv:2111.05328v1, 2021. \url{https://arxiv.org/abs/2111.05328v1}
\bibitem{ref105} Learning to Generate Noise for Multi-Attack Robustness. arXiv:2006.12135v4, 2020. \url{https://arxiv.org/abs/2006.12135v4}
\bibitem{ref106} RDA Robust Domain Adaptation via Fourier Adversarial Attacking. arXiv:2106.02874v3, 2021. \url{https://arxiv.org/abs/2106.02874v3}
\bibitem{ref107} Improving Noise Robustness of Contrastive Speech Representation Learning with Speech Reconstruction. arXiv:2110.15430v1, 2021. \url{https://arxiv.org/abs/2110.15430v1}
\bibitem{ref108} Enhancing LLM Intelligence with ARM-RAG Auxiliary Rationale Memory for Retrieval Augmented Generation. arXiv:2311.04177v1, 2023. \url{https://arxiv.org/abs/2311.04177v1}
\bibitem{ref109} Self-RAG Learning to Retrieve, Generate, and Critique through Self-Reflection. arXiv:2310.11511v1, 2023. \url{https://arxiv.org/abs/2310.11511v1}
\bibitem{ref110} Simple Domain Adaptation for Sparse Retrievers. arXiv:2401.11509v1, 2024. \url{https://arxiv.org/abs/2401.11509v1}
\bibitem{ref111} Linguistically Informed Masking for Representation Learning in the Patent Domain. arXiv:2106.05768v1, 2021. \url{https://arxiv.org/abs/2106.05768v1}
\bibitem{ref112} Learning Invariant Representation with Consistency and Diversity for Semi-supervised Source Hypothesis Transfer. arXiv:2107.03008v2, 2021. \url{https://arxiv.org/abs/2107.03008v2}
\bibitem{ref113} Low-confidence Samples Matter for Domain Adaptation. arXiv:2202.02802v2, 2022. \url{https://arxiv.org/abs/2202.02802v2}
\bibitem{ref114} Universal Semi-Supervised Domain Adaptation by Mitigating Common-Class Bias. arXiv:2403.11234v1, 2024. \url{https://arxiv.org/abs/2403.11234v1}
\bibitem{ref115} Pred\&Guide Labeled Target Class Prediction for Guiding Semi-Supervised Domain Adaptation. arXiv:2211.11975v1, 2022. \url{https://arxiv.org/abs/2211.11975v1}
\bibitem{ref116} Multi-modal Domain Adaptation for REG via Relation Transfer. arXiv:2309.13247v1, 2023. \url{https://arxiv.org/abs/2309.13247v1}
\bibitem{ref117} Discovering Domain Disentanglement for Generalized Multi-source Domain Adaptation. arXiv:2207.05070v1, 2022. \url{https://arxiv.org/abs/2207.05070v1}
\bibitem{ref118} Minimum Class Confusion for Versatile Domain Adaptation. arXiv:1912.03699v3, 2019. \url{https://arxiv.org/abs/1912.03699v3}
\bibitem{ref119} Semi-Supervised Hypothesis Transfer for Source-Free Domain Adaptation. arXiv:2107.06735v1, 2021. \url{https://arxiv.org/abs/2107.06735v1}
\bibitem{ref120} Semi-supervised Domain Adaptive Medical Image Segmentation through Consistency Regularized Disentangled Contrastive Learning. arXiv:2307.02798v1, 2023. \url{https://arxiv.org/abs/2307.02798v1}
\bibitem{ref121} Using Explanations to Guide Models. arXiv:2303.11932v1, 2023. \url{https://arxiv.org/abs/2303.11932v1}
\bibitem{ref122} Constructing and Exploring Intermediate Domains in Mixed Domain Semi-supervised Medical Image Segmentation. arXiv:2404.08951v1, 2024. \url{https://arxiv.org/abs/2404.08951v1}
\bibitem{ref123} HaluEval A Large-Scale Hallucination Evaluation Benchmark for Large Language Models. arXiv:2305.11747v3, 2023. \url{https://arxiv.org/abs/2305.11747v3}
\bibitem{ref124} HaluEval-Wild Evaluating Hallucinations of Language Models in the Wild. arXiv:2403.04307v1, 2024. \url{https://arxiv.org/abs/2403.04307v1}
\bibitem{ref125} NoMIRACL Knowing When You Don't Know for Robust Multilingual Retrieval-Augmented Generation. arXiv:2312.11361v2, 2023. \url{https://arxiv.org/abs/2312.11361v2}
\bibitem{ref126} DiaHalu A Dialogue-level Hallucination Evaluation Benchmark for Large Language Models. arXiv:2403.00896v1, 2024. \url{https://arxiv.org/abs/2403.00896v1}
\bibitem{ref127} TofuEval Evaluating Hallucinations of LLMs on Topic-Focused Dialogue Summarization. arXiv:2402.13249v2, 2024. \url{https://arxiv.org/abs/2402.13249v2}
\bibitem{ref128} AttributionBench How Hard is Automatic Attribution Evaluation. arXiv:2402.15089v1, 2024. \url{https://arxiv.org/abs/2402.15089v1}
\bibitem{ref129} Measuring and Reducing LLM Hallucination without Gold-Standard Answers via Expertise-Weighting. arXiv:2402.10412v1, 2024. \url{https://arxiv.org/abs/2402.10412v1}
\bibitem{ref130} UFO a Unified and Flexible Framework for Evaluating Factuality of Large Language Models. arXiv:2402.14690v1, 2024. \url{https://arxiv.org/abs/2402.14690v1}
\bibitem{ref131} Chainpoll A high efficacy method for LLM hallucination detection. arXiv:2310.18344v1, 2023. \url{https://arxiv.org/abs/2310.18344v1}
\bibitem{ref132} MMBench Is Your Multi-modal Model an All-around Player. arXiv:2307.06281v3, 2023. \url{https://arxiv.org/abs/2307.06281v3}
\bibitem{ref133} VALOR-EVAL Holistic Coverage and Faithfulness Evaluation of Large Vision-Language Models. arXiv:2404.13874v1, 2024. \url{https://arxiv.org/abs/2404.13874v1}
\bibitem{ref134} INSIDE LLMs' Internal States Retain the Power of Hallucination Detection. arXiv:2402.03744v1, 2024. \url{https://arxiv.org/abs/2402.03744v1}
\bibitem{ref135} C-RAG Certified Generation Risks for Retrieval-Augmented Language Models. arXiv:2402.03181v3, 2024. \url{https://arxiv.org/abs/2402.03181v3}
\bibitem{ref136} Conformal Prediction with Missing Values. arXiv:2306.02732v1, 2023. \url{https://arxiv.org/abs/2306.02732v1}
\bibitem{ref137} Robust Conformal Prediction under Distribution Shift via Physics-Informed Structural Causal Model. arXiv:2403.15025v1, 2024. \url{https://arxiv.org/abs/2403.15025v1}
\bibitem{ref138} Conformal Predictive Systems Under Covariate Shift. arXiv:2404.15018v1, 2024. \url{https://arxiv.org/abs/2404.15018v1}
\bibitem{ref139} Cross-Validation Conformal Risk Control. arXiv:2401.11974v1, 2024. \url{https://arxiv.org/abs/2401.11974v1}
\bibitem{ref140} Fact-Checking the Output of Large Language Models via Token-Level Uncertainty Quantification. arXiv:2403.04696v1, 2024. \url{https://arxiv.org/abs/2403.04696v1}
\bibitem{ref141} Sequence-Level Certainty Reduces Hallucination In Knowledge-Grounded Dialogue Generation. arXiv:2310.18794v3, 2023. \url{https://arxiv.org/abs/2310.18794v3}
\bibitem{ref142} In-Context Sharpness as Alerts An Inner Representation Perspective for Hallucination Mitigation. arXiv:2403.01548v3, 2024. \url{https://arxiv.org/abs/2403.01548v3}
\bibitem{ref143} KCTS Knowledge-Constrained Tree Search Decoding with Token-Level Hallucination Detection. arXiv:2310.09044v1, 2023. \url{https://arxiv.org/abs/2310.09044v1}
\bibitem{ref144} SAC3 Reliable Hallucination Detection in Black-Box Language Models via Semantic-aware Cross-check Consistency. arXiv:2311.01740v2, 2023. \url{https://arxiv.org/abs/2311.01740v2}
\bibitem{ref145} Enhancing Uncertainty-Based Hallucination Detection with Stronger Focus. arXiv:2311.13230v1, 2023. \url{https://arxiv.org/abs/2311.13230v1}
\bibitem{ref146} FACTOID FACtual enTailment fOr hallucInation Detection. arXiv:2403.19113v1, 2024. \url{https://arxiv.org/abs/2403.19113v1}
\bibitem{ref147} Improved Beam Search for Hallucination Mitigation in Abstractive Summarization. arXiv:2212.02712v2, 2022. \url{https://arxiv.org/abs/2212.02712v2}
\bibitem{ref148} Elastic Weight Removal for Faithful and Abstractive Dialogue Generation. arXiv:2303.17574v1, 2023. \url{https://arxiv.org/abs/2303.17574v1}
\bibitem{ref149} OPERA Alleviating Hallucination in Multi-Modal Large Language Models via Over-Trust Penalty and Retrospection-Allocation. arXiv:2311.17911v3, 2023. \url{https://arxiv.org/abs/2311.17911v3}
\bibitem{ref150} Retrieve Only When It Needs Adaptive Retrieval Augmentation for Hallucination Mitigation in Large Language Models. arXiv:2402.10612v1, 2024. \url{https://arxiv.org/abs/2402.10612v1}
\bibitem{ref151} Characterizing Attribution and Fluency Tradeoffs for Retrieval-Augmented Large Language Models. arXiv:2302.05578v2, 2023. \url{https://arxiv.org/abs/2302.05578v2}
\bibitem{ref152} AutoDAN Interpretable Gradient-Based Adversarial Attacks on Large Language Models. arXiv:2310.15140v2, 2023. \url{https://arxiv.org/abs/2310.15140v2}
\bibitem{ref153} A Survey of Black-Box Adversarial Attacks on Computer Vision Models. arXiv:1912.01667v3, 2019. \url{https://arxiv.org/abs/1912.01667v3}
\bibitem{ref154} Adversarial Learning in Statistical Classification A Comprehensive Review of Defenses Against Attacks. arXiv:1904.06292v3, 2019. \url{https://arxiv.org/abs/1904.06292v3}
\bibitem{ref155} Generalizable Black-Box Adversarial Attack with Meta Learning. arXiv:2301.00364v1, 2023. \url{https://arxiv.org/abs/2301.00364v1}
\bibitem{ref156} Spanning Attack Reinforce Black-box Attacks with Unlabeled Data. arXiv:2005.04871v2, 2020. \url{https://arxiv.org/abs/2005.04871v2}
\bibitem{ref157} An Efficient and Margin-Approaching Zero-Confidence Adversarial Attack. arXiv:1910.00511v1, 2019. \url{https://arxiv.org/abs/1910.00511v1}
\bibitem{ref158} Local Black-box Adversarial Attacks A Query Efficient Approach. arXiv:2101.01032v1, 2021. \url{https://arxiv.org/abs/2101.01032v1}
\bibitem{ref159} Investigating the prompt leakage effect and black-box defenses for multi-turn LLM interactions. arXiv:2404.16251v2, 2024. \url{https://arxiv.org/abs/2404.16251v2}
\bibitem{ref160} Dialectical Alignment Resolving the Tension of 3H and Security Threats of LLMs. arXiv:2404.00486v1, 2024. \url{https://arxiv.org/abs/2404.00486v1}
\bibitem{ref161} Maatphor Automated Variant Analysis for Prompt Injection Attacks. arXiv:2312.11513v1, 2023. \url{https://arxiv.org/abs/2312.11513v1}
\bibitem{ref162} Combining Confidence Elicitation and Sample-based Methods for Uncertainty Quantification in Misinformation Mitigation. arXiv:2401.08694v2, 2024. \url{https://arxiv.org/abs/2401.08694v2}
\bibitem{ref163} The State of Human-centered NLP Technology for Fact-checking. arXiv:2301.03056v1, 2023. \url{https://arxiv.org/abs/2301.03056v1}
\bibitem{ref164} A Stackelberg Game Perspective on the Conflict Between Machine Learning and Data Obfuscation. arXiv:1606.06771v2, 2016. \url{https://arxiv.org/abs/1606.06771v2}
\bibitem{ref165} Research Challenges in Designing Differentially Private Text Generation Mechanisms. arXiv:2012.05403v1, 2020. \url{https://arxiv.org/abs/2012.05403v1}
\bibitem{ref166} Reducing Privacy Risks in Online Self-Disclosures with Language Models. arXiv:2311.09538v2, 2023. \url{https://arxiv.org/abs/2311.09538v2}
\bibitem{ref167} Privacy Preserving Recalibration under Domain Shift. arXiv:2008.09643v1, 2020. \url{https://arxiv.org/abs/2008.09643v1}
\bibitem{ref168} Entity-Based Knowledge Conflicts in Question Answering. arXiv:2109.05052v2, 2021. \url{https://arxiv.org/abs/2109.05052v2}
\bibitem{ref169} On Hallucination and Predictive Uncertainty in Conditional Language Generation. arXiv:2103.15025v1, 2021. \url{https://arxiv.org/abs/2103.15025v1}
\bibitem{ref170} Quantifying and Attributing the Hallucination of Large Language Models via Association Analysis. arXiv:2309.05217v1, 2023. \url{https://arxiv.org/abs/2309.05217v1}
\bibitem{ref171} HypoTermQA Hypothetical Terms Dataset for Benchmarking Hallucination Tendency of LLMs. arXiv:2402.16211v1, 2024. \url{https://arxiv.org/abs/2402.16211v1}
\bibitem{ref172} Fakes of Varying Shades How Warning Affects Human Perception and Engagement Regarding LLM Hallucinations. arXiv:2404.03745v1, 2024. \url{https://arxiv.org/abs/2404.03745v1}
\bibitem{ref173} Biomedical knowledge graph-enhanced prompt generation for large language models. arXiv:2311.17330v1, 2023. \url{https://arxiv.org/abs/2311.17330v1}
\bibitem{ref174} REALM RAG-Driven Enhancement of Multimodal Electronic Health Records Analysis via Large Language Models. arXiv:2402.07016v1, 2024. \url{https://arxiv.org/abs/2402.07016v1}
\bibitem{ref175} Benchmarking Retrieval-Augmented Generation for Medicine. arXiv:2402.13178v2, 2024. \url{https://arxiv.org/abs/2402.13178v2}
\bibitem{ref176} Improving Medical Reasoning through Retrieval and Self-Reflection with Retrieval-Augmented Large Language Models. arXiv:2401.15269v2, 2024. \url{https://arxiv.org/abs/2401.15269v2}
\bibitem{ref177} ActiveRAG Revealing the Treasures of Knowledge via Active Learning. arXiv:2402.13547v1, 2024. \url{https://arxiv.org/abs/2402.13547v1}
\bibitem{ref178} Development and Testing of a Novel Large Language Model-Based Clinical Decision Support Systems for Medication Safety in 12 Clinical Specialties. arXiv:2402.01741v2, 2024. \url{https://arxiv.org/abs/2402.01741v2}
\bibitem{ref179} Development and Testing of Retrieval Augmented Generation in Large Language Models -- A Case Study Report. arXiv:2402.01733v1, 2024. \url{https://arxiv.org/abs/2402.01733v1}
\bibitem{ref180} IryoNLP at MEDIQA-CORR 2024 Tackling the Medical Error Detection \& Correction Task On the Shoulders of Medical Agents. arXiv:2404.15488v1, 2024. \url{https://arxiv.org/abs/2404.15488v1}
\bibitem{ref181} Zero-shot Medical Entity Retrieval without Annotation Learning From Rich Knowledge Graph Semantics. arXiv:2105.12682v1, 2021. \url{https://arxiv.org/abs/2105.12682v1}
\bibitem{ref182} Unlocking Multi-View Insights in Knowledge-Dense Retrieval-Augmented Generation. arXiv:2404.12879v1, 2024. \url{https://arxiv.org/abs/2404.12879v1}
\bibitem{ref183} JMLR Joint Medical LLM and Retrieval Training for Enhancing Reasoning and Professional Question Answering Capability. arXiv:2402.17887v3, 2024. \url{https://arxiv.org/abs/2402.17887v3}
\bibitem{ref184} NeCo@ALQAC 2023 Legal Domain Knowledge Acquisition for Low-Resource Languages through Data Enrichment. arXiv:2309.05500v1, 2023. \url{https://arxiv.org/abs/2309.05500v1}
\bibitem{ref185} MultiLegalPile A 689GB Multilingual Legal Corpus. arXiv:2306.02069v2, 2023. \url{https://arxiv.org/abs/2306.02069v2}
\bibitem{ref186} MultiEURLEX -- A multi-lingual and multi-label legal document classification dataset for zero-shot cross-lingual transfer. arXiv:2109.00904v2, 2021. \url{https://arxiv.org/abs/2109.00904v2}
\bibitem{ref187} Cross-Lingual Retrieval Augmented Prompt for Low-Resource Languages. arXiv:2212.09651v4, 2022. \url{https://arxiv.org/abs/2212.09651v4}
\bibitem{ref188} Cross-Lingual Question Answering over Knowledge Base as Reading Comprehension. arXiv:2302.13241v1, 2023. \url{https://arxiv.org/abs/2302.13241v1}
\bibitem{ref189} Cross-Lingual Transfer Robustness to Lower-Resource Languages on Adversarial Datasets. arXiv:2403.20056v1, 2024. \url{https://arxiv.org/abs/2403.20056v1}
\bibitem{ref190} M3Exam A Multilingual, Multimodal, Multilevel Benchmark for Examining Large Language Models. arXiv:2306.05179v2, 2023. \url{https://arxiv.org/abs/2306.05179v2}
\bibitem{ref191} Boosting Zero-shot Cross-lingual Retrieval by Training on Artificially Code-Switched Data. arXiv:2305.05295v2, 2023. \url{https://arxiv.org/abs/2305.05295v2}
\bibitem{ref192} MULTI3NLU++ A Multilingual, Multi-Intent, Multi-Domain Dataset for Natural Language Understanding in Task-Oriented Dialogue. arXiv:2212.10455v2, 2022. \url{https://arxiv.org/abs/2212.10455v2}
\bibitem{ref193} MIA 2022 Shared Task Evaluating Cross-lingual Open-Retrieval Question Answering for 16 Diverse Languages. arXiv:2207.00758v1, 2022. \url{https://arxiv.org/abs/2207.00758v1}
\bibitem{ref194} iRAG An Incremental Retrieval Augmented Generation System for Videos. arXiv:2404.12309v1, 2024. \url{https://arxiv.org/abs/2404.12309v1}
\bibitem{ref195} VAST A Vision-Audio-Subtitle-Text Omni-Modality Foundation Model and Dataset. arXiv:2305.18500v2, 2023. \url{https://arxiv.org/abs/2305.18500v2}
\bibitem{ref196} Accommodating Audio Modality in CLIP for Multimodal Processing. arXiv:2303.06591v1, 2023. \url{https://arxiv.org/abs/2303.06591v1}
\bibitem{ref197} A Comprehensive Empirical Study of Vision-Language Pre-trained Model for Supervised Cross-Modal Retrieval. arXiv:2201.02772v2, 2022. \url{https://arxiv.org/abs/2201.02772v2}
\bibitem{ref198} Video-Teller Enhancing Cross-Modal Generation with Fusion and Decoupling. arXiv:2310.04991v3, 2023. \url{https://arxiv.org/abs/2310.04991v3}
\bibitem{ref199} MuMUR Multilingual Multimodal Universal Retrieval. arXiv:2208.11553v7, 2022. \url{https://arxiv.org/abs/2208.11553v7}
\bibitem{ref200} Lyrics Boosting Fine-grained Language-Vision Alignment and Comprehension via Semantic-aware Visual Objects. arXiv:2312.05278v2, 2023. \url{https://arxiv.org/abs/2312.05278v2}
\bibitem{ref201} VLCap Vision-Language with Contrastive Learning for Coherent Video Paragraph Captioning. arXiv:2206.12972v2, 2022. \url{https://arxiv.org/abs/2206.12972v2}
\bibitem{ref202} From Image to Video, what do we need in multimodal LLMs. arXiv:2404.11865v1, 2024. \url{https://arxiv.org/abs/2404.11865v1}
\bibitem{ref203} Vision-Language Integration in Multimodal Video Transformers (Partially) Aligns with the Brain. arXiv:2311.07766v1, 2023. \url{https://arxiv.org/abs/2311.07766v1}
\bibitem{ref204} ROS package search for robot software development a knowledge graph-based approach. arXiv:2312.14781v1, 2023. \url{https://arxiv.org/abs/2312.14781v1}
\bibitem{ref205} Retrieval Augmented Generation using Engineering Design Knowledge. arXiv:2307.06985v7, 2023. \url{https://arxiv.org/abs/2307.06985v7}
\bibitem{ref206} RAG-Driver Generalisable Driving Explanations with Retrieval-Augmented In-Context Learning in Multi-Modal Large Language Model. arXiv:2402.10828v1, 2024. \url{https://arxiv.org/abs/2402.10828v1}
\bibitem{ref207} RSG Fast Learning Adaptive Skills for Quadruped Robots by Skill Graph. arXiv:2311.06015v1, 2023. \url{https://arxiv.org/abs/2311.06015v1}
\bibitem{ref208} Automated Creation and Human-assisted Curation of Computable Scientific Models from Code and Text. arXiv:2202.13739v1, 2022. \url{https://arxiv.org/abs/2202.13739v1}
\bibitem{ref209} NASA Science Mission Directorate Knowledge Graph Discovery. arXiv:2303.10871v1, 2023. \url{https://arxiv.org/abs/2303.10871v1}
\bibitem{ref210} Introduction of an Assistance System to Support Domain Experts in Programming Low-code to Leverage Industry 5.0. arXiv:2212.04830v1, 2022. \url{https://arxiv.org/abs/2212.04830v1}
\bibitem{ref211} CLAPNQ Cohesive Long-form Answers from Passages in Natural Questions for RAG systems. arXiv:2404.02103v1, 2024. \url{https://arxiv.org/abs/2404.02103v1}
\bibitem{ref212} \$k\$NN-Adapter Efficient Domain Adaptation for Black-Box Language Models. arXiv:2302.10879v1, 2023. \url{https://arxiv.org/abs/2302.10879v1}
\bibitem{ref213} RAG vs Fine-tuning Pipelines, Tradeoffs, and a Case Study on Agriculture. arXiv:2401.08406v3, 2024. \url{https://arxiv.org/abs/2401.08406v3}
\bibitem{ref214} Enhancing Q\&A with Domain-Specific Fine-Tuning and Iterative Reasoning A Comparative Study. arXiv:2404.11792v2, 2024. \url{https://arxiv.org/abs/2404.11792v2}
\bibitem{ref215} Investigating the performance of Retrieval-Augmented Generation and fine-tuning for the development of AI-driven knowledge-based systems. arXiv:2403.09727v1, 2024. \url{https://arxiv.org/abs/2403.09727v1}
\bibitem{ref216} A Fine-tuning Enhanced RAG System with Quantized Influence Measure as AI Judge. arXiv:2402.17081v1, 2024. \url{https://arxiv.org/abs/2402.17081v1}
\bibitem{ref217} Compression, Generalization and Learning. arXiv:2301.12767v2, 2023. \url{https://arxiv.org/abs/2301.12767v2}
\bibitem{ref218} T-RAG Lessons from the LLM Trenches. arXiv:2402.07483v1, 2024. \url{https://arxiv.org/abs/2402.07483v1}
\bibitem{ref219} UniMS-RAG A Unified Multi-source Retrieval-Augmented Generation for Personalized Dialogue Systems. arXiv:2401.13256v1, 2024. \url{https://arxiv.org/abs/2401.13256v1}
\bibitem{ref220} Persona-DB Efficient Large Language Model Personalization for Response Prediction with Collaborative Data Refinement. arXiv:2402.11060v1, 2024. \url{https://arxiv.org/abs/2402.11060v1}
\bibitem{ref221} The Good and The Bad Exploring Privacy Issues in Retrieval-Augmented Generation (RAG). arXiv:2402.16893v1, 2024. \url{https://arxiv.org/abs/2402.16893v1}
\bibitem{ref222} Follow My Instruction and Spill the Beans Scalable Data Extraction from Retrieval-Augmented Generation Systems. arXiv:2402.17840v1, 2024. \url{https://arxiv.org/abs/2402.17840v1}
\bibitem{ref223} Building and Maintaining a Third-Party Library Supply Chain for Productive and Secure SGX Enclave Development. arXiv:2005.04367v1, 2020. \url{https://arxiv.org/abs/2005.04367v1}
\bibitem{ref224} A Comparative Qualitative and Quantitative Analysis of the Performance of Security Options for Message Protocols Fog Computing Scenario. arXiv:2210.12917v2, 2022. \url{https://arxiv.org/abs/2210.12917v2}
\bibitem{ref225} Evolving the Digital Industrial Infrastructure for Production Steps Taken and the Road Ahead. arXiv:2305.10151v1, 2023. \url{https://arxiv.org/abs/2305.10151v1}
\bibitem{ref226} H\$\_2\$O Heavy-Hitter Oracle for Efficient Generative Inference of Large Language Models. arXiv:2306.14048v3, 2023. \url{https://arxiv.org/abs/2306.14048v3}
\bibitem{ref227} TriForce Lossless Acceleration of Long Sequence Generation with Hierarchical Speculative Decoding. arXiv:2404.11912v2, 2024. \url{https://arxiv.org/abs/2404.11912v2}
\bibitem{ref228} A Modern Primer on Processing in Memory. arXiv:2012.03112v3, 2020. \url{https://arxiv.org/abs/2012.03112v3}
\bibitem{ref229} Heterogeneous Data-Centric Architectures for Modern Data-Intensive Applications Case Studies in Machine Learning and Databases. arXiv:2205.14664v1, 2022. \url{https://arxiv.org/abs/2205.14664v1}
\bibitem{ref230} Intelligent Architectures for Intelligent Machines. arXiv:2008.06112v1, 2020. \url{https://arxiv.org/abs/2008.06112v1}
\bibitem{ref231} Sequence can Secretly Tell You What to Discard. arXiv:2404.15949v1, 2024. \url{https://arxiv.org/abs/2404.15949v1}
\bibitem{ref232} Application-Driven Near-Data Processing for Similarity Search. arXiv:1606.03742v2, 2016. \url{https://arxiv.org/abs/1606.03742v2}
\bibitem{ref233} Evaluating the Potential of Disaggregated Memory Systems for HPC applications. arXiv:2306.04014v2, 2023. \url{https://arxiv.org/abs/2306.04014v2}
\bibitem{ref234} Cello Efficient Computer Systems Optimization with Predictive Early Termination and Censored Regression. arXiv:2204.04831v1, 2022. \url{https://arxiv.org/abs/2204.04831v1}
\bibitem{ref235} EdgeBERT Sentence-Level Energy Optimizations for Latency-Aware Multi-Task NLP Inference. arXiv:2011.14203v5, 2020. \url{https://arxiv.org/abs/2011.14203v5}
\bibitem{ref236} A Moveable Beast Partitioning Data and Compute for Computational Storage. arXiv:2212.11459v2, 2022. \url{https://arxiv.org/abs/2212.11459v2}
\bibitem{ref237} A simple multiforce layout for multiplex networks. arXiv:1607.03914v2, 2016. \url{https://arxiv.org/abs/1607.03914v2}
\bibitem{ref238} How faithful are RAG models Quantifying the tug-of-war between RAG and LLMs' internal prior. arXiv:2404.10198v1, 2024. \url{https://arxiv.org/abs/2404.10198v1}
\bibitem{ref239} Not All Contexts Are Equal Teaching LLMs Credibility-aware Generation. arXiv:2404.06809v1, 2024. \url{https://arxiv.org/abs/2404.06809v1}
\bibitem{ref240} RAGAR, Your Falsehood RADAR RAG-Augmented Reasoning for Political Fact-Checking using Multimodal Large Language Models. arXiv:2404.12065v1, 2024. \url{https://arxiv.org/abs/2404.12065v1}
\bibitem{ref241} How Stable is Knowledge Base Knowledge. arXiv:2211.00989v1, 2022. \url{https://arxiv.org/abs/2211.00989v1}
\bibitem{ref242} On Expansion and Contraction of DL-Lite Knowledge Bases. arXiv:2001.09365v1, 2020. \url{https://arxiv.org/abs/2001.09365v1}
\bibitem{ref243} Is Your LLM Outdated Benchmarking LLMs \& Alignment Algorithms for Time-Sensitive Knowledge. arXiv:2404.08700v1, 2024. \url{https://arxiv.org/abs/2404.08700v1}
\bibitem{ref244} Unsupervised Context Retrieval for Long-tail Entities. arXiv:1908.01798v1, 2019. \url{https://arxiv.org/abs/1908.01798v1}
\bibitem{ref245} 18 Million Links in Commit Messages Purpose, Evolution, and Decay. arXiv:2305.16591v1, 2023. \url{https://arxiv.org/abs/2305.16591v1}
\bibitem{ref246} Sustainability of Stack Exchange Q\textbackslash\{\}\&A communities the role of trust. arXiv:2205.07745v1, 2022. \url{https://arxiv.org/abs/2205.07745v1}
\bibitem{ref247} What About Feedback. arXiv:1412.4179v1, 2014. \url{https://arxiv.org/abs/1412.4179v1}
\bibitem{ref248} Intuitive or Dependent Investigating LLMs' Behavior Style to Conflicting Prompts. arXiv:2309.17415v3, 2023. \url{https://arxiv.org/abs/2309.17415v3}
\bibitem{ref249} PoisonedRAG Knowledge Poisoning Attacks to Retrieval-Augmented Generation of Large Language Models. arXiv:2402.07867v1, 2024. \url{https://arxiv.org/abs/2402.07867v1}
\bibitem{ref250} Learning to Poison Large Language Models During Instruction Tuning. arXiv:2402.13459v1, 2024. \url{https://arxiv.org/abs/2402.13459v1}
\bibitem{ref251} PINCH An Adversarial Extraction Attack Framework for Deep Learning Models. arXiv:2209.06300v2, 2022. \url{https://arxiv.org/abs/2209.06300v2}
\bibitem{ref252} Why Should Adversarial Perturbations be Imperceptible Rethink the Research Paradigm in Adversarial NLP. arXiv:2210.10683v1, 2022. \url{https://arxiv.org/abs/2210.10683v1}
\bibitem{ref253} Zero-shot sampling of adversarial entities in biomedical question answering. arXiv:2402.10527v1, 2024. \url{https://arxiv.org/abs/2402.10527v1}
\bibitem{ref254} Data. arXiv:1801.04992v2, 2017. \url{https://arxiv.org/abs/1801.04992v2}
\bibitem{ref255} Self-Knowledge Guided Retrieval Augmentation for Large Language Models. arXiv:2310.05002v1, 2023. \url{https://arxiv.org/abs/2310.05002v1}
\bibitem{ref256} Hybrid Reasoning Based on Large Language Models for Autonomous Car Driving. arXiv:2402.13602v3, 2024. \url{https://arxiv.org/abs/2402.13602v3}
\bibitem{ref257} LMDrive Closed-Loop End-to-End Driving with Large Language Models. arXiv:2312.07488v2, 2023. \url{https://arxiv.org/abs/2312.07488v2}
\bibitem{ref258} Prioritizing Safeguarding Over Autonomy Risks of LLM Agents for Science. arXiv:2402.04247v2, 2024. \url{https://arxiv.org/abs/2402.04247v2}
\bibitem{ref259} LanguageMPC Large Language Models as Decision Makers for Autonomous Driving. arXiv:2310.03026v2, 2023. \url{https://arxiv.org/abs/2310.03026v2}
\bibitem{ref260} Correcting misinformation on social media with a large language model. arXiv:2403.11169v2, 2024. \url{https://arxiv.org/abs/2403.11169v2}
\bibitem{ref261} The Flaws of Policies Requiring Human Oversight of Government Algorithms. arXiv:2109.05067v4, 2021. \url{https://arxiv.org/abs/2109.05067v4}
\bibitem{ref262} On the Quest for Effectiveness in Human Oversight Interdisciplinary Perspectives. arXiv:2404.04059v1, 2024. \url{https://arxiv.org/abs/2404.04059v1}
\bibitem{ref263} Human-in-the-loop Evaluation for Early Misinformation Detection A Case Study of COVID-19 Treatments. arXiv:2212.09683v4, 2022. \url{https://arxiv.org/abs/2212.09683v4}
\bibitem{ref264} 'Team-in-the-loop' organisational oversight of high-stakes AI. arXiv:2303.14007v1, 2023. \url{https://arxiv.org/abs/2303.14007v1}
\bibitem{ref265} Algorithmic Arbitrariness in Content Moderation. arXiv:2402.16979v1, 2024. \url{https://arxiv.org/abs/2402.16979v1}
\bibitem{ref266} Justice in Misinformation Detection Systems An Analysis of Algorithms, Stakeholders, and Potential Harms. arXiv:2204.13568v2, 2022. \url{https://arxiv.org/abs/2204.13568v2}
\bibitem{ref267} 'Put the Car on the Stand' SMT-based Oracles for Investigating Decisions. arXiv:2305.05731v2, 2023. \url{https://arxiv.org/abs/2305.05731v2}
\bibitem{ref268} How Platform-User Power Relations Shape Algorithmic Accountability A Case Study of Instant Loan Platforms and Financially Stressed Users in India. arXiv:2205.05661v1, 2022. \url{https://arxiv.org/abs/2205.05661v1}
\end{thebibliography}

\end{document}
